% !TEX TS-program = pdflatex
\pdfoutput=1
\documentclass[10pt,twocolumn,letterpaper]{article}

% ============================================================
% Packages
% ============================================================
\usepackage[margin=0.75in,columnsep=0.25in]{geometry}
\usepackage{amsmath,amsfonts,amssymb,amsthm,booktabs}
\usepackage{float}
\usepackage{algorithm,algpseudocode}
\usepackage[utf8]{inputenc}
\usepackage[numbers,sort&compress]{natbib}
\usepackage{hyperref}
\usepackage{xcolor}
\usepackage{enumitem}
\usepackage{microtype}
\usepackage{graphicx}
\usepackage{titlesec}
\usepackage[font=small,labelfont=bf,skip=4pt]{caption}
\usepackage{setspace}
\usepackage{etoolbox}
\usepackage{fancyhdr}
\usepackage{times}
% Safe figure include: compile even if figure files are missing
\makeatletter
\newcommand{\safeincludegraphics}[2][]{%
	\begingroup
	% Make \IfFileExists search the same paths as \includegraphics
	\let\input@path\Ginput@path
	\def\@foundfile{}%
	% Try without extension first (rare), then common extensions
	\IfFileExists{#2}{\def\@foundfile{#2}}{%
		\IfFileExists{#2.pdf}{\def\@foundfile{#2.pdf}}{%
			\IfFileExists{#2.png}{\def\@foundfile{#2.png}}{%
				\IfFileExists{#2.jpg}{\def\@foundfile{#2.jpg}}{%
					\IfFileExists{#2.jpeg}{\def\@foundfile{#2.jpeg}}{%
						\IfFileExists{#2.eps}{\def\@foundfile{#2.eps}}{}%
					}%
				}%
			}%
		}%
	}%
	\ifx\@foundfile\@empty
	\fbox{\parbox[c][0.18\textheight][c]{0.9\linewidth}{\centering
			Missing figure: \texttt{\detokenize{#2}}}}%
	\else
	\includegraphics[#1]{\@foundfile}%
	\fi
	\endgroup
}
\makeatother

\usepackage{multirow}
\usepackage{subcaption}
\graphicspath{
	{../results/paper/convergence_grid/}
	{../results/paper/perf_profiles/}
	{../results/supplement/single/}
	{../results_comm/paper/comm_profile/}
	{../results_ada/paper/ada_mechanism/}
	{../results_robust/paper/}
	{../results_robust/paper/param_sweep/}
	{../results_scale/paper/}
}

\emergencystretch=3em
\allowdisplaybreaks
\tolerance=1500
\hbadness=1500
\medmuskip=3mu plus 1mu minus 2mu
\thickmuskip=4mu plus 2mu minus 2mu

% ============================================================
% Compact two-column layout (AI conference style)
% ============================================================
\setstretch{0.97}

\titlespacing*{\section}{0pt}{8pt plus 2pt minus 2pt}{4pt plus 1pt minus 1pt}
\titlespacing*{\subsection}{0pt}{6pt plus 2pt minus 1pt}{3pt plus 1pt minus 1pt}
\titlespacing*{\subsubsection}{0pt}{4pt plus 1pt minus 1pt}{2pt plus 1pt minus 1pt}
\titleformat{\section}{\normalfont\large\bfseries}{\thesection}{0.6em}{}
\titleformat{\subsection}{\normalfont\normalsize\bfseries}{\thesubsection}{0.5em}{}
\titleformat{\subsubsection}{\normalfont\small\bfseries}{\thesubsubsection}{0.4em}{}

\setlength{\abovedisplayskip}{5pt plus 2pt minus 2pt}
\setlength{\belowdisplayskip}{5pt plus 2pt minus 2pt}
\setlength{\abovedisplayshortskip}{2pt plus 1pt minus 1pt}
\setlength{\belowdisplayshortskip}{2pt plus 1pt minus 1pt}
\setlength{\jot}{3pt}

\setlength{\textfloatsep}{8pt plus 2pt minus 2pt}
\setlength{\floatsep}{6pt plus 2pt minus 2pt}
\setlength{\intextsep}{6pt plus 2pt minus 2pt}
\setlength{\dbltextfloatsep}{8pt plus 2pt minus 2pt}
\setlength{\dblfloatsep}{6pt plus 2pt minus 2pt}

\setlength{\abovecaptionskip}{4pt}
\setlength{\belowcaptionskip}{2pt}

\setlength{\parskip}{1pt plus 0.5pt minus 0.5pt}

\setlist{nosep,leftmargin=1em}

\AtBeginEnvironment{proof}{\vspace{-2pt}}
\AtEndEnvironment{proof}{\vspace{-2pt}}

\setlength{\bibsep}{1pt plus 0.5pt}

\hypersetup{
	colorlinks=true,
	linkcolor=blue!80!black,
	filecolor=magenta,
	urlcolor=cyan!80!black,
	citecolor=blue!80!black,
}

% Avoid hyperref PDF string warnings for math in bookmarks
\pdfstringdefDisableCommands{%
	\def\cO{O}%
	\def\one{1}%
	\def\J{J}%
	\def\Pmat{P}%
	\def\norm#1{||#1||}%
	\def\ip#1#2{<#1,#2>}%
}

% ============================================================
% Theorem-like (compact spacing)
% ============================================================
\newtheoremstyle{compact}
  {4pt}{4pt}{\itshape}{}{\bfseries}{.}{0.5em}{}
\newtheoremstyle{compactremark}
  {4pt}{4pt}{}{}{\itshape}{.}{0.5em}{}
\theoremstyle{compact}
\newtheorem{theorem}{Theorem}[section]
\newtheorem{lemma}[theorem]{Lemma}
\newtheorem{proposition}[theorem]{Proposition}
\newtheorem{corollary}[theorem]{Corollary}
\newtheorem{assumption}[theorem]{Assumption}
\theoremstyle{compactremark}
\newtheorem{remark}[theorem]{Remark}

% ============================================================
% Macros
% ============================================================
\newcommand{\R}{\mathbb{R}}
\newcommand{\one}{\mathbf{1}}
\newcommand{\ip}[2]{\left\langle #1,#2\right\rangle}
\newcommand{\norm}[1]{\left\|#1\right\|}
\newcommand{\normtwo}[1]{\left\|#1\right\|_2}
\newcommand{\normF}[1]{\left\|#1\right\|_F}
\newcommand{\cO}{\mathcal{O}}
\newcommand{\J}{\left(\frac{1}{N}\one\one^\top\right)}
\newcommand{\Pmat}{I-\J}

\newcommand{\barx}{\bar x}
\newcommand{\barg}{\bar g}
\newcommand{\bars}{\bar s}

% ============================================================
% Title (compact AI conference style)
% ============================================================
\makeatletter
\def\@maketitle{%
  \newpage\null
  \begin{center}%
    {\Large\bfseries Distributed Gradient--Regularized Newton:\\[3pt]
     Scheduled Consensus and \texorpdfstring{$\mathcal{O}(\varepsilon^{-1})$}{O(eps\^{}-1)} Global Iteration Complexity\par}%
    \vskip 10pt
    {\normalsize
      \textbf{Wei Hu}\textsuperscript{1,2}\qquad
      \textbf{Pengcheng Xie}\textsuperscript{3}\qquad
      \textbf{Ya-xiang Yuan}\textsuperscript{1,*}\qquad
      \textbf{Li Zhang}\textsuperscript{1}\par}%
    \vskip 4pt
    {\small\itshape
      \textsuperscript{1}LSEC, ICMSEC, Academy of Mathematics and Systems Science,
      Chinese Academy of Sciences, Beijing, China\\
      \textsuperscript{2}University of Chinese Academy of Sciences, Beijing, China\qquad
      \textsuperscript{3}Lawrence Berkeley National Laboratory, Berkeley, CA, USA\\[2pt]
      \textsuperscript{*}Corresponding author.
      \texttt{huwei@amss.ac.cn},
      \texttt{pxie98@gmail.com},
      \texttt{yyx@lsec.cc.ac.cn},
      \texttt{zlzhanglily@163.com}\par}%
  \end{center}%
  \vskip 8pt}
\makeatother
\title{}
\author{}
\date{}

%%%%%%%%%%%%%%%%%%%%%%%%%%%%%%%%%%%%%%%%%%%%%%%%%%%%%%%%%%%%
\begin{document}

	\maketitle
	\vspace{-4pt}

\renewcommand{\abstractname}{\normalsize\bfseries Abstract}
\begin{abstract}\small\noindent
	We propose \textsc{DisGrem}, the first fully decentralized
	second-order method that provably matches the centralized
	regularized Newton iteration complexity---$\cO(\varepsilon^{-1})$
	for $\|\nabla f(\bar x_k)\|\le\varepsilon$---\emph{without line
	search, without stepsize tuning}, and \emph{without any static
	Hessian-heterogeneity constant}.
	Each agent maintains gradient and Hessian trackers, solves
	a vanishing-regularization Newton system
	$\lambda_{i,k}=\sqrt{M\|\tilde g_{i,k}\|}$ controlled by a single
	parameter~$M$, and communicates via two-stage gossip mixing.
	An eigenvalue-shift stabilizer ensures well-posed local solves
	even when tracked Hessians are transiently indefinite.
	
	Two analytical ideas drive the result:
	(i)~a \emph{reference-step} construction reduces the average
	iterate to an inexact regularized Newton recursion, and
	(ii)~an \emph{increment-based dispersion analysis} bounds
	tracker mismatch through Lipschitz \emph{differences} rather
	than static heterogeneity, yielding \emph{transient} (polynomially
	decaying) dispersion with no accuracy floor.
	Under a logarithmic mixing schedule the total communication
	cost is $\cO(\varepsilon^{-1}\log\varepsilon^{-1})$ rounds; with
	a fixed depth $N_c$ the method converges to an
	$\cO(\rho^{N_c})$-neighborhood.
	Strong convexity further yields local $Q$-superlinear ($3/2$-order)
	gradient convergence.
	
	Two practical variants---\textsc{CeDisGrem} (compressed Hessian
	exchange) and \textsc{AdaDisGrem} (online secant-based $M$
	adaptation)---and experiments on nine objectives against six
	baselines confirm that the \textsc{DisGrem} family reaches
	relF~$\sim\!10^{-6}$--$10^{-15}$ on most test cases, while
	every baseline stagnates or diverges on at least one problem.
	A 100-run robustness study finds \textsc{AdaDisGrem} attains
	the highest average success rate (80\,\%).
\end{abstract}
	
	%%%%%%%%%%%%%%%%%%%%%%%%%%%%%%%%%%%%%%%%%%%%%%%%%%%%%%%%%%%%
	\section{Introduction}
	%%%%%%%%%%%%%%%%%%%%%%%%%%%%%%%%%%%%%%%%%%%%%%%%%%%%%%%%%%%%
	Decentralized optimization addresses problems of the form
	\begin{equation}\label{eq:intro_prob}
		\min_{x\in\R^d} f(x)\;:=\;\frac{1}{N}\sum_{i=1}^N f_i(x),
	\end{equation}
	where each $f_i$ is held privately at node~$i$ of a connected network
	and nodes exchange information only with their neighbors.
	Communication is the principal bottleneck: every outer iteration
	requires one or more rounds of costly neighbor exchanges, so
	reducing the total number of iterations directly reduces the
	communication bill.
	First-order decentralized methods can require many iterations on
	ill-conditioned problems, multiplying this burden.
	Second-order information reduces outer iterations in principle,
	yet decentralized second-order methods face their own
	obstacles---curvature data are expensive to maintain, and global
	convergence typically demands careful stepsize control or a
	line search procedure.
	
	The regularized Newton method of \citet{mishchenko2023rn} offers a
	different route: a vanishing regularization
	$\lambda_k\asymp\sqrt{\|g_k\|}$ delivers the optimal $\cO(1/k^2)$
	functional rate under convexity and Lipschitz Hessians, with no
	stepsize to tune.
	The recent work of \citet{doikov2024gradient} further solidifies
	this framework by showing that gradient-norm regularization can
	replace the more expensive cubic model while retaining the same
	complexity guarantees; \citet{ding2025nonconvex} extend the
	approach to non-convex objectives through negative-curvature
	exploitation.
	A natural question is:
	\emph{can the same per-iteration complexity be preserved in a
	decentralized network, despite disagreement and tracking errors?}
	%
	Previous decentralized second-order methods have \emph{not} answered
	this in the affirmative.
	The fundamental obstacle is that the self-regulating stepsize
	$\lambda_k\asymp\sqrt{\|g_k\|}$ relies on globally consistent
	gradient and curvature information;
	when each agent sees only a \emph{local} approximation polluted by
	consensus error, the delicate balance between regularization and
	curvature breaks down, and existing methods fall back on line
	search or static heterogeneity assumptions to
	recover convergence.
	The core difficulty is twofold.
	First, each agent inverts a \emph{local} regularized system, and
	the resulting average step generally differs from the Newton step
	of the global objective---the gap depends on how well the agents
	agree, and must be controlled without introducing static
	heterogeneity constants.
	Second, the Hessian trackers can become transiently indefinite
	due to imperfect consensus, threatening the well-posedness of
	the local solves and the quality of the curvature information.
	\textsc{DisGrem} addresses the first challenge through an
	increment-based dispersion recursion that bounds tracker
	mismatch via Lipschitz \emph{differences}, and the second
	through an eigenvalue-shift stabilizer that ensures
	well-posedness at every iteration.
	To our knowledge, this is the first work to show that the
	centralized regularized Newton rate survives decentralization via
	gradient/Hessian tracking and two-stage mixing, and to do so
	without any static Hessian-heterogeneity constant.
	
\noindent\textbf{Contributions.}
The resulting algorithm, \textsc{DisGrem} (Distributed
Gradient-Regularized Newton Method), requires only a single
scaling parameter $M\ge L_2$---no stepsize schedule or line
search.
Our main contributions are as follows.

\begin{itemize}[leftmargin=1.2em,itemsep=2pt]
	\item \textbf{Algorithm design.}
	A vanishing regularization
	$\lambda_{i,k}=\sqrt{M\|\tilde g_{i,k}\|}$, where
	$\tilde g_{i,k}$ is agent~$i$'s pre-mixed gradient tracker,
	drives each local Newton step; as the tracked gradient shrinks,
	the regularization vanishes and the step approaches a pure
	Newton direction.
	Two-stage mixing decouples consensus accuracy from the local
	solve, and an eigenvalue-shift stabilizer keeps every local
	system well-posed even when Hessian trackers are transiently
	indefinite.
	
	\item \textbf{Reference-step framework.}
	A reference-step construction controls the gap between the
	average local step and the global Newton step, reducing the
	average iterate's update to an \textbf{inexact regularized Newton} step.
	
	\item \textbf{Increment-based dispersion analysis.}
	Earlier decentralized second-order analyses rely on a static
	Hessian heterogeneity constant that creates irreducible accuracy
	floors.
	We replace it with a dispersion recursion driven by
	Lipschitz-continuous \emph{differences}
	$\nabla F(X_{k+1}){-}\nabla F(X_k)$ and
	$\nabla^2 F(X_{k+1}){-}\nabla^2 F(X_k)$; no static
	heterogeneity constant appears in the final bounds, and
	the resulting tracker dispersion is \emph{transient},
	decaying polynomially.
	
	\item \textbf{Convergence rates.}
	A Lyapunov argument proves iterate boundedness from a coercivity
	condition alone (Lemma~\ref{lem:lyap_bound}), removing the
	common but hard-to-verify ``bounded level set'' assumption.
	We then obtain
	$\cO(\varepsilon^{-1})$ iterations for
	$\|\nabla f(\bar x_k)\|\le\varepsilon$
	(Theorem~\ref{thm:main})
	and $\cO(\varepsilon^{-1}\log\varepsilon^{-1})$ total mixing
	rounds with a logarithmic schedule.
	Fixing the mixing depth to a constant $N_c$ yields convergence
	to an $\cO(\rho^{N_c})$-neighborhood
	(Corollary~\ref{cor:fixed_nc}).
	Under additional strong convexity, the gradient converges
	locally at a $Q$-superlinear rate of order $3/2$
	(Theorem~\ref{thm:superlinear}).
	
	\item \textbf{Practical variants.}
	\textsc{CeDisGrem} applies low-rank compression to Hessian
	corrections, cutting per-iteration communication.
	\textsc{AdaDisGrem} adapts $M$ on the fly via secant-based
	curvature estimation, removing per-function tuning.
	
	\item \textbf{Numerical evidence.}
	A suite of nine benchmarks ($d{=}30$, 20 Monte~Carlo runs)
	against six first- and second-order baselines confirms that
	the \textsc{DisGrem} family is the most uniformly reliable
	group in our tests.
	Ablation studies isolate the effects of compression and lazy
	Hessian updates; \textsc{AdaDisGrem} is shown to be
	insensitive to its initial $M$ over a $20\times$ range.
\end{itemize}

\vspace{1pt}
\noindent\textbf{Paper organization.}
Section~\ref{sec:related} surveys related work.
Section~\ref{sec:prelim} fixes notation, states the standing
assumptions, and establishes iterate boundedness via a Lyapunov
argument.
The algorithm and its practical variants are described in
Section~\ref{sec:alg}.
Section~\ref{sec:theory} develops the convergence theory; a proof
roadmap at the beginning of that section describes the five-stage
pipeline.
Section~\ref{sec:experiments} presents the numerical experiments.
Proofs deferred from the main text are collected in
Appendices~\ref{app:dispersion} and~\ref{app:proofs}.

	%%%%%%%%%%%%%%%%%%%%%%%%%%%%%%%%%%%%%%%%%%%%%%%%%%%%%%%%%%%%
	\section{Related Work}\label{sec:related}
	%%%%%%%%%%%%%%%%%%%%%%%%%%%%%%%%%%%%%%%%%%%%%%%%%%%%%%%%%%%%
	
	\noindent\textbf{First-order decentralized methods.}
	Decentralized (sub)gradient descent goes back to
	\citet{nedic2009subgradient}.
	EXTRA~\citep{shi2015extra} and exact
	diffusion~\citep{yuan2019exactdiffusion} remove the bias of
	constant-stepsize DGD through a two-step correction and a
	primal-dual reformulation, respectively.
	Gradient tracking~\citep{nedic2017diging,xin2019distributed,pu2021distributed}
	achieves exact convergence with a single doubly stochastic matrix by
	accumulating gradient increments;
	\citet{alghunaim2020decentralized} unifies primal-dual and
	gradient-tracking perspectives.
	These methods converge at $\cO(1/k)$ or linearly for strongly
	convex objectives, but the per-iteration improvement scales
	with $(1{-}\rho)^{-1}$ or $(1{-}\rho)^{-2}$---a severe
	penalty on weakly connected graphs.
	Accelerated variants~\citep{jakovetic2014fast,li2021decentralized}
	sharpen the rate but remain first-order.
	Communication compression
	\citep{koloskova2019decentralized,beznosikov2023biased}
	is an orthogonal axis of improvement that can be combined with
	any of the above; we explore this idea for Hessian data in
	Section~\ref{ssec:variants}.
	
	\noindent\textbf{Second-order and quasi-Newton decentralized methods.}
	Network Newton~\citep{mokhtari2015networknewton_arxiv} approximates
	the global Newton direction via a truncated Hessian power series
	over multi-hop paths; Newton
	tracking~\citep{mokhtari2020newtontracking} embeds Newton-type
	directions into gradient tracking.
	DQM~\citep{eisen2017dqn} and other decentralized quasi-Newton
	schemes replace the exact Hessian with L-BFGS-type surrogates;
	\citet{bajovic2017newton} study distributed Newton-type corrections
	in sensor-network settings, and \citet{zhang2021distributed} propose
	a variance-reduced Newton method for finite-sum objectives.
	ESOM~\citep{mokhtari2016esom} combines exact second-order
	information with an ADMM-type consensus step.
	SONATA~\citep{sun2022sonata} solves a sequence of strongly
	convex local surrogates and covers non-convex objectives, while
	Network-GIANT~\citep{ye2023networkgiant} constructs a
	global Newton direction through harmonic-mean Hessian consensus.
	%
	A critical weakness shared by all the above methods is the need for
	explicit stepsize tuning or line search mechanisms, and \emph{none}
	has been shown to match the centralized Newton rate in a rigorous
	global sense.
	More fundamentally, Hessian heterogeneity across nodes---a key
	difficulty---is handled either by ignoring it or by bounding it with
	a \emph{static} constant $\sigma_H$ that creates an irreducible
	accuracy floor: the final iterate satisfies
	$\|\nabla f(\bar x)\|\le\cO(\sigma_H)$ regardless of the number
	of iterations.
	Our increment-based dispersion analysis eliminates this floor
	entirely.
	
	\noindent\textbf{Communication-efficient and inversion-free
	second-order methods.}
	The $\cO(d^2)$ per-iteration cost of transmitting Hessian data and
	the $\cO(d^3)$ cost of solving local Newton systems are two
	computational bottlenecks of decentralized second-order methods.
	On the communication side,
	\citet{zhang2024distributed} combine lazy Hessian updates with
	compression for distributed cubic Newton, but rely on a central
	parameter server.
	Our \textsc{CeDisGrem} achieves analogous communication savings
	in a fully peer-to-peer topology.
	On the computation side,
	DINAS~\citep{crane2025dinas} proposes an adaptive-stepsize
	distributed inexact Newton method that avoids local Hessian
	inversion via iterative solvers such as conjugate gradients,
	and INDO~\citep{yuan2023indo} introduces an inversion-free
	second-order method for consensus optimization.
	Incorporating such inexact solvers into the DisGrem framework
	is a promising direction for scaling to large $d$
	(Section~\ref{sec:conclusion}).
	
	\noindent\textbf{Regularized Newton and cubic regularization.}
	Our starting point is the centralized regularized Newton method
	of \citet{mishchenko2023rn}:
	setting $\lambda_k\asymp\sqrt{\|\nabla f(x_k)\|}$ yields the
	tight $\cO(1/k^2)$ functional rate under convexity and
	Lipschitz Hessians, with no stepsize to choose.
	The idea has roots in the cubic regularization of
	\citet{nesterov2006cubic,cartis2011arc}, which achieves the
	same rate through an adaptive cubic model.
	\citet{doikov2024gradient} rigorously establish
	that gradient-norm regularization can replace the more
	expensive cubic model across a broad class of problems,
	cementing gradient regularization as a first-class alternative.
	Very recently, \citet{ding2025nonconvex} extend gradient-regularized
	Newton to non-convex objectives by incorporating negative-curvature
	directions---a centralized breakthrough that appeared only in 2025,
	underscoring why our convex-first approach is well-motivated.
	\citet{doikov2021optimization} extend second-order methods to
	higher-order tensor regularization, and
	\citet{hanzely2022newton} study Newton-type methods with lazy
	Hessian updates in the centralized setting, an idea we also
	explore in Section~\ref{ssec:commexp}.
	Our work is, to our knowledge, the first to extend the
	centralized regularized Newton framework to a fully
	decentralized setting with gradient and Hessian tracking,
	preserving the $\cO(\varepsilon^{-1})$ rate while eliminating
	all static Hessian-heterogeneity dependence.
	
	%%%%%%%%%%%%%%%%%%%%%%%%%%%%%%%%%%%%%%%%%%%%%%%%%%%%%%%%%%%%
	\section{Preliminaries and Assumptions}\label{sec:prelim}
	%%%%%%%%%%%%%%%%%%%%%%%%%%%%%%%%%%%%%%%%%%%%%%%%%%%%%%%%%%%%
	\subsection{Network model and notation}
	Communication uses a symmetric, nonnegative, doubly stochastic mixing matrix $W\in\R^{N\times N}$:
	$W_{ij}\ge0$, $W\one=\one$, and $\one^\top W=\one^\top$.
	Define the averaging projector $\J=\frac{1}{N}\one\one^\top$ and $\Pmat=I-\J$.
	Let $\rho := \normtwo{W-\J}\in[0,1)$, strictly less than $1$ for a connected undirected network.
	For stacked vectors $Z=[z_1;\dots;z_N]\in\R^{Nd}$ we use the Euclidean norm $\|Z\|$.
	We write $\bar z:=\frac1N\sum_{i=1}^N z_i$ for averages.
	We also use stacked primal variables $X=[x_1;\dots;x_N]\in\R^{Nd}$ and define the separable network objective
	\[
	F(X):=\sum_{i=1}^N f_i(x_i).
	\]
	Then $\nabla F(X)=[\nabla f_1(x_1);\dots;\nabla f_N(x_N)]\in\R^{Nd}$ and
	$\nabla^2 F(X)=\mathrm{blkdiag}(\nabla^2 f_1(x_1),\dots,\nabla^2 f_N(x_N))$.
	For convenience, we also define the stacked Hessian vectorization
	\[
	\mathcal H(X):=\big[\mathrm{vec}(\nabla^2 f_1(x_1));\dots;\mathrm{vec}(\nabla^2 f_N(x_N))\big]\in\R^{Nd^2}.
	\]
	When evaluated at a consensus point $x$, we write $\mathcal H(\one\otimes x)$.
	For the consensus problem \eqref{eq:intro_prob}, we write
	$f(x)=\frac1N\sum_{i=1}^N f_i(x)$, $g(x)=\nabla f(x)$, and $H(x)=\nabla^2 f(x)$.
	
	\subsection{Standing assumptions}
	\begin{assumption}\label{ass:standing}
		\begin{enumerate}[nosep,leftmargin=1em]
			\item \textbf{Convexity and smoothness.}
			Each $f_i$ is convex and twice continuously differentiable.
			Its gradient is $L_1$-Lipschitz and its Hessian is $L_2$-Lipschitz:
			\[
			\norm{\nabla f_i(x)\!-\!\nabla f_i(y)}\le L_1\norm{x\!-\!y},
			\]
			\vspace{-8pt}
			\[
			\normtwo{\nabla^2 f_i(x)\!-\!\nabla^2 f_i(y)}\le L_2\norm{x\!-\!y}.
			\]
			\item \textbf{Connected network.} $W$ is symmetric, doubly stochastic, and $\rho=\normtwo{W-\J}<1$.
			\item \textbf{Coercivity (compact sublevel sets).}
			The average objective $f$ is coercive: $f(x)\to+\infty$ as $\|x\|\to\infty$.
			Equivalently, for every $c\in\R$ the sublevel set $\{x:f(x)\le c\}$ is bounded.
			\item \textbf{Regularization scaling.} The user-chosen scaling factor satisfies $M\ge L_2$.
		\end{enumerate}
	\end{assumption}
	
	\begin{remark}[PSD Hessians]
	Since each $f_i$ is convex (item~1), $\nabla^2 f_i(x)\succeq 0$
	for all~$x$ and all~$i$.
	We state this explicitly because the eigenvalue-shift stabilizer
	$\delta_{i,k}$ is designed to handle indefinite tracker matrices
	that may arise transiently; on convex problems, the true Hessians
	are always PSD.
	\end{remark}
	
	\begin{remark}[Coercivity as a sufficient condition]
	The proof only requires that the initial sublevel set
	$\{x : f(x) \le f(\bar x_0)\}$ is bounded;
	coercivity is a convenient sufficient condition that holds
	whenever the objective includes an~$\ell_2$ regularizer.
	When it does not hold, adding a small regularizer
	$(\mu/2)\|x\|^2$ ($\mu>0$) restores coercivity with an $\cO(\mu)$ bias;
	alternatively, one may directly assume bounded sublevel sets.
	\end{remark}
	
	% ============================================================
	% Dispersion notation and the Lyapunov-based boundedness result
	% ============================================================

	\noindent\textbf{Dispersion notation.}
	For a stacked vector $Z=[z_1;\dots;z_N]\in\R^{Nd}$
	we write $\bar z=\frac1N\sum_{i=1}^N z_i$ and define the RMS dispersion
	(also called \emph{consensus error} or \emph{disagreement} in the
	decentralized optimization literature, see e.g.\
	\citet{nedic2009subgradient,nedic2017diging})
	\(
	  D(Z):=\bigl(\tfrac1N\sum_{i=1}^N\|z_i-\bar z\|^2\bigr)^{1/2}.
	\)
	In particular, $D(X_k)$ and $D(\tilde G_k)$ measure the primal
	iterate and gradient-tracker disagreements at iteration $k$.

	\begin{lemma}[Iterate boundedness via a Lyapunov argument]\label{lem:lyap_bound}
		Assume Assumption~\ref{ass:standing}(1)--(3) (convexity, connected network, coercivity).
		By coercivity (Assumption~\ref{ass:standing}(3)), the minimizer set is non-empty;
		let $x_\star\in\arg\min f$ and define
		$\Phi_0:=f(\bar x_0)-f_\star$, where $\bar x_0:=\frac1N\sum_{i=1}^N x_{i,0}$.
		Consider the generalized Lyapunov function
		\[
		\mathcal L_k \;:=\; f(\bar x_k)-f_\star \;+\; c_1 D(X_k)^2 \;+\; c_2 D(\tilde G_k)^2,
		\]
		where $c_1,c_2>0$ are constants depending only on $L_1,L_2,M,\rho$
		(see Appendix~\ref{app:proofs} for explicit expressions), and
		$D(X_k), D(\tilde G_k)$ denote the RMS primal and gradient-tracker dispersions.
		Then the following hold:
		\begin{enumerate}[leftmargin=1.2em,itemsep=2pt]
			\item $\mathcal L_k$ is uniformly bounded: $\sup_{k\ge 0}\mathcal L_k \le \bar{\mathcal L} < \infty$ for some finite $\bar{\mathcal L}$ depending on $\mathcal L_0$ and problem parameters.
			\item Consequently, there exists a finite constant $D>0$
			(depending on $\mathcal L_0$, $f_\star$, and the coercivity modulus of $f$)
			such that
			\begin{equation}\label{eq:bounded_levelset}
				\norm{\bar x_k\!-\!x_\star}\!\le\! D,\;
				\norm{x_{i,k}\!-\!x_\star}\!\le\! D,\;
				\norm{\tilde x_{i,k}\!-\!x_\star}\!\le\! D,
			\end{equation}
			for all $k\ge 0$ and all $i\in\{1,\dots,N\}$.
		\end{enumerate}
	\end{lemma}

	\begin{proof}[Proof sketch (full details in Appendix~\ref{app:proofs})]
	The argument relies on three ingredients developed later in the
	paper: a per-step size bound
	(Lemma~\ref{lem:M-vs-s}, Section~\ref{sec:theory}),
	and primal/gradient dispersion contraction estimates
	(Lemmas~\ref{lem:DX_rec}--\ref{lem:DG_rec},
	Appendix~\ref{app:dispersion}).
	We summarize the logic here; the reader may defer the details
	until after the full analysis has been developed.

	We construct an \emph{invariant set} argument parameterized by
	a single scalar $R>0$.
	Define
	\begin{multline*}
	  \mathcal S_R := \bigl\{(X_k,\tilde G_k) :
	  f(\bar x_k)\le f_\star+R,\\
	  D(X_k)\le B_X(R),\;
	  D(\tilde G_k)\le B_G(R)\bigr\},
	\end{multline*}
	where $B_X(R),B_G(R)$ are the geometric-series fixed points of the
	dispersion recursions (Lemmas~\ref{lem:DX_rec},~\ref{lem:DG_rec}),
	with all Lipschitz constants evaluated on the sublevel set
	$S(R):=\{x:f(x)\le f_\star+R\}$, which is compact by coercivity.

	\emph{One-step invariance.}
	Suppose the current state lies in $\mathcal S_R$.
	The key observation is that every local step satisfies
	$\|s_{i,k}\|\le\sqrt{\|\tilde g_{i,k}\|/M}$
	(Lemma~\ref{lem:M-vs-s})
	\emph{unconditionally}---no inexact Newton condition is needed.
	This yields:
	\begin{itemize}[leftmargin=1.2em,itemsep=1pt]
	  \item \emph{Bounded function-value change.}
	        By $\bar x_{k+1}=\bar x_k+\bar s_k$ and $L_1$-smoothness,
	        \(
	        f(\bar x_{k+1})-f(\bar x_k)
	        \le \|\nabla f(\bar x_k)\|\,\|\bar s_k\|
	            +\tfrac{L_1}{2}\|\bar s_k\|^2
	        \le R_\Phi(R),
	        \)
	        where $R_\Phi(R)<\infty$ depends on $R$ through the Lipschitz
	        constants and the gradient/step bounds on $\mathcal S_R$.
	        Monotone descent is \emph{not} required.
	  \item \emph{Dispersion contraction.}
	        Primal and gradient-tracker dispersions each contract by
	        $\rho^{2t_k}<1$ plus a bounded forcing term controlled by
	        the step bound above and the true-Hessian bound
	        $M_H(R):=\sup_{x\in S(R)}\max_i\|\nabla^2 f_i(x)\|_2$
	        (finite by compactness of $S(R)$).
	\end{itemize}
	All constants ($R_\Phi$, $B_X$, $B_G$) are explicit functions of
	$R$ alone (Appendix~\ref{app:proofs},
	eqs.~\eqref{eq:rphi_R}--\eqref{eq:disp_fp});
	the quantity $D$ being proved finite does \emph{not} enter these
	definitions, so no circularity arises.
	Choosing $c_1,c_2$ so that the Lyapunov function
	$\mathcal L_k=f(\bar x_k)-f_\star+c_1D(X_k)^2+c_2D(\tilde G_k)^2$
	absorbs the cross-coupling terms, and then taking $R$ large enough
	that $\mathcal L_0\le R$ and $R_\Phi(R)$ is dominated by the
	dispersion contraction, one verifies that $\mathcal S_R$ is
	forward-invariant and
	$\sup_{k}\mathcal L_k\le\bar{\mathcal L}<\infty$.

	Boundedness of $\mathcal L_k$ confines $\bar x_k$ to a compact sublevel set,
	and each $x_{i,k}$ deviates from $\bar x_k$ by at most
	$\sqrt{N}D(X_k)\le\sqrt{N\bar{\mathcal L}/c_1}$,
	yielding~\eqref{eq:bounded_levelset}.
	\end{proof}

	Once Lemma~\ref{lem:lyap_bound} establishes the existence of $D$,
	we define all subsequent constants on the compact set $\{x:\|x-x_\star\|\le D\}$.
	In particular, define the uniform local Hessian bound
	\begin{equation}\label{eq:MHmax_def}
		M_{H,\max}\!:=\!\max_{1\le i\le N}\sup\{\|\nabla^2 f_i(x)\|_2:\|x\!-\!x_\star\|\!\le\! D\},
	\end{equation}
	which is finite by twice continuous differentiability on the compact set \eqref{eq:bounded_levelset}.
	Clearly $\|\nabla^2 f(x)\|_2\le M_{H,\max}$ for all $\|x-x_\star\|\le D$.

	\begin{remark}[Bounded Hessians on compact sets]
		Since each $f_i$ is twice continuously differentiable (Assumption~\ref{ass:standing}(1))
		and the iterates are confined to a compact set by Lemma~\ref{lem:lyap_bound},
		the quantity $\sup_{x\in\mathcal B}\max_{1\le i\le N}\|\nabla^2 f_i(x)\|_2$
		is automatically finite for any bounded $\mathcal B\subset\R^d$.
		No separate assumption is needed.
	\end{remark}

	\begin{remark}[No static Hessian-heterogeneity assumption]
		Several decentralized second-order analyses posit a bounded
		\emph{static} heterogeneity constant
		$\sigma_H:=\sup_{\|x-x_\star\|\le D}
		  \frac{1}{\sqrt N}\|(\Pmat\otimes I_{d^2})\mathcal H(\one\otimes x)\|$.
		We avoid this entirely: Hessian-tracker dispersion is controlled
		by an increment-based recursion
		(Appendix~\ref{app:dispersion}, Lemma~\ref{lem:DH_rec})
		that tracks mismatch through Lipschitz-continuous
		\emph{differences} $\nabla^2 F(X_{k+1})-\nabla^2 F(X_k)$,
		so no accuracy floor linked to~$\sigma_H$ appears in the
		final bounds.
	\end{remark}

	
	%%%%%%%%%%%%%%%%%%%%%%%%%%%%%%%%%%%%%%%%%%%%%%%%%%%%%%%%%%%%
	
\section{Algorithm}\label{sec:alg}
	%%%%%%%%%%%%%%%%%%%%%%%%%%%%%%%%%%%%%%%%%%%%%%%%%%%%%%%%%%%%
	
	\subsection{Design motivation}\label{ssec:design}
	
	Lifting regularized Newton to a decentralized setting requires
	addressing three coupled difficulties.
	
	\textbf{(i) Averaging vs.\ inversion.}
	Centralized regularized Newton inverts the average system
	$(\nabla^2 f(\bar x_k)+\lambda_k I)^{-1}$.
	In a network, each agent solves a \emph{local} system; the
	resulting average step
	$\bar s_k=\frac1N\sum_i s_{i,k}$ generally differs from the
	``ideal'' centralized step
	$-(\frac1N\sum_i A_{i,k})^{-1}(\frac1N\sum_i g_{i,k})$,
	because the average of matrix inverses is not the inverse of
	the average matrix.
	Pre-mixing drives the local inputs
	$(x_{i,k},g_{i,k},H_{i,k})$ closer to their averages,
	shrinking the resulting \emph{step dispersion}
	(Section~\ref{sec:theory}, Lemma~\ref{lem:step-disp}).
	
	\textbf{(ii) Hessian maintenance and positivity.}
	Tracking propagates Hessian increments
	$\nabla^2 f_j(x_{j,k+1})-\nabla^2 f_j(x_{j,k})$ through the
	network.
	Floating-point effects can render the tracked matrices
	non-symmetric or slightly indefinite, making the local systems
	ill-conditioned.
	We address this in two steps: symmetrize the Hessian data
	(both the pre-mixed Hessian for the local solve in step~B
	and the post-mixed tracker in step~D of Algorithm~\ref{alg:disgrem}),
	then add an \emph{eigenvalue-shift stabilizer}
	$\delta_{i,k}:=\max\{0,-\lambda_{\min}(\mathrm{sym}(\tilde H_{i,k}))\}$.
	The system matrix
	$\mathrm{sym}(\tilde H_{i,k})+(\lambda_{i,k}+\delta_{i,k})I$
	is then uniformly SPD, so every local solve is well-posed
	without an explicit PSD projection.
	
	\textbf{(iii) Stepsize-free regularization.}
	Setting $\lambda_{i,k}=\sqrt{M\|\tilde g_{i,k}\|}$ mirrors
	the centralized choice
	$\lambda_k\asymp\sqrt{\|\nabla f(\bar x_k)\|}$ that yields
	$\cO(1/k^2)$ convergence with no line search.
	The constant $M\ge L_2$ acts as a damping parameter: a larger $M$
	widens the admissible heterogeneity range and improves
	stability on ill-conditioned or non-convex problems
	(Sections~\ref{ssec:commexp}--\ref{ssec:ada}), while the
	regularizer still vanishes as $\|\nabla f\|\to 0$.
	We note that $M$ plays a role analogous to a trust-region radius
	in classical second-order methods: it controls the step magnitude
	implicitly (Lemma~\ref{lem:M-vs-s}: $\|s_{i,k}\|\le\sqrt{\|\tilde g_{i,k}\|/M}$),
	but does not require per-iteration adjustment or a line search subroutine.
	
	\subsection{\textsc{DisGrem}: full algorithm}\label{ssec:disgrem}
	
	Each node $i$ stores a primal variable $x_{i,k}\in\R^d$, a gradient tracker $g_{i,k}\in\R^d$,
	and a Hessian tracker $H_{i,k}\in\R^{d\times d}$.
	$\mathrm{sym}(A):=\tfrac12(A+A^\top)$ denotes symmetrization.
	
	\begin{algorithm}[t]
\caption{\textsc{DisGrem}}
\label{alg:disgrem}
\small
\begin{algorithmic}[1]
\State \textbf{Input:} $\{x_{i,0}\}$, $W$, $L_1,L_2$, $M\ge L_2$.
\State \textbf{Init:} $g_{i,0}\!\gets\!\nabla f_i(x_{i,0})$, $H_{i,0}\!\gets\!\nabla^2 f_i(x_{i,0})$.
\For{$k=0,1,2,\dots$}
\State \textbf{(A) Pre-mixing} ($\tau_k\!\ge\!1$ rounds):
\Statex\quad $\tilde x_{i,k}\!\gets\!\textstyle\sum_j [W^{\tau_k}]_{ij}x_{j,k}$
\Statex\quad $\tilde g_{i,k}\!\gets\!\textstyle\sum_j [W^{\tau_k}]_{ij}g_{j,k}$
\Statex\quad $\tilde H_{i,k}\!\gets\!\textstyle\sum_j [W^{\tau_k}]_{ij}H_{j,k}$
\State \textbf{(B) Local Newton step} ($\forall\,i$):
\Statex\quad $\lambda_{i,k}\!\gets\!\sqrt{M\|\tilde g_{i,k}\|}$
\Statex\quad $\delta_{i,k}\!\gets\!\max\{0,-\lambda_{\min}(\mathrm{sym}(\tilde H_{i,k}))\}$
\Statex\quad \textbf{Solve} $(\mathrm{sym}(\tilde H_{i,k})\!+\!(\lambda_{i,k}\!+\!\delta_{i,k})I)s_{i,k}\!=\!-\tilde g_{i,k}$
\Statex\quad Set $y_{i,k+1}\!\gets\!\tilde x_{i,k}+s_{i,k}$
\State \textbf{(C) Post-mixing} ($t_k\!\ge\!1$ rounds):
\Statex\quad $x_{i,k+1}\!\gets\!\textstyle\sum_j [W^{t_k}]_{ij}y_{j,k+1}$
\State \textbf{(D) Tracker updates:}
\Statex\quad $g_{i,k+1}\!\gets\!\textstyle\sum_j [W^{t_k}]_{ij}(\tilde g_{j,k}+\nabla f_j(x_{j,k+1})\!-\!\nabla f_j(x_{j,k}))$
\Statex\quad $\widehat H_{i,k+1}\!\gets\!\textstyle\sum_j [W^{t_k}]_{ij}(\tilde H_{j,k}+\nabla^2 f_j(x_{j,k+1})\!-\!\nabla^2 f_j(x_{j,k}))$
\Statex\quad $H_{i,k+1}\!\gets\!\mathrm{sym}(\widehat H_{i,k+1})$
\EndFor
\end{algorithmic}
\end{algorithm}
	
\subsection{Communication cost per iteration}\label{ssec:comm}

Each outer iteration of \textsc{DisGrem} involves $(\tau_k + t_k)$ rounds of neighbor communication.
In each round every agent $i$ sends messages to and receives messages from its $|\mathcal N_i|$ neighbors.
In \textsc{DisGrem}, the natural payload exchanged per neighbor per round consists of:
\begin{itemize}[leftmargin=1.2em,itemsep=1pt]
	\item two $d$-dimensional vectors (e.g., mixing of $(x,g)$ in pre-mixing, or mixing of $(y,g)$ in post-mixing): $2d$ floats,
	\item one $d\times d$ symmetric matrix (Hessian tracker or Hessian increment): $\tfrac{d(d+1)}{2}$ floats.
\end{itemize}
Equivalently, the two vectors can be concatenated into a single $2d$-vector; this does not change the byte complexity.
In practice, the post-mixing step~(C) and the tracker-update step~(D)
use the \emph{same} $t_k$~rounds of gossip: every agent packs
$(y_i,\; \tilde g_i+\Delta\nabla f_i,\; \tilde H_i+\Delta\nabla^2 f_i)$
into a single message and applies $W^{t_k}$ once, so the total rounds
per outer iteration are $(\tau_k+t_k)$, not $(\tau_k+2t_k)$.

Let $n_v$ denote the number of $d$-vectors transmitted per round (for \textsc{DisGrem}, $n_v=2$).
With double precision (8 bytes per float), the \emph{sent} message volume per agent per outer iteration is
\begin{equation}\label{eq:comm_cost_iter}
	C_{\mathrm{iter}}^{\mathrm{send}}(k)
	= (\tau_k+t_k)\cdot |\mathcal N_i|\cdot \Bigl(n_v d + \tfrac{d(d+1)}{2}\Bigr)\cdot 8 \ \ \text{bytes}.
\end{equation}
If one counts both sending and receiving volume, then
$C_{\mathrm{iter}}^{\mathrm{send+recv}}(k)=2\,C_{\mathrm{iter}}^{\mathrm{send}}(k)$.
All communication figures reported in Section~\ref{sec:experiments}
use the send+recv per-agent convention.

For the default experimental setting $N{=}10$, $d{=}30$, and an Erd\H{o}s--R\'enyi (ER)
graph with $p{=}0.5$ (average degree $|\mathcal N_i|\approx 4.5$),
taking $\tau_k=t_k=3$ and $n_v=2$ yields
\[
C_{\mathrm{iter}}^{\mathrm{send}}\!\approx\! 6\!\times\! 4.5\!\times\!(60\!+\!465)\!\times\! 8
\!\approx\! 113\,\text{kB/agent/iter}.
\]
This is consistent with the empirical communication magnitudes reported in
Table~\ref{tab:summary} (up to graph-degree variability and implementation details).
	
	\subsection{Practical variants}\label{ssec:variants}
	
\noindent\textbf{\textsc{CeDisGrem}: compressed Hessian communication.}
The dominant communication cost of \textsc{DisGrem} is the $\frac{d(d+1)}{2}$-dimensional
Hessian increment $\nabla^2 f_j(x_{j,k+1})-\nabla^2 f_j(x_{j,k})$.
\textsc{CeDisGrem} replaces each symmetric Hessian increment
$\Delta H_{j,k}:=\nabla^2 f_j(x_{j,k+1})-\nabla^2 f_j(x_{j,k})$
by a rank-$r$ symmetric truncation based on the eigen-decomposition.
Let $\Delta H_{j,k}=V\operatorname{diag}(\mu_1,\dots,\mu_d)V^\top$ with
$|\mu_1|\ge \cdots \ge |\mu_d|$ and $V=[v_1,\dots,v_d]$ orthogonal.
Define the truncated approximation
\[
\hat\Delta H_{j,k}:=\sum_{\ell=1}^r \mu_\ell v_\ell v_\ell^\top
= V_r\Lambda_r V_r^\top,
\]
where $V_r=[v_1,\dots,v_r]\in\R^{d\times r}$ and $\Lambda_r=\operatorname{diag}(\mu_1,\dots,\mu_r)\in\R^{r\times r}$.
This $\hat\Delta H_{j,k}$ is the best rank-$r$ symmetric approximation of $\Delta H_{j,k}$ in Frobenius norm
(equivalently, the best rank-$r$ approximation for a symmetric matrix is obtained by keeping the $r$
eigenpairs with largest absolute eigenvalues).
Thus, transmitting $(V_r,\Lambda_r)$ costs $r(d+1)$ floats instead of $\frac{d(d+1)}{2}$ floats per neighbor per round.

The truncation error satisfies
\begin{align*}
\normF{\Delta H_{j,k}\!-\!\hat\Delta H_{j,k}}
&=\Bigl(\textstyle\sum_{\ell=r+1}^d \mu_\ell^2\Bigr)^{\!1/2}\\
&\le |\mu_{r+1}(\Delta H_{j,k})|\sqrt{d\!-\!r},
\end{align*}
where $\mu_{r+1}(\Delta H_{j,k})$ denotes the $(r+1)$-th largest absolute eigenvalue.
This error is absorbed into the Hessian-tracker dispersion and thus into the regularization automatically.
Empirically, $r=\lceil d/5\rceil$ reduces communication by 30--50\% with negligible iteration overhead
(Section~\ref{ssec:commexp}).

\begin{remark}[Local computation cost of compression]
The exact eigendecomposition of $\Delta H_{j,k}\in\R^{d\times d}$
costs $\cO(d^3)$ per agent per iteration---the same order as the
Newton-system solve.
When $d$ is large, randomized SVD
(cost $\cO(d^2 r)$) or Nystr\"om-type approximations can replace
the exact decomposition.
The additional approximation error enters the tracker dispersion
additively and is expected to be absorbed by the existing
regularization margin, although a formal analysis is deferred to
future work.
\end{remark}
	
	\noindent\textbf{\textsc{AdaDisGrem}: adaptive scaling factor.}
	The default choice $M=\mathrm{const}$ requires prior knowledge of $L_2$.
	\textsc{AdaDisGrem} maintains a running estimate $\hat M_{i,k}$ of
	each agent's local Hessian Lipschitz constant using a secant condition.
	Specifically, define the local secant ratio
	\begin{gather*}
	\hat L_{i,k}\gets
	\frac{\normtwo{\nabla^2 f_i(x_{i,k})-\nabla^2 f_i(x_{i,k-1})}}
	     {\norm{x_{i,k}-x_{i,k-1}}},\\
	\hat M_{i,k}\gets \max\!\bigl(\gamma\,\hat M_{i,k-1},\;
	\zeta\,\min(\hat L_{i,k},\,\eta_c\,\hat M_{i,0})\bigr),
	\end{gather*}
	with decay factor $\gamma\in(0,1)$ and cap $\eta_c>0$.
	The vanishing regularizer becomes $\lambda_{i,k}=\sqrt{\hat M_{i,k}\|\tilde g_{i,k}\|}$.
	The decay factor ensures that $\hat M_{i,k}$ decreases when the local curvature
	stabilizes, enabling faster Newton steps near the solution; the cap prevents
	transient large Hessian increments from inflating the estimate.
	This mechanism removes the need to specify $L_2$ \emph{a priori}
	at the cost of a mild adaptation lag.
	Note that $\hat L_{i,k}$ is a per-iterate secant approximation to the
	\emph{local} Hessian Lipschitz constant; it is distinct from the
	static proxy $H_{\max}^0 := \max_i\|\nabla^2 f_i(x_0)\|_2$ used to
	scale the default $M$ and baseline stepsizes (Section~\ref{ssec:practical}).
	Because the secant ratio is a one-dimensional finite-difference
	approximation of the operator-norm Lipschitz constant, it may
	underestimate $L_2$ when the dominant curvature variation occurs
	along directions orthogonal to $x_{i,k}-x_{i,k-1}$.
	Near a strongly convex minimizer with $\nabla^2 f(x_\star)\succ 0$,
	the secant ratio $\hat L_{i,k}$ converges to the directional second
	derivative of the Hessian map along the Newton direction, which
	equals $L_2$ when the curvature variation is isotropic.
	In general, $\hat L_{i,k}$ provides a lower bound on $L_2$
	restricted to the one-dimensional secant subspace; the
	max-smoothing ($\hat M_{i,k}\ge\zeta\hat L_{i,k}$) and the cap
	$\eta_c$ together ensure that the estimate does not collapse to zero.
	A full theoretical treatment (in particular, establishing
	$\hat M_{i,k}\ge L_2$ eventually under suitable regularity
	conditions) is left for future work;
	Section~\ref{ssec:ada} provides an extensive empirical study showing
	that AdaDisGrem is robust to the initial $\hat M_{i,0}$ across a $20\times$ range.
	
	\vspace{1pt}
	\noindent\textbf{\textsc{CeAdaDisGrem}: combined variant.}
	\textsc{CeAdaDisGrem} applies both Hessian compression (Ce) and
	adaptive $M$ (Ada) simultaneously.
	Because the compressed Hessian increments introduce approximation noise
	into the secant-based Lipschitz estimation, the interaction between
	the two mechanisms is nontrivial: Section~\ref{ssec:commexp} shows
	that this combination works well at moderate precision but may slow
	convergence at very high precision.
	
	\subsection{Practical guidelines for parameter selection}\label{ssec:practical}
	
	\begin{itemize}[leftmargin=1.2em,itemsep=2pt]
		\item \textbf{Scaling factor $M$.}
		Theorem~\ref{thm:main} requires $M\ge L_2$
		(the Hessian Lipschitz constant, \emph{not} the gradient Lipschitz
		constant $L_1$).
		When $L_2$ is unknown---as is the common case---$M$ can be treated
		as a robustness parameter: larger $M$ increases damping and ensures the
		regularized Newton step remains well-conditioned, at the cost of
		slower asymptotic convergence.
		A practical default is $M = M_{\mathrm{fac}} \cdot H_{\max}^{0}$ where
		$H_{\max}^{0} := \max_i\|\nabla^2 f_i(x_0)\|_2$ is a readily
		computable Hessian-norm proxy;
		Table~\ref{tab:config} lists per-function values of $M_{\mathrm{fac}}$
		ranging from $0.1$ (well-conditioned convex) to $15$ (multimodal non-convex).
		We stress that $H_{\max}^{0}$ is \emph{not} an estimate of the Hessian
		Lipschitz constant~$L_2$; the theoretical condition $M\ge L_2$
		is not directly enforced by this default.
		\textsc{AdaDisGrem} addresses the gap via online secant-based
		estimation of~$L_2$ (Section~\ref{ssec:ada}), and the
		adaptive-mechanism study shows that convergence is robust
		across a $20\times$ range of initial $M$ values.
		\item \textbf{Consensus rounds $N_c$.}
		The theory requires $\tau_k, t_k = \cO(\log 1/\|\tilde{\bar g}_k\|)$, which grows
		logarithmically.  In practice, a fixed $N_c = 3$--$5$ rounds per stage is sufficient
		for most networks with $\rho \le 0.95$.
		For strongly sparse networks ($\rho > 0.97$), increasing $N_c$ to $8$--$10$
		significantly improves stability.
		\item \textbf{Hessian compression rank $r$ (\textsc{CeDisGrem}).}
		Set $r = \lceil 0.2d \rceil$ as a default; increase $r$ if convergence slows.
		For diagonal-dominant Hessians (e.g.\ $\ell_2$-regularized losses), even $r=1$
		is often sufficient.
		Alternatively, element-wise top-$k$ sparsification of $\Delta H_{j,k}$
		can replace the low-rank truncation at comparable cost; we compare both in
		Section~\ref{ssec:commexp}.
		\item \textbf{Lazy Hessian updates ($K_{\mathrm{lazy}}$).}
		The Hessian tracker update (step~D of Algorithm~\ref{alg:disgrem}) can be
		performed every $K_{\mathrm{lazy}}$-th iteration instead of every iteration,
		reusing the previous Hessian for $K_{\mathrm{lazy}}{-}1$ consecutive steps.
		This reduces communication (no Hessian payload in skipped rounds) at the cost
		of stale curvature information.
		Section~\ref{ssec:commexp} studies the trade-off and recommends
		$K_{\mathrm{lazy}}\in\{5,10\}$ for slowly varying problems.
	\end{itemize}
	
	\noindent
	Section~\ref{sec:experiments} tests these guidelines and all four
	variants empirically, with convergence benchmarks on nine test functions
	(Section~\ref{ssec:convergence}), a communication-cost ablation study
	of compression and lazy Hessian updates (Section~\ref{ssec:commexp}),
	and an adaptive-mechanism study quantifying the robustness of
	\textsc{AdaDisGrem} to the initial $M$ (Section~\ref{ssec:ada}).
	
	%%%%%%%%%%%%%%%%%%%%%%%%%%%%%%%%%%%%%%%%%%%%%%%%%%%%%%%%%%%%
	
\section{Theory and Convergence Analysis}\label{sec:theory}
%%%%%%%%%%%%%%%%%%%%%%%%%%%%%%%%%%%%%%%%%%%%%%%%%%%%%%%%%%%%
All results in this section assume Assumption~\ref{ass:standing};
full proofs appear in Appendix~\ref{app:proofs}, and the explicit
dispersion/burn-in construction is detailed in
Appendix~\ref{app:dispersion}.

\noindent\textbf{Proof outline.}
The convergence argument proceeds through the following chain of
reductions:
\begin{enumerate}[leftmargin=1.8em,itemsep=1pt,label=\textbf{(\roman*)}]
	\item \emph{Averaging identities}
	  (\S\ref{ssec:avg})---show that the average iterate evolves as
	  $\bar x_{k+1}=\bar x_k+\bar s_k$ and that tracked averages
	  equal true averages.
	\item \emph{Reference step and bridge bounds}
	  (\S\ref{ssec:refstep}--\S\ref{ssec:bridge})---introduce a
	  ``centralized reference'' step $s_k^{\mathrm{ref}}$ and bound
	  its gap from $\bar s_k$ in terms of tracker dispersions.
	\item \emph{Stabilizer and step-dispersion control}
	  (\S\ref{ssec:stab}--\S\ref{ssec:stepdisp})---relate the
	  eigenvalue shift $\delta_{i,k}$ and the step dispersion to
	  gradient/Hessian tracker dispersions.
	\item \emph{Inexact regularized Newton condition}
	  (\S\ref{ssec:inexact})---combine the above to verify that the
	  average iterate satisfies an inexact Newton residual bound.
	\item \emph{Descent and complexity}
	  (\S\ref{ssec:descent})---derive a $3/2$-recursion on the
	  optimality gap and conclude the $\cO(\varepsilon^{-1})$ rate.
\end{enumerate}
The reader interested only in the final result may skip to
Theorem~\ref{thm:main} (p.~\pageref{thm:main}) and the subsequent
discussion of strong convexity and fixed mixing depth.

\subsection{Basic inequalities}\label{ssec:basic}
\begin{lemma}[Third-order Taylor bound]\label{lem:taylor3}
Let $h:\R^d\to\R$ have $L_2$-Lipschitz Hessian. Then for all $x,s$,
\[
h(x+s)\le h(x)+\ip{\nabla h(x)}{s}+\frac12 s^\top\nabla^2 h(x)s+\frac{L_2}{6}\norm{s}^3.
\]
\end{lemma}

\begin{lemma}[Consensus contraction]\label{lem:consensus}
For any $z\in\R^N$ and integer $t\ge 1$,
\[
\normtwo{(W^t-\J)z}\le \rho^t\normtwo{z}.
\]
Consequently, for any stacked $Z\in\R^{Nd}$,
\[
\norm{((W^t-\J)\otimes I_d)Z}\le \rho^t\norm{Z}.
\]
\end{lemma}

\begin{lemma}[Smooth convex gap--gradient inequality]\label{lem:gap-grad}
Let $h$ be convex with $L$-Lipschitz gradient. Then for any minimizer $x_\star$,
\[
\norm{\nabla h(x)}^2\le 2L\big(h(x)-h(x_\star)\big).
\]
\end{lemma}

\subsection{Averaging identities}\label{ssec:avg}
The starting point of the analysis is that double-stochastic mixing
preserves averages, so the mean iterate and mean trackers can be
studied as if they were computed by a single ``virtual agent.''

Define averages $\bar x_k:=\frac1N\sum_{i=1}^N x_{i,k}$ and similarly for other variables.
Define pre-mixed averages
\begin{gather*}
\tilde{\bar x}_k:=\tfrac1N\!\sum_i \tilde x_{i,k},\quad
\tilde{\bar g}_k:=\tfrac1N\!\sum_i \tilde g_{i,k},\\
\tilde{\bar H}_k:=\tfrac1N\!\sum_i \tilde H_{i,k}.
\end{gather*}
Define the average step $\bar s_k:=\frac1N\sum_i s_{i,k}$.

\begin{lemma}[Average preservation]\label{lem:avg-pres}
For all $k\ge0$, $\tilde{\bar x}_k=\bar x_k$ and $\bar x_{k+1}=\bar x_k+\bar s_k$.
\end{lemma}
\begin{proof}
Immediate from $W\one=\one$ (double stochasticity),
which implies $\frac1N\one^\top(W^t\otimes I_d)=\frac1N\one^\top\otimes I_d$
for any integer $t\ge 1$.
\end{proof}

\begin{lemma}[Average tracking identity for gradients]\label{lem:avg-track-g}
Assume Algorithm~\ref{alg:disgrem} is initialized by $g_{i,0}=\nabla f_i(x_{i,0})$.
Then for all $k\ge0$,
\[
\bar g_k:=\frac1N\sum_{i=1}^N g_{i,k}=\frac1N\sum_{i=1}^N\nabla f_i(x_{i,k}),
\qquad
\tilde{\bar g}_k=\bar g_k.
\]
\end{lemma}
\begin{proof}
By induction using $W\one=\one$ and the telescoping form of the tracker update (step~D).
\end{proof}

\begin{lemma}[Average tracking identity for Hessians]\label{lem:avg-track-h}
Assume Algorithm~\ref{alg:disgrem} is initialized by $H_{i,0}=\nabla^2 f_i(x_{i,0})$.
Then for all $k\ge0$,
\[
\bar H_k:=\frac1N\sum_{i=1}^N H_{i,k}=\frac1N\sum_{i=1}^N\nabla^2 f_i(x_{i,k}),
\qquad
\tilde{\bar H}_k=\bar H_k.
\]
\end{lemma}
\begin{proof}
Identical argument to Lemma~\ref{lem:avg-track-g}, applied to the Hessian tracker.
\end{proof}

\subsection{Well-posed local solves and a reference step}\label{ssec:refstep}
With the averaging identities in hand, we next introduce the
\emph{reference step} $s_k^{\mathrm{ref}}$: the Newton step that a
centralized agent would compute using the averaged Hessian and gradient.
The gap $\|\bar s_k - s_k^{\mathrm{ref}}\|$ then quantifies how much the
decentralized updates deviate from this ideal step.

Algorithm~\ref{alg:disgrem} sets
\begin{gather*}
\lambda_{i,k}:=\sqrt{M\|\tilde g_{i,k}\|},\quad
\delta_{i,k}:=\max\{0,-\lambda_{\min}(\mathrm{sym}(\tilde H_{i,k}))\},\\
\tilde\lambda_{i,k}:=\lambda_{i,k}+\delta_{i,k}.
\end{gather*}
Define
\[
A_{i,k}:=\mathrm{sym}(\tilde H_{i,k})+\tilde\lambda_{i,k}I,
\qquad
s_{i,k}:=-A_{i,k}^{-1}\tilde g_{i,k}.
\]
Define averaged quantities
\begin{gather*}
\tilde{\bar\lambda}_k:=\tfrac1N\textstyle\sum_{i=1}^N \tilde\lambda_{i,k},\quad
A_k^{\mathrm{ref}}:=\mathrm{sym}(\tilde{\bar H}_k)+\tilde{\bar\lambda}_k I,\\
s_k^{\mathrm{ref}}:=-(A_k^{\mathrm{ref}})^{-1}\tilde{\bar g}_k.
\end{gather*}

\begin{lemma}[Local cubic regime: $M$ controls stepsizes]\label{lem:M-vs-s}
For all $i,k$, the local step satisfies
\[
M\norm{s_{i,k}}\le \lambda_{i,k}
\qquad\text{and hence}\qquad
L_2\norm{s_{i,k}}\le \lambda_{i,k}.
\]
\end{lemma}

\begin{lemma}[Average cubic regime]\label{lem:avg-regime}
For all $k$,
\[
L_2\|\bar s_k\|\le \frac1N\sum_{i=1}^N \lambda_{i,k}\le \tilde{\bar\lambda}_k.
\]
\end{lemma}

\subsection{Bridge bounds: tracked averages vs.\ true quantities}\label{ssec:bridge}
The reference step uses the tracked averages
$\tilde{\bar g}_k,\tilde{\bar H}_k$ rather than the true gradient
$g_k:=\nabla f(\bar x_k)$ and true Hessian
$\nabla^2 f(\bar x_k)$.
By Lemma~\ref{lem:avg-track-g},
$\tilde{\bar g}_k = \frac{1}{N}\sum_{i}\nabla f_i(x_{i,k})$,
which equals $\nabla f(\bar x_k)$ \emph{only at exact consensus}
($x_{i,k}=\bar x_k$ for all~$i$).
Away from consensus, the discrepancy is controlled by the
``bridge'' lemmas below, which quantify it in terms of the
post-mixing spatial disagreement $D_X(k)$.
All constants depending on the bounded sublevel set
(e.g.\ $L_1,L_2$) are evaluated on the invariant set
established by Lemma~\ref{lem:lyap_bound}, so no circularity arises.

Define the post-mixing disagreement (RMS)
\[
D_X(k):=\sqrt{\frac1N\sum_{i=1}^N \|x_{i,k}-\bar x_k\|^2}
= D(X_k),
\]
where $D(\cdot)$ is the dispersion defined in Section~\ref{sec:prelim}.
We use $D_X(k)$ when emphasizing the iterate index and $D(X_k)$ elsewhere.

\begin{lemma}[Gradient bridge (average tracked gradient vs.\ true average gradient)]\label{lem:bridge-grad}
For every $k$,
\[
\|\nabla f(\bar x_k)-\tilde{\bar g}_k\|\le L_1\,D_X(k).
\]
\end{lemma}

\begin{lemma}[Hessian bridge (average tracked Hessian vs.\ true average Hessian)]\label{lem:bridge-hess}
For every $k$,
\[
\|\nabla^2 f(\bar x_k)-\mathrm{sym}(\tilde{\bar H}_k)\|_2\le L_2\,D_X(k).
\]
\end{lemma}

\subsection{The eigenvalue-shift stabilizer \texorpdfstring{$\delta_{i,k}$}{delta\_ik}}\label{ssec:stab}
The stabilizer $\delta_{i,k}$ ensures that every local linear system
is positive definite.
Because it acts as an additive bias in $\tilde\lambda_{i,k}$, we need
to show that its average $\bar\delta_k$ is controlled by quantities
that vanish as consensus improves.

Define the stabilizer average
\[
\bar\delta_k:=\frac1N\sum_{i=1}^N \delta_{i,k}.
\]

\begin{lemma}[Stabilizer controlled by consensus/tracking dispersion]\label{lem:delta_control}
For every $k$ and every $i$,
\[
0\le \delta_{i,k}\le \big\|\mathrm{sym}(\tilde H_{i,k})-\nabla^2 f(\bar x_k)\big\|_2.
\]
Consequently,
\begin{gather*}
\bar\delta_k\le \Delta_k^H + L_2D_X(k),\\
\Delta_k^H:=\tfrac1N\!\sum_{i=1}^N\big\|\mathrm{sym}(\tilde H_{i,k})-\mathrm{sym}(\tilde{\bar H}_k)\big\|_2 .
\end{gather*}
\end{lemma}

\subsection{Step dispersion via a resolvent identity}\label{ssec:stepdisp}
The inexact Newton condition (Proposition~\ref{prop:inexact-rn} below)
requires three quantities to be small: the step dispersion
$\|\bar s_k - s_k^{\mathrm{ref}}\|$, the gradient bridge error
$L_1 D_X(k)$, and the Hessian bridge error $L_2 D_X(k)$.
The latter two are already controlled by
Lemmas~\ref{lem:bridge-grad}--\ref{lem:bridge-hess};
this subsection handles the first.
The key tool is a resolvent identity that factorizes the
difference of two linear-system solutions through their
coefficient matrices.

Define dispersions
\begin{gather*}
\Delta_k^g:=\tfrac1N\textstyle\sum_{i=1}^N \|\tilde g_{i,k}-\tilde{\bar g}_k\|,\quad
\Delta_k^\lambda:=\tfrac1N\textstyle\sum_{i=1}^N |\tilde\lambda_{i,k}-\tilde{\bar\lambda}_k|,\\
\underline\lambda_k:=\min_{1\le i\le N}\tilde\lambda_{i,k}.
\end{gather*}

\begin{lemma}[Step dispersion]\label{lem:step-disp}
For every $k$ with $\underline\lambda_k>0$,
\[
\|\bar s_k-s_k^{\mathrm{ref}}\|
\le \frac{1}{\underline\lambda_k}\Delta_k^g
+\frac{\|\tilde{\bar g}_k\|}{\underline\lambda_k\,\tilde{\bar\lambda}_k}\big(\Delta_k^H+\Delta_k^\lambda\big).
\]
\end{lemma}

\begin{lemma}[Dispersion of $\tilde\lambda_{i,k}=\sqrt{M\|\tilde g_{i,k}\|}+\delta_{i,k}$]\label{lem:lambda_dispersion}
Assume $\tilde{\bar g}_k\neq 0$ and the \emph{relative gradient dispersion} condition
\begin{equation}\label{eq:rel_grad_disp}
\max_{1\le i\le N}\|\tilde g_{i,k}-\tilde{\bar g}_k\|\le \alpha_d\|\tilde{\bar g}_k\|
\qquad\text{for some }\alpha_d\in(0,1)
\end{equation}
(this condition is verified for all $k\ge K_0(\varepsilon)$ by the
burn-in analysis in Proposition~\ref{prop:burnin}).
Then
\begin{equation}\label{eq:lambda_disp_bound}
\Delta_k^\lambda
\le
\frac{\sqrt{M}}{\sqrt{1-\alpha_d}}\cdot\frac{\Delta_k^g}{\sqrt{\|\tilde{\bar g}_k\|}}
+2\bar\delta_k.
\end{equation}
\end{lemma}

% ============================================================
\subsection{Inexact regularized Newton condition for the average step}\label{ssec:inexact}
% ============================================================
Collecting the bridge, stabilizer, and step-dispersion bounds, we can
now verify that the average iterate $\bar x_k$ satisfies a standard
\emph{inexact regularized Newton} condition once the burn-in phase
is complete (i.e., for $k\ge K_0(\varepsilon)$, where $K_0$ is the
burn-in index from Proposition~\ref{prop:burnin}).
This is the key step: it allows us to invoke the one-step descent
lemma from the centralized analysis.

Define
\begin{gather*}
g_k:=\nabla f(\bar x_k),\quad
\lambda_k:=\tilde{\bar\lambda}_k,\\
r_k:=(\nabla^2 f(\bar x_k)+\lambda_k I)\bar s_k+g_k.
\end{gather*}

\begin{proposition}[Inexact RN condition]\label{prop:inexact-rn}
Assume the following hold at iteration $k$ for some $\eta\in(0,1)$:
\begin{enumerate}[leftmargin=1.2em,itemsep=2pt]
\item \textbf{Dispersion control:}
\[
\|\bar s_k-s_k^{\mathrm{ref}}\|\le \frac{\eta}{8}\cdot \frac{\lambda_k}{M_{H,\max}+\lambda_k}\,\|\bar s_k\|.
\]
\item \textbf{Gradient bridge accuracy:}
\[
L_1D_X(k)\le \frac{\eta}{8}\lambda_k\|\bar s_k\|.
\]
\item \textbf{Hessian bridge accuracy:}
\[
L_2D_X(k)\le \frac{\eta}{8}\lambda_k.
\]
\end{enumerate}
Then
\[
\|r_k\|\le \eta\,\lambda_k\|\bar s_k\|,
\qquad
L_2\|\bar s_k\|\le \lambda_k.
\]
\end{proposition}

\subsection{Descent and global rates}\label{ssec:descent}
Steps~(i)--(iv) of the proof pipeline are now complete: we have shown
that the average iterate satisfies an inexact regularized Newton
condition (Proposition~\ref{prop:inexact-rn}) once the burn-in phase
ends.
The remaining step~(v) follows the centralized analysis of
\citet{mishchenko2023rn}: a one-step descent lemma yields a
$3/2$-power recursion on the optimality gap, which is then solved to
give the global complexity bound.

Define $\Phi_k:=f(\bar x_k)-f_\star$.

\begin{lemma}[One-step descent]\label{lem:descent}
Assume
\[
\|(\nabla^2 f(\bar x_k)+\lambda_k I)\bar s_k+g_k\|\le \eta\lambda_k\|\bar s_k\|,
\qquad
L_2\|\bar s_k\|\le \lambda_k,
\]
for some $\eta\in(0,1)$.
Then
\[
f(\bar x_{k+1})\le f(\bar x_k) - \Big(1-\eta-\frac{1}{6}\Big)\lambda_k\|\bar s_k\|^2.
\]
\end{lemma}

\begin{lemma}[Gradient bound]\label{lem:grad_bound}
Under the conditions of Lemma~\ref{lem:descent},
\[
\|g_{k+1}\|\le \Big(1+\eta+\frac{1}{2}\Big)\lambda_k\|\bar s_k\|.
\]
\end{lemma}

Define
\begin{gather*}
\mathcal I:=\{k\ge0:\ \|g_{k+1}\|\ge \tfrac14\|g_k\|\},\\
\mathcal S:=\{k\ge0:\ \|g_{k+1}\|< \tfrac14\|g_k\|\}.
\end{gather*}

\begin{lemma}[Steady indices give a $3/2$ recursion]\label{lem:steady}
Assume Proposition~\ref{prop:inexact-rn} holds for all $k$ with some $\eta\le 1/12$.
Assume additionally that there exists $C_\lambda>0$ such that
\[
\lambda_k\le C_\lambda\sqrt{\|g_k\|}\qquad\text{for all $k$ with $g_k\neq 0$}.
\]
(This condition is ensured by Lemma~\ref{lem:lambda_scaling}
in Appendix~\ref{app:proofs} under
schedule~\eqref{eq:log_schedule_main}, together with the burn-in
analysis of Proposition~\ref{prop:burnin}.)
Then there exists $\tau>0$ such that for all $k\in\mathcal I$,
\[
\Phi_k-\Phi_{k+1}\ge \tau\,\Phi_k^{3/2}.
\]
\end{lemma}

\begin{lemma}[Solving the $3/2$ recursion]\label{lem:solve32}
If $\Phi_{k+1}\le \Phi_k-\tau\Phi_k^{3/2}$ for some $\tau>0$ and all $k$, then
\[
\Phi_k\le \frac{4}{\tau^2(k+2)^2}.
\]
\end{lemma}

Throughout we use the \emph{logarithmic mixing schedule}
\begin{equation}\label{eq:log_schedule_main}
  \tau_k \;=\; t_k \;=\; \Bigl\lceil \frac{p\,\log(k+2)}{-\log\rho} \Bigr\rceil,
  \qquad k\ge 0,
\end{equation}
so that $\rho^{\tau_k}\le(k+2)^{-p}$; details and motivation are given in
Appendix~\ref{app:dispersion}.

\begin{theorem}[Global iteration complexity]\label{thm:main}
Assume Assumption~\ref{ass:standing}.
Fix $p>2$ and run \textsc{DisGrem} with the logarithmic mixing schedule~\eqref{eq:log_schedule_main}.
Then for every target accuracy $\varepsilon>0$, there exist a finite burn-in index $K_0(\varepsilon)$
(depending on the spectral gap, Lipschitz constants, initial dispersions, and $\varepsilon$;
see Proposition~\ref{prop:burnin})
and a constant $C_K>0$ (independent of $\varepsilon$) such that the total number of
outer iterations to achieve $\|\nabla f(\bar x_k)\|\le \varepsilon$ satisfies
\[
K(\varepsilon)\;\le\; K_0(\varepsilon) + C_K\,\varepsilon^{-1}.
\]
In particular:
\begin{enumerate}[leftmargin=1.2em,itemsep=2pt]
\item The post-burn-in phase requires at most $\cO(\varepsilon^{-1})$ iterations,
matching the centralized regularized Newton complexity of \citet{mishchenko2023rn}.
\item Under the schedule~\eqref{eq:log_schedule_main}, the total number of neighbor communication rounds is
$\cO\big((K_0(\varepsilon)+\varepsilon^{-1})\log(K_0(\varepsilon)+\varepsilon^{-1})\big)$.
\item The burn-in $K_0(\varepsilon)$ is the smallest index at which
polynomial-decaying tracker dispersions
(Proposition~\ref{prop:dispersion_decay}) become small enough to satisfy the
inexact regularized Newton condition (Proposition~\ref{prop:inexact-rn}).
Since the dispersions decay as $\cO(k^{-(p-1)})$, inverting the decay gives
$K_0(\varepsilon)=\cO\!\bigl(\varepsilon^{-1/(2(p-1))}\bigr)$,
which is $o(\varepsilon^{-1})$ for $p>2$; the total complexity is therefore
dominated by the post-burn-in $\cO(\varepsilon^{-1})$ term.
\end{enumerate}
\end{theorem}

\begin{remark}[Burn-in order and network dependence]\label{rem:burnin_order}
The dispersions decay as $\cO(k^{-(p-1)})$
(Proposition~\ref{prop:dispersion_decay}), and the inexact Newton
conditions (Proposition~\ref{prop:inexact-rn}) require them to be at
most $\cO(\sqrt\varepsilon)$.
Inverting the polynomial decay gives
$K_0(\varepsilon)=\cO\!\bigl(\varepsilon^{-1/(2(p-1))}\bigr)$.
For the default choice $p>2$, the exponent $\frac{1}{2(p-1)}<1$, so
$K_0(\varepsilon)=o(\varepsilon^{-1})$ and the total complexity
$K_0(\varepsilon)+\cO(\varepsilon^{-1})$ is dominated by the
$\cO(\varepsilon^{-1})$ post-burn-in term.
For instance, $p=3$ yields
$K_0(\varepsilon)=\cO(\varepsilon^{-1/4})$.

The dependence on the spectral gap enters through the mixing schedule:
$\tau_k=\lceil p\log(k{+}2)/({-\log\rho})\rceil$ implies that each
outer iteration uses $\cO\!\bigl((1{-}\rho)^{-1}\log k\bigr)$
communication rounds (since $-1/\log\rho\approx(1{-}\rho)^{-1}$ for
$\rho$ near~1).
Consequently, the total communication-round budget is
$\cO\!\bigl((1{-}\rho)^{-1}\,\varepsilon^{-1}\log^2\varepsilon^{-1}\bigr)$
(the extra $\log$ factor relative to item~2 of Theorem~\ref{thm:main}
accounts for the $(1{-}\rho)^{-1}$ per-iteration cost, which the
theorem statement absorbs into the round count),
exhibiting a \emph{linear} penalty in $(1{-}\rho)^{-1}$---the same
order as first-order gradient-tracking methods
\citep{nedic2017diging}.
On very sparse graphs ($\rho\to 1$), the per-iteration communication
overhead grows, but the iteration count remains network-independent.
\end{remark}

\begin{remark}[Proof structure: two self-contained stages]
The proof of Theorem~\ref{thm:main} combines two independent
analyses, distinct from the five-step pipeline~(i)--(v) above
(which describes the \emph{per-iteration} reduction chain used
in Stage~2).
\emph{Stage~1---dispersion decay}
(Proposition~\ref{prop:dispersion_decay},
Appendix~\ref{app:dispersion}): under the logarithmic schedule,
tracker dispersions decay polynomially; this requires only
bounded iterates (Lemma~\ref{lem:lyap_bound}), not any gradient
or function-value condition.
\emph{Stage~2---descent}: once dispersions are small enough
(after $K_0(\varepsilon)$ iterations), the per-iteration chain
(i)--(v) verifies an inexact regularized Newton recursion
(Proposition~\ref{prop:inexact-rn}), and the $3/2$-recursion
(Lemma~\ref{lem:solve32}) gives $\cO(\varepsilon^{-1})$ further
iterations.
The two stages are self-contained; there is no circularity.
A fully explicit version of Theorem~\ref{thm:main} appears as
Theorem~\ref{thm:end2end} in Appendix~\ref{app:dispersion}.
\end{remark}

\subsection{Strong convexity and local superlinear convergence}
The results above require only convexity.
When the objective is locally strongly convex near the optimum,
we can further show that the stabilizer term $\bar\delta_k$ becomes
negligible relative to $\|g_k\|$, recovering the exact Newton
regime and, with it, local $Q$-superlinear convergence.

\begin{assumption}\label{ass:strong}
In addition to Assumption~\ref{ass:standing}, $f$ is $\mu$-strongly convex on the bounded set
$\{x:\|x-x_\star\|\le D\}$ (where $D$ is given by Lemma~\ref{lem:lyap_bound}):
$\nabla^2 f(x)\succeq \mu I$ for all $\|x-x_\star\|\le D$.
\end{assumption}

\begin{lemma}[Stabilizer becomes negligible relative to $\|g_k\|$]\label{lem:delta_o_g}
Assume Assumptions~\ref{ass:standing} and~\ref{ass:strong}.
Fix any exponent $\gamma>0$.
Let
\[
G_g:=\sup_{k\ge 0}\ \max_{1\le i\le N}\ \|\tilde g_{i,k}\|\in(0,\infty),
\]
which is finite by Lemma~\ref{lem:tilde_track_bounded}.
Suppose that for all $k$ large enough (and $k\ge 1$) the mixing depths satisfy
\begin{equation}\label{eq:delta_schedule}
\max\{\rho^{\tau_k},\rho^{t_{k-1}}\}\le \tfrac{1}{C_\delta}\|g_k\|^{1+\gamma},
\end{equation}
where $C_\delta:=\hat C_H+L_2(2D+\sqrt{G_g/M})$ and $\hat C_H:=B_H+2\sqrt{d}\,L_2D$,
where $B_H:=\sup_{k\ge0}D_H(\tilde{\mathcal H}_k)<\infty$
is the uniform Hessian-tracker dispersion bound from Lemma~\ref{lem:tilde_hess_bounded}.
Then, for all sufficiently large $k$,
\[
\bar\delta_k = o(\|g_k\|)\qquad\text{and more precisely}\qquad \bar\delta_k\le \|g_k\|^{1+\gamma}.
\]
\end{lemma}

\begin{theorem}[Local $Q$-superlinear ($3/2$-order) gradient rate]\label{thm:superlinear}
Assume Assumptions~\ref{ass:standing} and~\ref{ass:strong}.
Assume that Proposition~\ref{prop:inexact-rn} holds for all sufficiently large $k$ with some $\eta\in(0,1/12]$
(guaranteed by Theorem~\ref{thm:main} for $k\ge K_0(\varepsilon)$),
and that the mixing depth $\tau_k$ is chosen so that \eqref{eq:delta_schedule} holds with some $\gamma>0$.
Then there exist constants $C_{\mathrm{sc}}>0$ and $k_{\mathrm{sc}}$ such that for all $k\ge k_{\mathrm{sc}}$,
\[
\|\nabla f(\bar x_{k+1})\|\le C_{\mathrm{sc}}\,\|\nabla f(\bar x_k)\|^{3/2}.
\]
\end{theorem}

\subsection{Practical regime: fixed mixing depth}\label{ssec:fixed_nc}

The logarithmic schedule~\eqref{eq:log_schedule_main} yields exact convergence but requires
$\tau_k,t_k\to\infty$.
In practice, one fixes the mixing depth $\tau_k=t_k=N_c$ for all~$k$.
The following corollary characterizes the resulting convergence behavior.

\begin{corollary}[Neighborhood convergence with fixed mixing depth]\label{cor:fixed_nc}
Assume Assumption~\ref{ass:standing} and fix a constant mixing depth $\tau_k=t_k=N_c\ge 1$ for all $k$.
Define the residual consensus error
\[
\epsilon_{\mathrm{mix}} \;:=\; \rho^{N_c}\in(0,1).
\]
Then under the same convexity and Lipschitz conditions as Theorem~\ref{thm:main}:
\begin{enumerate}[leftmargin=1.2em,itemsep=2pt]
\item The tracker dispersions stabilize at
\[
\limsup_{k\to\infty} D(\tilde G_k)
\le \tfrac{\epsilon_{\mathrm{mix}}}{1-\epsilon_{\mathrm{mix}}^2}2L_1D
=: \bar\Delta_G,
\]
\[
\limsup_{k\to\infty} D_H(\tilde{\mathcal H}_k)
\le \tfrac{\epsilon_{\mathrm{mix}}}{1-\epsilon_{\mathrm{mix}}^2}2\sqrt{d}\,L_2D
=: \bar\Delta_H.
\]
\item There exists a neighborhood radius $\varepsilon_{\mathrm{nbhd}}=\cO(\epsilon_{\mathrm{mix}})$
(depending on $L_1,L_2,M,D$) such that
\[
\liminf_{k\to\infty}\|\nabla f(\bar x_k)\|\;\le\; \varepsilon_{\mathrm{nbhd}}.
\]
Specifically, for every $\varepsilon>\varepsilon_{\mathrm{nbhd}}$, the iterate reaches
$\|\nabla f(\bar x_k)\|\le \varepsilon$ in $\cO(\varepsilon^{-1})$ iterations.
\item When $\varepsilon_{\mathrm{nbhd}}\ll\varepsilon$
(equivalently, $N_c\gg \log(\varepsilon^{-1})/\log(\rho^{-1})$),
the fixed-$N_c$ regime is indistinguishable from the exact-convergence regime
for all practical purposes.
\end{enumerate}
\end{corollary}

\begin{proof}
Item~1 follows by setting $a_k=\rho^{2N_c}=\epsilon_{\mathrm{mix}}^2$ in the
dispersion recursions (Lemmas~\ref{lem:DG_rec}--\ref{lem:DH_rec}) and solving the
resulting geometric series.
Item~2 follows because the residual dispersions $\bar\Delta_G, \bar\Delta_H$
prevent the three conditions of Proposition~\ref{prop:inexact-rn} from being
satisfied below $\varepsilon_{\mathrm{nbhd}}$,
where $\varepsilon_{\mathrm{nbhd}}$ is determined by inverting the burn-in
conditions (Proposition~\ref{prop:burnin}) with
$D_X(\infty)\le C_X^{+}\epsilon_{\mathrm{mix}}/(1-\epsilon_{\mathrm{mix}})$
in place of zero.
Item~3 is immediate from comparing $\varepsilon_{\mathrm{nbhd}}=\cO(\epsilon_{\mathrm{mix}})$
with the target~$\varepsilon$.
\end{proof}

\begin{remark}[Interpretation of $\liminf$ and min-gradient form]
The $\liminf$ form in item~2 guarantees infinitely many visits to
the $\varepsilon_{\mathrm{nbhd}}$-ball; an equivalent statement is
$\min_{0\le j\le k}\|\nabla f(\bar x_j)\|
=\cO(\max(\varepsilon_{\mathrm{nbhd}}, k^{-1}))$,
the form commonly reported in non-convex complexity analyses.
\end{remark}

\begin{remark}[Conservatism of the neighborhood bound and practical recommendations]\label{rem:practical_implications}
With $\rho\approx 0.92$ and $N_c=3$ the worst-case bound gives
$\epsilon_{\mathrm{mix}}=\rho^3\approx 0.78$, yet fixed $N_c{=}3$
reaches relF~$\lesssim 10^{-10}$ on several benchmarks
(Section~\ref{sec:experiments}).
The gap arises because Corollary~\ref{cor:fixed_nc} replaces the
per-iteration driving term $L_1\|X_{k+1}{-}X_k\|$ in the
dispersion recursion (Lemma~\ref{lem:DG_rec}) with a \emph{uniform}
upper bound, whereas in practice
$\|X_{k+1}{-}X_k\|\to 0$ as stationarity is approached, so the
effective contraction is much tighter than~$\rho^{2N_c}$.
Formalizing this self-reinforcing effect is an open question.
Increasing $N_c$ to $5$--$8$ further tightens the neighborhood at
modest communication cost.
The logarithmic schedule~\eqref{eq:log_schedule_main} guarantees
exact convergence but prescribes $\tau_k=\cO(\log k/(1{-}\rho))$
rounds per iteration---e.g.\ ${\approx}75$ at $k{=}500$---which is
far more expensive than the fixed $N_c{=}3$ used in all experiments.
We recommend fixed $N_c$ for implementation and view the
logarithmic schedule primarily as a device for theoretical analysis.
\end{remark}

%%%%%%%%%%%%%%%%%%%%%%%%%%%%%%%%%%%%%%%%%%%%%%%%%%%%%%%%%%%%


% ============================================================

\section{Numerical Experiments}\label{sec:experiments}
	%%%%%%%%%%%%%%%%%%%%%%%%%%%%%%%%%%%%%%%%%%%%%%%%%%%%%%%%%%%%
	
	The preceding theory establishes exact convergence under a
	logarithmic mixing schedule and neighborhood convergence under
	a fixed depth $N_c$.
	All experiments below use the practical fixed-$N_c$ regime;
	we test whether the theoretical advantages translate
	into concrete gains over state-of-the-art baselines.
	
	We compare the \textsc{DisGrem} family with six first- and
	second-order decentralized baselines on nine objectives that
	cover well-conditioned and ill-conditioned convex problems,
	logistic regression on real-world data, and four non-convex landscapes.
	Every experiment uses a fixed mixing depth $N_c{=}3$
	(Corollary~\ref{cor:fixed_nc}), and all numbers are averaged
	over 20 independent Monte~Carlo trials (random starting points
	and random graphs).
	
	% ─────────────────────────────────────────────────────────────────────
	\subsection{Experimental setup}\label{ssec:setup}
	% ─────────────────────────────────────────────────────────────────────
	
	\noindent\textbf{Network.}
	$N{=}10$ agents communicate over an ER random graph
	with edge probability $p{=}0.5$ (regenerated per trial).
	The Metropolis--Hastings doubly stochastic mixing matrix $W$ yields
	spectral gap $1{-}\rho\approx 0.08$
	($\rho\approx 0.92$; typical range $\rho\in[0.88,0.96]$ across 20 trials),
	representing a moderately sparse topology.
	All algorithms perform $N_c{=}3$ consensus rounds per outer iteration.
	
	\noindent\textbf{Algorithms (10 methods in three groups).}
	
	\noindent\textit{Proposed (\textsc{DisGrem} family, 4 variants)}:
	\textbf{DisGrem} (Algorithm~\ref{alg:disgrem}),
	\textbf{CeDisGrem} (communication-efficient: low-rank Hessian compression
	with default rank $r{=}\lceil d/5\rceil$, Section~\ref{ssec:variants}),
	\textbf{AdaDisGrem} (adaptive $M$ via secant-based Lipschitz estimation),
	and \textbf{CeAdaDisGrem} (adaptive + compressed).
	The regularization scaling is $M = M_{\mathrm{fac}}\cdot H_{\max}^{0}$ where
	$H_{\max}^{0} = \max_i \|\nabla^2 f_i(x_0)\|_2$;
	per-function $M_{\mathrm{fac}}$ values are listed in Table~\ref{tab:config}.
	Note that the theoretical condition $M\ge L_2$ is not explicitly
	enforced by the heuristic default $M=M_{\mathrm{fac}}\cdot H_{\max}^0$;
	the adaptive variant \textsc{AdaDisGrem} addresses this gap via
	online secant-based estimation.
	The fixed-$M$ results should therefore be viewed as demonstrating
	empirical robustness beyond the strict theoretical regime.
	
	\noindent\textit{First-order baselines (2)}:
	\textbf{EXTRA}~\citep{shi2015extra} and
	\textbf{DIGing}~\citep{nedic2017diging}.
	
	\noindent\textit{Second-order baselines (4)}:
	\textbf{DQM}~\citep{eisen2017dqn},
	\textbf{ESOM}~\citep{mokhtari2016esom},
	\textbf{SONATA}~\citep{sun2022sonata}, and
	\textbf{Network-GIANT}~\citep{ye2023networkgiant}.
	First-order baselines use a Lipschitz-scaled stepsize
	$\alpha=\alpha_{\mathrm{base}}/H_{\max}^{0}$;
	for each function, $\alpha_{\mathrm{base}}$ is selected from
	$\{0.01, 0.1, 0.5, 1.0\}$ as the value giving the fastest
	convergence without divergence over 5 preliminary runs
	(Table~\ref{tab:config}).
	Second-order baselines (DQM, ESOM, SONATA, Network-GIANT) use their
	published default parameters; SONATA and Network-GIANT solve local
	subproblems to machine precision via a direct solver.
	All baselines use the same $N_c{=}3$ consensus rounds per outer
	iteration as the \textsc{DisGrem} family.
	Code for all experiments will be made publicly available upon
	publication.
	
	\noindent\textbf{Test objectives (9 functions).}
	All synthetic functions use $d{=}30$; the two logistic regression problems
	inherit $d{=}22$ from the \textsc{svmguide3} dataset.
	See Table~\ref{tab:config} for full parameter specifications.
	\textit{Convex} (5):
	(i)~\textsc{Ridge} ($\ell_2$-regularized quadratic, $\lambda{=}10^{-3}$);
	(ii)~\textsc{QuadBad} (ill-conditioned quadratic, $\kappa{=}10^3$);
	(iii)~\textsc{LogSumExp};
	(iv)~\textsc{Huber} (smooth pseudo-Huber loss, $\delta{=}1$);
	(v)~\textsc{LogReg-real} ($\ell_2$-regularized logistic regression on the
	\textsc{svmguide3} LibSVM dataset~\citep{chang2011libsvm}, $m{=}1243$, $d{=}22$).
	\textit{Non-convex} (4; \textbf{not covered by the theory}---included
	to probe algorithmic robustness beyond the convex regime):
	(vi)~\textsc{LinLog} ($\sum\ln(1{+}|r_i|)$, with flat curvature regions);
	(vii)~\textsc{Rosenbrock};
	(viii)~\textsc{Styblinski--Tang} (multimodal);
	(ix)~\textsc{LogReg-NCVR} (non-convex regularized logistic regression on
	\textsc{svmguide3}).
	On non-convex objectives the eigenvalue-shift stabilizer $\delta_{i,k}$
	ensures the local linear system remains SPD, so the algorithm operates
	as a damped/stabilized Newton-like method; however, the convergence
	guarantees of Theorems~\ref{thm:main}--\ref{thm:superlinear} do not apply.
	We note that $\delta_{i,k}$ serves two structurally similar but
	theoretically distinct roles:
	in the convex analysis it compensates for transient indefiniteness of
	tracked Hessian matrices (which are PSD in the exact case);
	on non-convex objectives it additionally absorbs true negative curvature,
	effectively acting as a Levenberg--Marquardt damping term.
	A unified non-convex convergence analysis is left for future work.
	For convex objectives, $f_{\mathrm{ref}}:=f^\star$ (the global minimum,
	whose existence is guaranteed by coercivity).
	For non-convex objectives, $f_{\mathrm{ref}}$ denotes the best value
	found by multi-start L-BFGS-B (50 restarts, tolerance $10^{-15}$);
	it is \emph{not} a certified global optimum.
	
	\begin{table*}[t]
		\centering
		\caption{Per-function experimental configuration.
			$M_{\mathrm{fac}}$: regularization scaling for \textsc{DisGrem}
			($M = M_{\mathrm{fac}}\cdot H_{\max}^{0}$);
			$\alpha_{\mathrm{base}}$: stepsize coefficient for first-order
			baselines (actual step $\alpha_{\mathrm{base}}/H_{\max}^{0}$);
			$K_{\max}$: iteration budget (shared by all algorithms);
			Decay: whether a $1/\sqrt{k}$ stepsize decay is applied to
			first-order baselines (relevant only for non-convex problems).}
		\label{tab:config}
		\setlength{\tabcolsep}{3pt}
		\footnotesize
		\begin{tabular}{lccccl}
			\toprule
			Function & $M_{\mathrm{fac}}$ & $\alpha_{\mathrm{base}}$ & $K_{\max}$ & Decay & Notes \\
			\midrule
			Ridge       & 0.1  & 0.20 &  200 & No  & $\lambda{=}10^{-3}$ \\
			QuadBad     & 0.1  & 0.10 & 1500 & No  & $\kappa{=}10^3$ \\
			LogSumExp   & 5.0  & 0.30 &  400 & No  & \\
			Huber       & 1.5  & 0.30 &  800 & No  & Pseudo-Huber, $\delta{=}1$ \\
			LogReg-real & 3.0  & 1.00 &  600 & No  & svmguide3, $d{=}22$ \\
			LinLog      & 1.0  & 0.20 & 1500 & No  & Non-convex \\
			Rosenbrock  & 3.0  & 0.10 &  300 & Yes & Non-convex \\
			Styblinski  & 15.0 & 0.05 &  100 & Yes & Multimodal \\
			LogReg-NCVR & 3.0  & 1.00 & 1000 & Yes & svmguide3, $d{=}22$ \\
			\bottomrule
		\end{tabular}
	\end{table*}
	
	\noindent\textbf{Metrics and reporting.}
	Define the consensus residual
	\[
	  \mathrm{cons}_k
	  \;:=\;
	  D(X_k)
	  \;=\;
	  \sqrt{\tfrac{1}{N}\textstyle\sum_{i=1}^{N}\|x_{i,k}-\bar x_k\|^{2}},
	\]
	and the \emph{combo} stopping criterion
	\[
	  \mathrm{combo}_k
	  \;:=\;
	  \|\nabla f(\bar x_k)\| + \mathrm{cons}_k,
	\]
	where $\nabla f(\bar x_k)=\frac{1}{N}\sum_{i}\nabla f_i(\bar x_k)$
	is the true gradient at the average iterate (computed offline,
	not from the tracker), ensuring a fair comparison across algorithms
	with different internal tracking mechanisms.
	Each algorithm runs to $K_{\max}$ iterations or until
	$\mathrm{combo}_k < 10^{-12}$.
	We report four metrics:
	(a)~$\mathrm{combo}=\min_k\,\mathrm{combo}_k$;
	(b)~$\mathrm{relF}=\min_k |f(\bar x_k)-f_{\mathrm{ref}}|/|f(\bar x_0)-f_{\mathrm{ref}}|$;
	(c)~wall-clock time (seconds; for reference only---all implementations
	use Python/NumPy with comparable vectorization, but we did not
	optimize any method's code to the same degree);
	(d)~cumulative communication cost (MB), defined as the total
	send\,+\,receive volume per agent, summed over all outer iterations and all
	neighbors (consistent with the per-iteration estimate in Section~\ref{ssec:comm}).
	All convergence curves show the 20-run mean with $\pm1$ standard
	deviation shading.
	Analytical gradients and Hessians are used for all functions.
	Because \textsc{DisGrem} transmits $\cO(d^2)$ Hessian data per
	iteration while first-order methods transmit $\cO(d)$ vectors,
	a direct per-iteration communication comparison would be
	unfair.
	We therefore compare total communication volume \emph{to reach
	a given precision $\epsilon$} (the performance-profile threshold),
	which is the relevant metric in communication-limited settings.
	A run is deemed \emph{successful} at threshold $\epsilon$ if it reaches
	$\mathrm{relF}\le\epsilon$ within $K_{\max}$ iterations without
	encountering NaN or overflow; for non-convex problems, relF is
	computed relative to $f_{\mathrm{ref}}$.
	Unless otherwise noted, the Hessian is updated every iteration
	($K_{\mathrm{lazy}}{=}1$); Section~\ref{ssec:commexp} studies the
	effect of lazy updates.
	
	\noindent\textbf{Theory--experiment gap for fixed $N_c$.}
	Corollary~\ref{cor:fixed_nc} gives a conservative
	$\cO(\rho^{N_c})$-neighborhood bound; the experiments reach
	much higher accuracy because the driving terms decay with the
	iterates (Remark~\ref{rem:practical_implications}).
	
	% ─────────────────────────────────────────────────────────────────────
	\subsection{Convergence benchmarks}\label{ssec:convergence}
	% ─────────────────────────────────────────────────────────────────────
	
	%% --- Four 3x3 convergence grid figures ---
	
	\begin{figure*}[t]
		\centering
		\safeincludegraphics[width=0.92\textwidth]{multiobj_steps_vs_relF}
		\caption{Relative function-value gap (relF) vs.\ iteration count across
			all nine test functions ($N{=}10$, $d{=}30$, 20 MC runs).
			Each subplot shows one function; solid curves represent the
			\textsc{DisGrem} family, dashed curves represent baselines.
			Shaded bands indicate $\pm 1$ standard deviation.
			For non-convex functions, $f_{\mathrm{ref}}$ is the best value found by
			multi-start L-BFGS-B (50 restarts, tolerance $10^{-15}$);
			it is not a certified global optimum.}
		\label{fig:relF_steps}
	\end{figure*}
	
	%% --- Performance and data profiles ---
	
	\begin{figure*}[t]
		\centering
		\safeincludegraphics[width=0.92\textwidth]{perf_profiles_panel}
		\caption{Performance profiles (top row) and data profiles (bottom row)
			at three precision levels $\epsilon\in\{10^{-3},10^{-6},10^{-9}\}$.
			Both profile types use iteration count as the budget axis.
			A solver's curve at ordinate~$\pi$ means it solves a
			$\pi$-fraction of the 9 test problems within the given budget.}
		\label{fig:profiles}
	\end{figure*}
	
	\begin{figure*}[t]
		\centering
		\safeincludegraphics[width=0.92\textwidth]{perf_profiles_comm_panel}
		\caption{Communication-budget profiles at three precision levels
			$\epsilon\in\{10^{-3},10^{-6},10^{-9}\}$.
			Top row: performance profiles where the ratio~$\tau$ is computed
			over cumulative communication cost (MB) instead of iterations.
			Bottom row: data profiles with communication cost on the
			$\kappa$-axis.
			Second-order methods incur a higher per-round payload
			($\cO(d^2)$ vs.\ $\cO(d)$ for first-order), yet the
			\textsc{DisGrem} variants remain competitive or superior
			at medium-to-high precision thanks to fewer total iterations.}
		\label{fig:profiles_comm}
	\end{figure*}
	
	Figure~\ref{fig:relF_steps} presents
	convergence trajectories (relF vs.\ iteration) for all nine functions and ten algorithms.
	Figures~\ref{fig:profiles} and~\ref{fig:profiles_comm} summarize
	solver quality via performance and data profiles using iteration count
	and communication cost as the budget metric, respectively.
	Table~\ref{tab:summary} reports key numerical results on four
	representative functions.
	Additional convergence views (combo metric, wall-clock time, communication cost)
	are provided in Appendix~\ref{app:extra_exp}.
	
	\begin{table*}[t]
		\centering
		\caption{Summary on four representative functions
			($N{=}10$, $d{=}30$, 20 MC runs).
			\textbf{Time}: mean wall-clock time (s).
			$K$: iterations used.
			\textbf{min(relF)}: best relative function-value gap achieved.
			\textbf{Comm}: mean cumulative communication (MB).}
		\label{tab:summary}
		\setlength{\tabcolsep}{3pt}
		\footnotesize
		\resizebox{0.95\textwidth}{!}{%
			\begin{tabular}{l rrrr rrrr rrrr rrrr}
				\toprule
				& \multicolumn{4}{c}{\textsc{Huber} ($K_{\max}{=}800$)}
				& \multicolumn{4}{c}{\textsc{LogSumExp} ($K_{\max}{=}400$)}
				& \multicolumn{4}{c}{\textsc{LinLog} ($K_{\max}{=}1500$)}
				& \multicolumn{4}{c}{\textsc{Styblinski--Tang} ($K_{\max}{=}100$)} \\
				\cmidrule(lr){2-5}\cmidrule(lr){6-9}\cmidrule(lr){10-13}\cmidrule(lr){14-17}
				Algorithm & Time & $K$ & min(relF) & MB
				& Time & $K$ & min(relF) & MB
				& Time & $K$ & min(relF) & MB
				& Time & $K$ & min(relF) & MB \\
				\midrule
				\multicolumn{17}{l}{\textit{Proposed (\textsc{DisGrem} family)}} \\
				\textbf{DisGrem}      & 1.4 & 800 & $9\text{e-}6$  & 242
				& 0.8 & 400 & $4\text{e-}10$ & 121
				& 3.1 & 1500 & $2\text{e-}6$  & 444
				& 0.07 & 94  & $1\text{e-}15$ & 30 \\
				\textbf{CeDisGrem}    & 1.3 & 800 & $1\text{e-}6$  & 239
				& 0.8 & 400 & $3\text{e-}11$ & 120
				& 2.9 & 1500 & $4\text{e-}6$  & 449
				& 0.05 & 93  & $5\text{e-}16$ & 28 \\
				\textbf{AdaDisGrem}   & 1.5 & 800 & $3\text{e-}6$  & 242
				& 0.9 & 400 & $3\text{e-}6$  & 121
				& 3.3 & 1500 & $9\text{e-}7$  & 444
				& 0.01 & 17  & $6\text{e-}16$ & 6 \\
				\textbf{CeAdaDisGrem} & 1.4 & 800 & $3\text{e-}6$  & 239
				& 0.8 & 400 & $5\text{e-}7$  & 120
				& 3.2 & 1500 & $2\text{e-}6$  & 446
				& 0.02 & 33  & $2\text{e-}13$ & 10 \\
				\midrule
				\multicolumn{17}{l}{\textit{First-order baselines}} \\
				EXTRA  & 8.2 & 800 & $6\text{e-}3$  & 8
				& 4.0 & 400 & $4\text{e-}1$  & 4
				& 21.7 & 1500 & $1\text{e-}8$  & 15
				& 0.56 & 100 & $8\text{e-}1$  & 1 \\
				DIGing & 7.4 & 800 & $3\text{e-}4$  & 24
				& 3.8 & 400 & $2\text{e-}1$  & 12
				& 10.9 & 989 & $1\text{e-}12$ & 31
				& 0.60 & 100 & $9\text{e-}1$  & 3 \\
				\midrule
				\multicolumn{17}{l}{\textit{Second-order baselines}} \\
				DQM    & 0.9 & 800 & $3\text{e-}1$  & 12
				& 0.5 & 400 & $3\text{e-}3$  & 6
				& 2.1 & 1500 & $2\text{e+}0$  & 23
				& 0.08 & 100 & $2\text{e+}2$  & 2 \\
				ESOM   & 1.7 & 800 & $3\text{e-}1$  & 73
				& 1.0 & 400 & $3\text{e-}3$  & 37
				& 3.5 & 1500 & $1\text{e+}1$  & 139
				& 0.15 & 100 & $3\text{e+}12$ & 9 \\
				SONATA & 1.0 & 800 & $2\text{e-}1$  & 37
				& 0.6 & 400 & $4\text{e-}2$  & 18
				& 2.3 & 1500 & $1\text{e+}0$  & 69
				& 0.03 & 12  & $5\text{e-}16$ & 1 \\
				Net-GIANT & 1.0 & 800 & $5\text{e-}2$ & 37
				& 0.7 & 400 & $2\text{e-}2$  & 18
				& 2.5 & 1500 & $1\text{e+}0$  & 69
				& 0.01 & 12  & $5\text{e-}16$ & 1 \\
				\bottomrule
		\end{tabular}}
	\end{table*}
	
	\noindent\textbf{Key observations.}
	\begin{enumerate}[leftmargin=1.4em,itemsep=2pt]
		\item \textbf{Highest precision on most functions.}
		On LogSumExp, DisGrem and CeDisGrem reach
		relF~$\sim10^{-10}$--$10^{-11}$; every baseline stagnates above
		$10^{-3}$ (first-order) or $10^{-2}$ (SONATA, Net-GIANT).
		On Huber the proposed methods reach $\sim10^{-6}$, well
		below the $>10^{-1}$ plateau of DQM/ESOM/SONATA.
		On the multimodal, non-convex Styblinski--Tang, AdaDisGrem
		hits $\sim10^{-16}$ in just 17 iterations.
		
		\item \textbf{Adaptive $M$ helps on heterogeneous curvature,
			less so on smooth problems.}
		AdaDisGrem gives the best relF on LinLog ($9{\times}10^{-7}$)
		and the fastest run on Styblinski--Tang (17 vs.\ 94 steps for
		fixed-$M$ DisGrem), indicating that the secant estimate
		tracks the local curvature well.
		On LogSumExp, however, AdaDisGrem ($\sim3{\times}10^{-6}$)
		lags behind fixed-$M$ DisGrem ($\sim4{\times}10^{-10}$):
		the estimate $\hat M_{i,k}$ overshoots early on when Hessian
		increments are large and the decay factor $\gamma$ cannot
		bring it down fast enough.
		In short, adaptation removes manual tuning
		(Section~\ref{ssec:ada}) but a well-chosen fixed $M$ can
		still win when the optimal value is known.
		
		\item \textbf{Best first-order baseline: EXTRA, but far less
			precise.}
		On Huber, EXTRA's best relF is $6{\times}10^{-3}$---three
		orders of magnitude above DisGrem's $9{\times}10^{-6}$.
		On LogSumExp it stalls at relF~$\approx0.44$ after 400 iterations.
		A notable exception is DIGing on LinLog: its gradient-tracking
		loop reaches relF~$\approx10^{-12}$, beating the DisGrem
		family ($\sim10^{-6}$).
		LinLog has very flat curvature regions where DisGrem's
		Hessian-based regularization ($M_{\mathrm{fac}}{=}1.0$)
		overestimates curvature; DIGing's simpler dynamics are less
		affected.
		That said, DIGing is $3$--$7\times$ slower in wall-clock time
		and fails on most other functions.
		
		\item \textbf{Second-order baselines: strong on some functions,
			unreliable overall.}
		SONATA and Net-GIANT reach $\sim10^{-16}$ on Styblinski--Tang
		in 12 steps, but stagnate or diverge on Huber and LinLog
		(relF~$>1$).
		DQM diverges on LinLog and Styblinski--Tang; ESOM diverges
		catastrophically on Styblinski--Tang (relF~$\sim10^{12}$).
		No baseline matches the DisGrem family's breadth across all
		nine test cases.
		
		\item \textbf{Communication cost: per-iteration overhead vs.\
			total budget.}
		Each DisGrem iteration exchanges $\cO(d^2)$ Hessian data on
		top of $\cO(d)$ gradient data, so the per-iteration byte count
		exceeds that of first-order methods
		(e.g.\ 121\,MB vs.\ 4\,MB on LogSumExp).
		But EXTRA stalls at relF~$\approx0.44$ whereas DisGrem reaches
		$\sim10^{-10}$---a precision EXTRA cannot attain no matter how
		many more iterations it runs.
		The more informative metric is \emph{communication to a
		target precision}; Section~\ref{ssec:commexp} gives a
		systematic comparison.
	\end{enumerate}
	
	\noindent\textbf{Non-convex problems (empirical).}
	Theorem~\ref{thm:main} assumes convexity; nevertheless, the
	DisGrem family \emph{empirically} performs well on all four
	non-convex benchmarks:
	AdaDisGrem reaches relF~$\sim10^{-7}$ on LinLog (best among
	proposed variants); all four variants achieve relF~$<10^{-12}$
	on Rosenbrock within 300 iterations; AdaDisGrem hits $\sim10^{-16}$
	on Styblinski--Tang in 17 steps; and the family converges to
	$\sim10^{-6}$ on LogReg-NCVR within 1000 steps.
	The regularized Newton structure appears to provide enough
	implicit damping for practical non-convex optimization, though
	a formal convergence analysis for non-convex objectives is
	left for future work.
	
	\noindent\textbf{Empirical evidence of the burn-in/descent phase transition.}
	On ill-conditioned problems (QuadBad, $\kappa{=}10^3$), the relF
	curve exhibits a clear two-phase pattern: an initial plateau of
	$\sim$50--100 iterations during which the tracker dispersions
	contract but the objective barely decreases, followed by a rapid
	descent phase consistent with the $\cO(1/k^2)$ theoretical rate.
	This matches the predicted burn-in/descent structure of
	Theorem~\ref{thm:main}: the burn-in phase corresponds to the
	$K_0(\varepsilon)$ iterations needed for dispersions to decay below
	the inexact Newton threshold, after which the $3/2$-recursion
	drives superlinear convergence.
	
	% ─────────────────────────────────────────────────────────────────────
	\subsection{Communication cost study}\label{ssec:commexp}
	% ─────────────────────────────────────────────────────────────────────
	
	The performance profiles in Section~\ref{ssec:convergence} use
	iteration count as the budget axis; this section complements them
	with a communication-volume perspective.
	We investigate the communication overhead of the DisGrem family
	through two complementary studies on four representative functions
	(Ridge, LogSumExp, Huber, LogReg-real), all at $d{=}30$ with 5 MC runs.
	
	\subsubsection{Benefit of communication-efficient (Ce) variants}
	
	\begin{figure*}[t]
		\centering
		\safeincludegraphics[width=0.92\textwidth]{ce_benefit}
		\caption{Communication cost (MB) required to reach each precision
			level $\epsilon\in\{10^{-3},\ldots,10^{-10}\}$ for the full and Ce
			variants of DisGrem and AdaDisGrem.
			Each subplot corresponds to one function.
			Missing markers indicate the algorithm did not reach the
			given precision within the iteration budget.}
		\label{fig:ce_benefit}
	\end{figure*}
	
	\begin{table*}[t]
		\centering
		\caption{Communication savings (\%) of Ce variants over their full
			counterparts at selected precision levels.
			Positive values indicate savings; negative values indicate that
			Ce requires \emph{more} communication (due to slower convergence
			erasing per-iteration byte savings).}
		\label{tab:ce_savings}
		\setlength{\tabcolsep}{3pt}
		\footnotesize
		\begin{tabular}{llrrrr}
			\toprule
			Function & Pair & $\epsilon{=}10^{-3}$ & $\epsilon{=}10^{-6}$ & $\epsilon{=}10^{-8}$ & $\epsilon{=}10^{-10}$ \\
			\midrule
			\multirow{2}{*}{Ridge}
			& DisGrem $\to$ Ce       & $+5\%$  & $+6\%$   & $-4\%$  & $+1\%$  \\
			& AdaDisGrem $\to$ Ce    & $+8\%$  & $-19\%$  & $-3\%$  & $+0\%$  \\
			\midrule
			\multirow{2}{*}{LogSumExp}
			& DisGrem $\to$ Ce       & $+5\%$  & $+5\%$   & $+5\%$  & $+5\%$  \\
			& AdaDisGrem $\to$ Ce    & $+7\%$  & $-128\%$ & ---     & ---     \\
			\midrule
			\multirow{2}{*}{Huber}
			& DisGrem $\to$ Ce       & $+31\%$ & $-11\%$  & ---     & ---     \\
			& AdaDisGrem $\to$ Ce    & $-75\%$ & $-128\%$ & ---     & ---     \\
			\midrule
			\multirow{2}{*}{LogReg-real}
			& DisGrem $\to$ Ce       & $+6\%$  & $+6\%$   & $-6\%$  & $+1\%$  \\
			& AdaDisGrem $\to$ Ce    & $-0\%$  & $-149\%$ & $-201\%$ & $-141\%$ \\
			\bottomrule
		\end{tabular}
	\end{table*}
	
	Figure~\ref{fig:ce_benefit} and Table~\ref{tab:ce_savings} compare
	the total communication cost needed to reach various precision levels.
	For the fixed-$M$ pair (DisGrem vs.\ CeDisGrem), the Ce variant
	consistently saves 5--6\% at moderate precision ($\epsilon \ge 10^{-6}$)
	on Ridge and LogSumExp.
	On Huber at low precision ($\epsilon{=}10^{-3}$), CeDisGrem saves
	31\% of communication.
	However, at high precision ($\epsilon \le 10^{-8}$), CeDisGrem
	occasionally requires slightly \emph{more} communication because the
	compressed Hessian introduces approximation error that slows late-stage
	convergence, offsetting the per-iteration byte savings.
	
	For the adaptive pair (AdaDisGrem vs.\ CeAdaDisGrem), the picture is
	more nuanced.
	At low precision, CeAdaDisGrem provides modest savings (7--8\% on
	Ridge and LogSumExp), but at high precision the Ce variant's
	convergence deteriorates significantly on Huber and LogReg-real
	(negative savings exceeding 100\%), indicating that Hessian
	compression interferes with the adaptive $M$ estimation.
	This suggests that the Ce mechanism is most beneficial when combined
	with fixed $M$, where the compression noise does not corrupt the
	curvature-tracking feedback loop.

	In summary, the Ce mechanism reliably reduces
	\emph{per-iteration payload} (by compressing the $\cO(d^2)$
	Hessian exchange), but total communication savings depend on
	the convergence-speed trade-off: with fixed $M$, savings of
	5--31\% are typical at moderate precision; with adaptive $M$,
	compression noise can slow convergence enough to \emph{increase}
	total bytes at high precision (Table~\ref{tab:ce_savings}).

	A detailed ablation of the Hessian re-use frequency ($K_{\mathrm{lazy}}$)
	and compression method is provided in Appendix~\ref{app:extra_exp}.
	
	% ─────────────────────────────────────────────────────────────────────
	\subsection{Adaptive mechanism study}\label{ssec:ada}
	% ─────────────────────────────────────────────────────────────────────
	
	We study the adaptive $M$ mechanism of AdaDisGrem on four
	representative functions (Ridge, LogSumExp, LogReg-real, LogReg-NCVR)
	with 5 Monte Carlo runs each.
	
	\begin{figure*}[t]
		\centering
		\safeincludegraphics[width=0.92\textwidth]{ada_vs_fixed_m}
		\caption{Convergence comparison: AdaDisGrem (bold blue) vs.\
			DisGrem at five manually chosen fixed $M$ values
			($0.1\times$, $0.3\times$, $1\times$, $3\times$, $10\times$
			the default $M^*$, grey gradient curves light-to-dark).
			Each subplot corresponds to one function.
			AdaDisGrem matches or outperforms the \emph{best} fixed-$M$
			choice without requiring per-function tuning.}
		\label{fig:ada_vs_fixed}
	\end{figure*}
	
	Figure~\ref{fig:ada_vs_fixed} compares AdaDisGrem against five manually chosen
	fixed-$M$ settings spanning a $100\times$ range.
	AdaDisGrem matches or exceeds the best fixed
	choice on every function while avoiding the divergence that extreme
	values can cause, effectively eliminating the need for per-function
	$M$~calibration.
	Additional results on $M$ trajectories and robustness to initial $M$
	are provided in Appendix~\ref{app:extra_exp}.
	
	% ─────────────────────────────────────────────────────────────────────
	\subsection{Robustness study}\label{ssec:robust}
	% ─────────────────────────────────────────────────────────────────────
	
	We investigate robustness along two complementary axes:
	(i)~sensitivity to the starting point, and
	(ii)~sensitivity to algorithmic parameters ($M_{\mathrm{fac}}$ for
	the DisGrem family, stepsize $\alpha$ for first-order methods, $M$
	for second-order methods).
	All runs use $N{=}10$, $d{=}30$, $N_c{=}3$ (same as the convergence
	benchmarks).
	
	\subsubsection{Starting-point robustness}
	
	Each algorithm is run from 100 independent random initial points
	sampled uniformly on a ball of radius $r$ centered at the default
	initialization.
	We test two regimes: \textit{near} ($r{=}1$) and \textit{far}
	($r{=}3$).
	A run is counted as \emph{successful} if it reaches relF~$<10^{-6}$
	within the iteration budget $K_{\max}$.
	(This threshold is deliberately less stringent than the combo~$<10^{-12}$
	criterion used in the convergence benchmarks, so that iteration counts
	and success rates between the two sections are not directly comparable.)
	Table~\ref{tab:robust} reports
	success rates (\%) across all nine functions and ten algorithms.
	Consequently, iteration counts in Table~\ref{tab:summary}
	and success rates in Table~\ref{tab:robust}
	should not be compared across sections.
	
	\begin{table*}[t]
		\centering
		\caption{Starting-point robustness: success rate (\%) at two initialization radii (100 MC runs each).
			Success criterion: relF~$<10^{-6}$.
			Bold: best in each row; ``---'' denotes 0\%.}
		\label{tab:robust}\label{tab:robust_near}\label{tab:robust_far}
		\setlength{\tabcolsep}{3pt}
		\scriptsize
		\begin{tabular}{l rrrr rr rrrr}
			\toprule
			& \multicolumn{4}{c}{\textit{DisGrem family}}
			& \multicolumn{2}{c}{\textit{1st-order}}
			& \multicolumn{4}{c}{\textit{2nd-order baselines}} \\
			\cmidrule(lr){2-5}\cmidrule(lr){6-7}\cmidrule(lr){8-11}
			Function & DG & CeDG & AdaDG & CeAdaDG
			& EXTRA & DIGing & DQM & ESOM & SON & N-GI \\
			\midrule
			\multicolumn{11}{l}{\textit{Near initialization ($r{=}1$)}} \\
			\midrule
			Ridge       & \textbf{100} & \textbf{100} & \textbf{100} & \textbf{100} & --- & --- & \textbf{100} & 59 & \textbf{100} & \textbf{100} \\
			QuadBad     & 99 & \textbf{100} & 99  & \textbf{100} & --- & --- & 98  & --- & 1   & 1   \\
			LogSumExp   & \textbf{100} & \textbf{100} & 41  & 92  & --- & --- & --- & --- & --- & --- \\
			Huber       & 26 & \textbf{40}  & 26  & 15  & --- & --- & --- & --- & --- & --- \\
			LinLog      & 60 & 12  & 57  & 9   & \textbf{100} & \textbf{100} & --- & --- & --- & --- \\
			LogReg-real & \textbf{100} & 97  & \textbf{100} & 56  & --- & --- & --- & --- & 53  & 49  \\
			Rosenbrock  & \textbf{100} & --- & \textbf{100} & --- & --- & --- & 97  & \textbf{100} & \textbf{100} & \textbf{100} \\
			Styblinski  & --- & --- & \textbf{100} & \textbf{100} & --- & --- & --- & --- & \textbf{100} & \textbf{100} \\
			LogReg-NCVR & \textbf{100} & \textbf{100} & \textbf{100} & 90  & --- & --- & --- & --- & 98  & 96  \\
			\textit{Avg} & \textit{76} & \textit{61} & \textbf{\textit{80}} & \textit{62}
			& \textit{11} & \textit{11} & \textit{33} & \textit{18}
			& \textit{50} & \textit{50} \\
			\midrule
			\multicolumn{11}{l}{\textit{Far initialization ($r{=}3$)}} \\
			\midrule
			Ridge       & \textbf{100} & \textbf{100} & \textbf{100} & \textbf{100} & --- & --- & \textbf{100} & 59  & \textbf{100} & \textbf{100} \\
			QuadBad     & 99  & \textbf{100} & 99  & \textbf{100} & --- & --- & 98  & 1   & 1   & 1   \\
			LogSumExp   & \textbf{100} & \textbf{100} & 42  & 53  & --- & --- & --- & --- & --- & --- \\
			Huber       & 26  & \textbf{37}  & 25  & 12  & --- & --- & --- & --- & --- & --- \\
			LinLog      & 62  & 23  & 55  & 21  & 89  & \textbf{100} & --- & --- & --- & --- \\
			LogReg-real & \textbf{100} & 77  & \textbf{100} & 37  & --- & --- & --- & --- & 51  & 48  \\
			Rosenbrock  & \textbf{100} & --- & \textbf{100} & --- & --- & --- & 82  & 96  & \textbf{100} & \textbf{100} \\
			Styblinski  & --- & --- & \textbf{100} & \textbf{100} & --- & --- & --- & --- & \textbf{100} & \textbf{100} \\
			LogReg-NCVR & \textbf{100} & 91  & \textbf{100} & 65  & --- & 1   & --- & --- & 99  & 97  \\
			\textit{Avg} & \textit{76} & \textit{59} & \textbf{\textit{80}} & \textit{54}
			& \textit{10} & \textit{11} & \textit{31} & \textit{17}
			& \textit{50} & \textit{50} \\
			\bottomrule
		\end{tabular}
	\end{table*}
	
	\noindent\textbf{Key observations.}
	\begin{enumerate}[leftmargin=1.4em,itemsep=2pt]
		\item \textbf{AdaDisGrem is the most robust variant}
		(80\,\% average success, near and far), followed by DisGrem (76\,\%).
		Both maintain nearly identical rates across the two initialization
		radii, indicating stable convergence basins.
		The adaptive mechanism is especially valuable on
		Styblinski--Tang, where fixed-$M$ variants score 0\% but
		AdaDisGrem reaches 100\%.

		\item \textbf{Hessian compression reduces robustness.}
		On Rosenbrock, CeDisGrem drops to 0\% while DisGrem achieves 100\%:
		the Hessian's distributed spectral mass is lost under rank-$r$
		truncation ($r{=}6$).
		Ce variants degrade by 2--8 percentage points on average when
		the starting radius increases; the compression noise accumulated
		in the Hessian tracker degrades both the Newton direction and,
		for CeAdaDisGrem, the $\hat M$ estimation.

		\item \textbf{Baselines are function-specific.}
		First-order methods succeed only on LinLog; DQM only on
		Ridge/QuadBad/Rosenbrock; SONATA/Net-GIANT only on
		Ridge/Rosenbrock/Styblinski--Tang.
		Huber is the hardest function overall (best rate: 40\,\%),
		due to its near-linear tails and flat curvature regions.
	\end{enumerate}
	
	\subsubsection{Parameter sensitivity}
	A comprehensive parameter sweep for all ten algorithms across four functions
	is provided in Appendix~\ref{app:extra_exp}.
	In summary, the DisGrem family is robust across a wide range of $M_{\mathrm{fac}}$,
	while first-order baselines are highly sensitive to stepsize
	and second-order baselines show mixed parameter robustness.
	
	%%%%%%%%%%%%%%%%%%%%%%%%%%%%%%%%%%%%%%%%%%%%%%%%%%%%%%%%%%%%
	
	\subsection{Scalability with problem dimension}\label{ssec:scalability}
	
	We repeat the benchmark at $d\in\{30,100,200\}$ on three functions.
	The DisGrem family reaches relF~$\le 10^{-12}$ within a nearly
	dimension-invariant iteration count, consistent with
	Theorem~\ref{thm:end2end}. Per-iteration wall-clock time scales as
	$\cO(d^3)$ due to the $d\times d$ system solve.
	Full results are in Appendix~\ref{app:extra_exp}.
	
	%%%%%%%%%%%%%%%%%%%%%%%%%%%%%%%%%%%%%%%%%%%%%%%%%%%%%%%%%%%%
	
\section{Conclusion}
	%%%%%%%%%%%%%%%%%%%%%%%%%%%%%%%%%%%%%%%%%%%%%%%%%%%%%%%%%%%%
	
	\noindent\textbf{Summary.}
	Decentralized second-order optimization has long been
	hampered by the tension between exploiting curvature and
	managing network disagreement.
	\textsc{DisGrem} resolves this tension: a reference-step
	framework recasts the average iterate's dynamics as an inexact
	regularized Newton recursion, while an increment-based
	dispersion analysis eliminates the heterogeneity-dependent
	accuracy floors present in prior work.
	The result is a $\cO(\varepsilon^{-1})$ iteration complexity
	that matches the centralized regularized Newton rate---without
	line search, without static Hessian-heterogeneity constants,
	and without assuming iterate boundedness \emph{a priori}.
	Under a logarithmic mixing schedule the total communication
	cost is $\cO(\varepsilon^{-1}\log\varepsilon^{-1})$ rounds;
	with a fixed mixing depth $N_c$ the algorithm converges to an
	$\cO(\rho^{N_c})$-stationarity neighborhood, which in practice
	is much tighter than the worst-case bound suggests.
	
	\noindent\textbf{Experimental findings.}
	Experiments on nine objectives ($d{=}30$, 20 Monte~Carlo runs,
	10 algorithms) support the theory:
	\begin{itemize}[leftmargin=1.2em,itemsep=2pt]
		\item The \textsc{DisGrem} family is the most uniformly reliable
		group among those tested, reaching relF~$\lesssim 10^{-6}$ on seven of
		nine problems. No baseline matches this breadth: each stagnates or
		diverges on at least one test case.
		\item \textsc{AdaDisGrem} attains the highest 100-run robustness
		success rate (80\,\%) and eliminates per-function tuning of~$M$.
		\textsc{CeDisGrem} reduces per-iteration communication by 5--31\,\%
		at moderate precision via Hessian compression, though combining
		compression with adaptive~$M$ can degrade high-precision convergence.
		\item Dimension-scalability experiments ($d\in\{30,100,200\}$)
		confirm that the iteration count to a given accuracy is nearly
		constant in~$d$, consistent with the dimension-free complexity bound.
	\end{itemize}
	
	\noindent\textbf{Limitations and future directions.}\label{sec:conclusion}
	\begin{itemize}[leftmargin=1.2em,itemsep=2pt]
		\item \textbf{Non-convex convergence theory.}
		The theoretical guarantees of Theorems~\ref{thm:main}--\ref{thm:superlinear}
		require convexity, yet the \textsc{DisGrem} family performs
		well empirically on all four non-convex benchmarks.
		The recent centralized breakthrough of
		\citet{ding2025nonconvex}---which extends gradient-regularized
		Newton to non-convex objectives via negative-curvature
		exploitation---provides a natural roadmap:
		integrating their technique with our decentralized tracking
		framework could yield the first decentralized regularized
		Newton method with provable non-convex guarantees,
		a highly compelling direction.
		\item \textbf{Inversion-free and inexact local solves.}
		Each local Newton system currently requires an $\cO(d^3)$ direct
		solve.
		Recent inversion-free approaches---notably
		DINAS~\citep{crane2025dinas} (adaptive inexact Newton via CG)
		and INDO~\citep{yuan2023indo} (inversion-free consensus
		optimization)---demonstrate that iterative solvers can
		replace direct factorization while preserving convergence.
		Embedding such inexact solvers into DisGrem would reduce the
		per-iteration cost to $\cO(d^2)$ or less, substantially
		improving scalability for high-dimensional problems.
		\item \textbf{Network-dependent constants.}
		The increment-based analysis eliminates static heterogeneity
		constants, yet the remaining constants in the complexity
		bounds---burn-in length $K_0(\varepsilon)$, contraction
		coefficients $C_X, C_G, C_H$---still depend implicitly on
		the spectral gap $1{-}\rho$, the Lipschitz constants
		$L_1, L_2$, and the initial tracker dispersions.
		Tightening these dependencies, especially the scaling
		with $\rho$, is an important direction.
		\item \textbf{Stochastic and mini-batch extensions.}
		The current analysis assumes exact local gradients and Hessians.
		Extending to stochastic approximations (e.g.\ via variance reduction) is important
		for large-scale machine learning applications.
		\item \textbf{Time-varying and directed graphs.}
		The mixing matrix $W$ is assumed symmetric and time-invariant.
		Handling directed graphs or dropping links would require a more delicate consensus
		analysis and may affect the $\cO(\varepsilon^{-1}\log(\varepsilon^{-1}))$ bound.
		\item \textbf{Lower bounds.}
		It is an open question whether the $\cO(\varepsilon^{-1})$ iteration
		complexity and the $\cO(\varepsilon^{-1}\log(\varepsilon^{-1}))$
		communication-round count are optimal for decentralized
		second-order methods under convexity, Lipschitz-continuous Hessians,
		and bounded spectral gap.
		\item \textbf{Convergence theory for AdaDisGrem.}
		The adaptive variant lacks formal convergence guarantees;
		establishing that the secant-based estimate $\hat M_{i,k}$ eventually
		satisfies $\hat M_{i,k}\ge L_2$ under suitable regularity conditions
		would close this gap.
	\end{itemize}
	
	%%%%%%%%%%%%%%%%%%%%%%%%%%%%%%%%%%%%%%%%%%%%%%%%%%%%%%%%%%%%

\bibliographystyle{plainnat}
{\small
	\begin{thebibliography}{99}
	\setlength{\itemsep}{0pt}\setlength{\parsep}{0pt}\setlength{\parskip}{0pt}
		
		\bibitem[Cartis et~al.(2011)Cartis, Gould, and Toint]{cartis2011arc}
		C.~Cartis, N.~I.~M. Gould, and Ph.~L. Toint.
		\newblock Adaptive cubic regularisation methods for unconstrained optimization. Part~I: motivation, convergence and numerical results.
		\newblock \emph{Mathematical Programming}, 127:245--295, 2011.
		
		\bibitem[Chang and Lin(2011)]{chang2011libsvm}
		C.-C.~Chang and C.-J.~Lin.
		\newblock {LIBSVM}: a library for support vector machines.
		\newblock \emph{ACM Transactions on Intelligent Systems and Technology}, 2(3):27:1--27:27, 2011.
		\newblock Software available at \url{https://www.csie.ntu.edu.tw/~cjlin/libsvmtools/datasets/}.
		
		\bibitem[Eisen et~al.(2017)Eisen, Mokhtari, and Ribeiro]{eisen2017dqn}
		M.~Eisen, A.~Mokhtari, and A.~Ribeiro.
		\newblock Decentralized quasi-{N}ewton methods.
		\newblock \emph{IEEE Transactions on Signal Processing}, 65(10):2613--2628, 2017.
		
		\bibitem[Jakoveti{\'c} et~al.(2014)Jakoveti{\'c}, Xavier, and Moura]{jakovetic2014fast}
		D.~Jakoveti{\'c}, J.~Xavier, and J.~M.~F. Moura.
		\newblock Fast distributed gradient methods.
		\newblock \emph{IEEE Transactions on Automatic Control}, 59(5):1131--1146, 2014.
		
		\bibitem[Mishchenko(2023)]{mishchenko2023rn}
		K.~Mishchenko.
		\newblock Regularized {N}ewton method with global convergence and complexity bounds.
		\newblock \emph{arXiv preprint arXiv:2112.02089v3}, 2023.
		
		\bibitem[Mokhtari et~al.(2015)Mokhtari, Ling, and Ribeiro]{mokhtari2015networknewton_arxiv}
		A.~Mokhtari, Q.~Ling, and A.~Ribeiro.
		\newblock Network {N}ewton---part {I}: algorithm and convergence.
		\newblock \emph{arXiv preprint arXiv:1504.06017}, 2015.
		
		\bibitem[Mokhtari et~al.(2020)Mokhtari, Koppel, and Ribeiro]{mokhtari2020newtontracking}
		A.~Mokhtari, A.~Koppel, and A.~Ribeiro.
		\newblock Newton tracking: a decentralized method for exact convex optimization.
		\newblock \emph{IEEE Transactions on Signal Processing}, 68:572--587, 2020.
		
		\bibitem[Nedi{\'c} and Ozdaglar(2009)]{nedic2009subgradient}
		A.~Nedi{\'c} and A.~Ozdaglar.
		\newblock Distributed subgradient methods for multi-agent optimization.
		\newblock \emph{IEEE Transactions on Automatic Control}, 54(1):48--61, 2009.
		
		\bibitem[Nedi{\'c} et~al.(2017)Nedi{\'c}, Olshevsky, and Shi]{nedic2017diging}
		A.~Nedi{\'c}, A.~Olshevsky, and W.~Shi.
		\newblock Achieving geometric convergence for distributed optimization over time-varying graphs.
		\newblock \emph{SIAM Journal on Optimization}, 27(4):2597--2633, 2017.
		
		\bibitem[Nesterov and Polyak(2006)]{nesterov2006cubic}
		Y.~Nesterov and B.~T. Polyak.
		\newblock Cubic regularization of {N}ewton method and its global performance.
		\newblock \emph{Mathematical Programming}, 108:177--205, 2006.
		
		\bibitem[Shi et~al.(2015)Shi, Ling, Wu, and Yin]{shi2015extra}
		W.~Shi, Q.~Ling, G.~Wu, and W.~Yin.
		\newblock {EXTRA}: an exact first-order algorithm for decentralized consensus optimization.
		\newblock \emph{SIAM Journal on Optimization}, 25(2):944--966, 2015.
		
		\bibitem[Sun et~al.(2022)Sun, Scutari, and Palomar]{sun2022sonata}
		Y.~Sun, G.~Scutari, and D.~P. Palomar.
		\newblock Distributed nonconvex optimization and learning based on successive convex approximation.
		\newblock \emph{IEEE Transactions on Signal Processing}, 70:5900--5915, 2022.
		
		\bibitem[Ye et~al.(2023)Ye, Luo, and Zhang]{ye2023networkgiant}
		H.~Ye, Z.~Luo, and T.~Zhang.
		\newblock Network-{GIANT}: fully distributed {N}ewton-type optimization via harmonic Hessian consensus.
		\newblock In \emph{Proceedings of the International Conference on Machine Learning (ICML)}, pp.~39947--39975, 2023.
		
		\bibitem[Yuan et~al.(2019)Yuan, Ying, and Sayed]{yuan2019exactdiffusion}
		K.~Yuan, B.~Ying, and A.~H. Sayed.
		\newblock Exact diffusion for distributed optimization and learning---Part {I}: algorithm development.
		\newblock \emph{IEEE Transactions on Signal Processing}, 67(3):708--723, 2019.
		
		% ── New references added during revision ──
		
		\bibitem[Crane et~al.(2025)Crane, Garc\'{\i}a~Trillos, and Kalantzis]{crane2025dinas}
		R.~Crane, N.~Garc\'{\i}a~Trillos, and V.~Kalantzis.
		\newblock {DINAS}: Distributed inexact {N}ewton method with adaptive step sizes.
		\newblock \emph{arXiv preprint arXiv:2402.09920}, 2025.
		
		\bibitem[Ding et~al.(2025)Ding, Gratton, and Gu\'erol]{ding2025nonconvex}
		T.~Ding, S.~Gratton, and M.~Gu\'erol.
		\newblock A gradient-regularized {N}ewton method for non-convex optimization with global convergence and second-order guarantees.
		\newblock \emph{arXiv preprint arXiv:2502.14556}, 2025.
		
		\bibitem[Doikov and Nesterov(2024)]{doikov2024gradient}
		N.~Doikov and Y.~Nesterov.
		\newblock Gradient regularization of {N}ewton method with {B}regman distances.
		\newblock \emph{Mathematical Programming}, 204:1--25, 2024.
		
		\bibitem[Yuan et~al.(2023b)Yuan, Li, and Ling]{yuan2023indo}
		G.~Yuan, X.~Li, and Q.~Ling.
		\newblock {INDO}: {IN}version-free {D}istributed second-{O}rder method for consensus optimization.
		\newblock \emph{arXiv preprint arXiv:2302.14172}, 2023.
		
		\bibitem[Zhang et~al.(2024)Zhang, Chen, and Ye]{zhang2024distributed}
		X.~Zhang, J.~Chen, and H.~Ye.
		\newblock Communication-efficient distributed optimization with compressed {H}essians.
		\newblock \emph{Neural Networks}, 176:106329, 2024.
		
		\bibitem[Alghunaim et~al.(2020)Alghunaim, Ryu, Yuan, and Sayed]{alghunaim2020decentralized}
		S.~A. Alghunaim, E.~K. Ryu, K.~Yuan, and A.~H. Sayed.
		\newblock Decentralized proximal gradient algorithms with linear convergence rates.
		\newblock \emph{IEEE Transactions on Automatic Control}, 66(6):2787--2794, 2020.
		
		\bibitem[Bajovi{\'c} et~al.(2017)Bajovi{\'c}, Jakoveti{\'c}, Kre\-ji{\'c}, and Krklec~Jerinki{\'c}]{bajovic2017newton}
		D.~Bajovi{\'c}, D.~Jakoveti{\'c}, N.~Kre\-ji{\'c}, and N.~Krklec~Jerinki{\'c}.
		\newblock Newton-like method with diagonal correction for distributed optimization.
		\newblock \emph{SIAM Journal on Optimization}, 27(2):1171--1203, 2017.
		
		\bibitem[Beznosikov et~al.(2023)Beznosikov, Samokhin, and Gasnikov]{beznosikov2023biased}
		A.~Beznosikov, V.~Samokhin, and A.~Gasnikov.
		\newblock Distributed methods with compressed communication for solving variational inequalities, with theoretical guarantees.
		\newblock \emph{Computational Mathematics and Mathematical Physics}, 63(5):804--824, 2023.
		
		\bibitem[Doikov and Nesterov(2021)]{doikov2021optimization}
		N.~Doikov and Y.~Nesterov.
		\newblock Optimization methods for large-scale machine learning with second-order guarantees.
		\newblock \emph{Mathematical Programming}, 186:297--327, 2021.
		
		\bibitem[Hanzely et~al.(2022)Hanzely, Doikov, Nesterov, and Richt{\'a}rik]{hanzely2022newton}
		F.~Hanzely, N.~Doikov, Y.~Nesterov, and P.~Richt{\'a}rik.
		\newblock Stochastic Newton and cubic Newton methods with simple local linear-algebra-free implementations.
		\newblock \emph{arXiv preprint arXiv:2012.11809v2}, 2022.
		
		\bibitem[Koloskova et~al.(2019)Koloskova, Stich, and Jaggi]{koloskova2019decentralized}
		A.~Koloskova, S.~U. Stich, and M.~Jaggi.
		\newblock Decentralized stochastic optimization and gossip algorithms with compressed communication.
		\newblock In \emph{Proceedings of the International Conference on Machine Learning (ICML)}, pp.~3478--3487, 2019.
		
		\bibitem[Li et~al.(2021)Li, Lin, and Fang]{li2021decentralized}
		H.~Li, Z.~Lin, and Y.~Fang.
		\newblock Accelerated gradient tracking over time-varying graphs for decentralized optimization.
		\newblock \emph{arXiv preprint arXiv:2104.02596}, 2021.
		
		\bibitem[Mokhtari et~al.(2016)Mokhtari, Shi, Ling, and Ribeiro]{mokhtari2016esom}
		A.~Mokhtari, W.~Shi, Q.~Ling, and A.~Ribeiro.
		\newblock {ESOM}: A second-order method for exact decentralized optimization over networks.
		\newblock \emph{IEEE Transactions on Signal and Information Processing over Networks}, 2(4):507--522, 2016.
		
		\bibitem[Pu and Nedi{\'c}(2021)]{pu2021distributed}
		S.~Pu and A.~Nedi{\'c}.
		\newblock Distributed stochastic gradient tracking methods.
		\newblock \emph{Mathematical Programming}, 187:409--457, 2021.
		
		\bibitem[Xin and Khan(2019)]{xin2019distributed}
		R.~Xin and U.~A. Khan.
		\newblock Distributed heavy-ball: a generalization and acceleration of first-order methods with gradient tracking.
		\newblock \emph{IEEE Transactions on Automatic Control}, 65(6):2627--2633, 2019.
		
		\bibitem[Zhang et~al.(2021)Zhang, Xiao, and Zhang]{zhang2021distributed}
		J.~Zhang, H.~Xiao, and S.~Zhang.
		\newblock Communication-efficient distributed optimization in networks with gradient tracking.
		\newblock In \emph{Advances in Neural Information Processing Systems (NeurIPS)}, 2021.
		
	\end{thebibliography}
}

\appendix
\footnotesize

\section{Additional Experimental Results}\label{app:extra_exp}

This appendix collects supplementary experimental figures referenced in Section~\ref{sec:experiments}.

\subsection{Additional convergence views}

\begin{figure*}[t]
	\centering
	\safeincludegraphics[width=0.92\textwidth]{multiobj_steps_vs_combo}
	\caption{Composite optimality measure (combo $=\|\nabla f(\bar x_k)\|+\mathrm{cons}_k$)
		vs.\ iteration count.
		Same layout and experimental setting as Figure~\ref{fig:relF_steps}.}
	\label{fig:combo_steps}
\end{figure*}

\begin{figure*}[t]
	\centering
	\safeincludegraphics[width=0.92\textwidth]{multiobj_timeCost_vs_relF}
	\caption{relF vs.\ wall-clock time (seconds).
		The \textsc{DisGrem} family achieves high-precision convergence within
		1--3 seconds on most functions, despite the higher per-iteration cost
		of second-order updates.}
	\label{fig:relF_time}
\end{figure*}

\begin{figure*}[t]
	\centering
	\safeincludegraphics[width=0.92\textwidth]{multiobj_commCost_vs_relF}
	\caption{relF vs.\ cumulative communication cost (MB).
		Second-order methods transmit $\cO(d^2)$ bytes per round
		(Hessian tracker), so their curves start further right than
		first-order baselines; nevertheless, the \textsc{DisGrem}
		family reaches high precision ($\mathord{\sim}10^{-6}$--$10^{-15}$)
		within a moderate communication budget on most test
		functions, whereas several baselines stagnate or diverge.}
	\label{fig:relF_comm}
\end{figure*}

\subsection{Lazy Hessian and compression ablation}

\begin{figure*}[t]
	\centering
	\safeincludegraphics[width=0.92\textwidth]{klazy_sweep}
	\caption{Effect of the Hessian re-use frequency ($K_{\mathrm{lazy}}$)
		on the communication--precision trade-off for CeDisGrem.
		$K_{\mathrm{lazy}}{=}1$ updates the Hessian every iteration;
		larger values reuse the previous Hessian for
		$K_{\mathrm{lazy}}{-}1$ consecutive iterations.
		Gradient-coloured curves from light to dark correspond to
		$K_{\mathrm{lazy}}\in\{1,5,10,20,40,80\}$.}
	\label{fig:klazy}
\end{figure*}

\begin{figure*}[t]
	\centering
	\safeincludegraphics[width=0.92\textwidth]{compress_sweep}
	\caption{Effect of compression method and rank on the
		communication--precision trade-off for CeDisGrem.
		Compared configurations: full (no compression),
		Top-$k$ ($k\in\{5\%,10\%,20\%,50\%\}$ of entries),
		Low-Rank ($r\in\{1,2,d/5,d/2\}$).
		Gradient-coloured curves from light to dark represent
		increasing compression aggressiveness.}
	\label{fig:compress}
\end{figure*}

On Ridge and LogReg-real, moderate values of
$K_{\mathrm{lazy}}\in\{5,10\}$ reduce communication by 40--60\% with
negligible precision loss (Figure~\ref{fig:klazy}).
On LogSumExp and Huber, the savings are smaller because the rapidly
changing Hessian landscape penalizes stale curvature information;
$K_{\mathrm{lazy}}\ge 20$ causes noticeable precision degradation.
Low-rank compression with $r=d/5$ achieves a good trade-off between
communication savings and convergence quality across all four functions (Figure~\ref{fig:compress}).

\subsection{Adaptive mechanism: trajectory and robustness}

\begin{figure*}[t]
	\centering
	\safeincludegraphics[width=0.92\textwidth]{ada_m_trajectory}
	\caption{Regularization parameter $M$ trajectory for AdaDisGrem
		(solid) compared to the fixed $M$ used by DisGrem (dashed
		horizontal line).
		Each subplot shows one function; curves are averaged over 5 MC runs.
		The adaptive mechanism adjusts $M$ to the local curvature,
		typically starting high and decaying as the iterate approaches
		the solution.}
	\label{fig:ada_trajectory}
\end{figure*}

\begin{figure*}[t]
	\centering
	\safeincludegraphics[width=0.92\textwidth]{ada_init_m_robust}
	\caption{Robustness of AdaDisGrem to the initial $M$ value.
		Each curve corresponds to a different initial
		$M_0 \in \{0.1M^*,\,0.5M^*,\,M^*,\,3M^*,\,10M^*\}$
		where $M^*$ is the default value from Table~\ref{tab:config}.
		All initializations converge to similar trajectories within
		50--100 iterations, spanning a $100\times$ range.}
	\label{fig:ada_init_m}
\end{figure*}

The adaptive $M$ tracks the local curvature: it starts near
or above the fixed default, then decays as the iterate approaches
stationarity (Figure~\ref{fig:ada_trajectory}).
All initializations in $[0.1M^*,10M^*]$ converge to a common
$M$ trajectory within 50--100 iterations (Figure~\ref{fig:ada_init_m}),
confirming that AdaDisGrem largely eliminates the need for careful $M$~calibration.

\subsection{Parameter sensitivity}

\begin{figure*}[t]
	\centering
	\safeincludegraphics[width=0.92\textwidth]{disgrem_family_combined}
	\caption{Parameter sweep for the \textsc{DisGrem} family
		(DisGrem and CeDisGrem) on four functions.
		Each subplot shows relF vs.\ iteration for 10 values of
		$M_{\mathrm{fac}}$ (gradient-coloured light-to-dark for
		increasing $M_{\mathrm{fac}}$).
		Rows: functions; columns: algorithms.}
	\label{fig:sweep_disgrem}
\end{figure*}

\begin{figure*}[t]
	\centering
	\safeincludegraphics[width=0.92\textwidth]{first-order_combined}
	\caption{Stepsize sensitivity for first-order baselines
		(EXTRA and DIGing) on four functions.
		Each subplot shows relF vs.\ iteration for 10 values of
		$\alpha$ (light-to-dark for increasing $\alpha$).
		Divergent runs are clipped at relF~$=10^4$.}
	\label{fig:sweep_first}
\end{figure*}

\begin{figure*}[t]
	\centering
	\safeincludegraphics[width=0.92\textwidth]{second-order_combined}
	\caption{Key-parameter sensitivity for second-order baselines
		on four functions: ADMM penalty~$c$ for DQM,
		penalty~$\rho$ for ESOM, regularisation~$\tau$ for SONATA,
		and step size~$\alpha$ for Network-GIANT.}
	\label{fig:sweep_second}
\end{figure*}

DisGrem and CeDisGrem converge across the full range of
$M_{\mathrm{fac}}$ on Ridge and the logistic regressions (Figure~\ref{fig:sweep_disgrem}).
On LogSumExp, $M_{\mathrm{fac}}<1$ causes divergence; moderate-to-large
values all converge at similar rates.
EXTRA and DIGing are highly sensitive to $\alpha$ (Figure~\ref{fig:sweep_first}):
even the best $\alpha$ stagnates at relF~$>10^{-1}$ on LogSumExp.
DQM and ESOM diverge for extreme parameter values, while SONATA and
Network-GIANT are more stable but stagnate on LogSumExp regardless of
parameter choice (Figure~\ref{fig:sweep_second}).

\subsection{Dimension scalability}

\begin{figure*}[t]
	\centering
	\safeincludegraphics[width=0.92\textwidth]{scalability_grid}
	\caption{Dimension scalability.
	Each row corresponds to one function (Ridge, LogSumExp, Rosenbrock);
	each column to a dimension ($d=30$, $100$, $200$).
	Curves show median relF over 5~Monte Carlo runs.
	The iteration count to high accuracy is nearly constant across
	dimensions for the \textsc{DisGrem} family, while EXTRA stagnates
	and SONATA converges more slowly.}
	\label{fig:scalability}
\end{figure*}

On Ridge and Rosenbrock, the DisGrem family reaches
relF~$\le 10^{-12}$ within a nearly dimension-invariant iteration
count (Figure~\ref{fig:scalability}), consistent with
Theorem~\ref{thm:end2end}.
On LogSumExp, the final accuracy degrades slightly at $d=200$, yet
DisGrem still outperforms both baselines by several orders of magnitude.
Per-iteration wall-clock time scales as $\cO(d^3)$ (Ridge: 2.3\,s at $d{=}30$ vs.\
118\,s at $d{=}200$), as expected from dense Newton steps.

\section{Dispersion Recursions, Burn-in Analysis, and Logarithmic Mixing}\label{app:dispersion}
% ============================================================

This appendix tightens the ``tracker dispersion'' part of the analysis \emph{step by step} and eliminates the
non-vanishing heterogeneity constants from earlier analyses. The key idea is simple:
we always bound tracker increments using Lipschitz continuity of \emph{differences}
$\nabla F(X_{k+1})-\nabla F(X_k)$ and $\nabla^2F(X_{k+1})-\nabla^2F(X_k)$.

\subsection{Dispersion operators and a contraction identity}

For $Z=[z_1;\dots;z_N]\in\R^{Nd}$ define the (RMS) dispersion
\begin{gather*}
D(Z):=\sqrt{\tfrac1N\textstyle\sum_{i=1}^N\|z_i\!-\!\bar z\|^2}
=\tfrac{1}{\sqrt N}\|(\Pmat\!\otimes\! I_d)Z\|,\\
\bar z:=\tfrac1N\textstyle\sum_i z_i.
\end{gather*}
For stacked matrices $\mathcal H=[H_1;\dots;H_N]\in\R^{Nd^2}$ (with Frobenius norm per block),
define
\[
D_H(\mathcal H):=\sqrt{\frac1N\sum_{i=1}^N\|H_i-\bar H\|_F^2}
=\frac{1}{\sqrt N}\big\|((\Pmat)\otimes I_{d^2})\mathcal H\big\|_F.
\]

\begin{lemma}[Mixing contracts dispersion]\label{lem:mix_contract}
For any stacked vectors $Z\in\R^{Nd}$ and any integer $t\ge1$,
\[
D\big((W^t\otimes I_d)Z\big)\le \rho^t\,D(Z).
\]
Likewise, for any stacked matrices $\mathcal H\in\R^{Nd^2}$,
\[
D_H\big((W^t\otimes I_{d^2})\mathcal H\big)\le \rho^t\,D_H(\mathcal H).
\]
\end{lemma}
\begin{proof}
We prove the vector case; the matrix case is identical with $I_d$ replaced by $I_{d^2}$.
Let $Z^\perp:=((\Pmat)\otimes I_d)Z$. Since $W$ is doubly stochastic, $W^t\J=\J$ and $\J W^t=\J$, hence
$\Pmat W^t=(W^t-\J)$. Therefore
\begin{align*}
&(\Pmat\!\otimes\! I_d)(W^t\!\otimes\! I_d)Z\\
&\quad=((W^t\!-\!\J)\!\otimes\! I_d)Z=((W^t\!-\!\J)\!\otimes\! I_d)Z^\perp,
\end{align*}
because $(W^t-\J)\J=0$ implies $(W^t-\J)=(W^t-\J)\Pmat$.
Taking norms and using Lemma~\ref{lem:consensus} gives
\begin{align*}
\|(\Pmat\!\otimes\! I_d)(W^t\!\otimes\! I_d)Z\|
&=\|((W^t\!-\!\J)\!\otimes\! I_d)Z^\perp\|\\
&\le \|W^t\!-\!\J\|_2\,\|Z^\perp\|
\le \rho^t\|Z^\perp\|.
\end{align*}
Divide by $\sqrt N$ to obtain the claim.
\end{proof}

\subsection{A uniform step bound}

\begin{lemma}[Uniform boundedness of pre-mixed trackers]\label{lem:tilde_track_bounded}
Assume Assumption~\ref{ass:standing} and $\tau_k,t_k\ge 1$ for all $k$.
Then there exists a finite constant $B_G>0$ such that
\[
\sup_{k\ge 0} D(\tilde G_k)\le B_G,
\]
and consequently there exists a finite constant $G_g>0$ such that
\[
\sup_{k\ge 0}\max_{1\le i\le N}\|\tilde g_{i,k}\|\le G_g.
\]
\end{lemma}
\begin{proof}
From \eqref{eq:DG_rec_const} we have
\[
D(\tilde G_{k+1})\le \rho^{\tau_{k+1}}\rho^{t_k}\big(D(\tilde G_k)+2L_1D\big)\le \rho^2\big(D(\tilde G_k)+2L_1D\big),
\]
since $\tau_{k+1},t_k\ge 1$ and $\rho\in[0,1)$.
Set $q:=\rho^2\in[0,1)$ and $b:=2L_1D$. Unrolling yields
\[
D(\tilde G_k)\le q^k D(\tilde G_0)+\frac{qb}{1-q}\quad\forall k\ge 0,
\]
so we may take $B_G:=D(\tilde G_0)+\frac{qb}{1-q}$.
Next, Lemma~\ref{lem:avg-track-g} and the iterate boundedness from Lemma~\ref{lem:lyap_bound} imply that
$\|\tilde{\bar g}_k\|=\|\bar g_k\|=\big\|\frac1N\sum_i \nabla f_i(x_{i,k})\big\|$ is uniformly bounded on the level set;
let $G_{\mathrm{avg}}:=\sup_{k\ge 0}\|\tilde{\bar g}_k\|<\infty$.
For any $i$,
\[
\|\tilde g_{i,k}\|\le \|\tilde{\bar g}_k\|+\|\tilde g_{i,k}-\tilde{\bar g}_k\|
\le G_{\mathrm{avg}}+\sqrt N\,D(\tilde G_k)\le G_{\mathrm{avg}}+\sqrt N\,B_G.
\]
Thus we may take $G_g:=G_{\mathrm{avg}}+\sqrt N\,B_G$.
\end{proof}

For later reference, we record the uniform gradient bound established above:
\begin{equation}\label{eq:step_unif_app}
G_g:=\sup_{k\ge 0}\ \max_{1\le i\le N}\ \|\tilde g_{i,k}\| < \infty,
\end{equation}
where finiteness follows from Lemma~\ref{lem:tilde_track_bounded}.


\begin{lemma}[Uniform step bound]\label{lem:step_bound_app}
For every $k\ge0$ and $i\in\{1,\dots,N\}$,
\[
\|s_{i,k}\|\le \sqrt{\frac{\|\tilde g_{i,k}\|}{M}}.
\]
Consequently, using the boundedness from Lemma~\ref{lem:lyap_bound},
\[
\|x_{i,k+1}-x_{i,k}\|\le 2D
\qquad\text{and}\qquad
\|X_{k+1}-X_k\|\le 2\sqrt N\,D,
\]
where $X_k=[x_{1,k};\dots;x_{N,k}]$.
\end{lemma}
\begin{proof}
Fix $i,k$. The local linear system is
\[
(\mathrm{sym}(\tilde H_{i,k})+(\lambda_{i,k}+\delta_{i,k})I)s_{i,k}=-\tilde g_{i,k}.
\]
By definition of $\delta_{i,k}$, the matrix $\mathrm{sym}(\tilde H_{i,k})+\delta_{i,k}I$ is PSD, hence
\[
\mathrm{sym}(\tilde H_{i,k})+(\lambda_{i,k}+\delta_{i,k})I \succeq \lambda_{i,k}I.
\]
Taking norms, we obtain
\[
\lambda_{i,k}\|s_{i,k}\|\le \|(\mathrm{sym}(\tilde H_{i,k})+(\lambda_{i,k}+\delta_{i,k})I)s_{i,k}\|=\|\tilde g_{i,k}\|.
\]
If $\tilde g_{i,k}=0$, then $s_{i,k}=0$ and the bound holds. Otherwise divide by $\lambda_{i,k}>0$:
\[
\|s_{i,k}\|\le \frac{\|\tilde g_{i,k}\|}{\lambda_{i,k}}
=\frac{\|\tilde g_{i,k}\|}{\sqrt{M\|\tilde g_{i,k}\|}}
=\sqrt{\frac{\|\tilde g_{i,k}\|}{M}}.
\]
For the second part, Lemma~\ref{lem:lyap_bound} gives $\|x_{i,k}\|\le \|x_\star\|+D$ and in particular
$\|x_{i,k+1}-x_{i,k}\|\le \|x_{i,k+1}-x_\star\|+\|x_{i,k}-x_\star\|\le 2D$.
Summing squares yields $\|X_{k+1}-X_k\|\le 2\sqrt N\,D$.
\end{proof}

\begin{lemma}[Post-mixing primal dispersion]\label{lem:DX_post}
For all $k\ge 0$,
\[
D(X_{k+1})\le \rho^{t_k}\Big(D(\tilde X_k)+\frac{\|S_k\|}{\sqrt N}\Big).
\]
\end{lemma}
\begin{proof}
Recall $X_{k+1}=(W^{t_k}\otimes I_d)(\tilde X_k+S_k)$ and that $D(Z)=\|Z^\perp\|/\sqrt N$ with $Z^\perp:=((I-\J)\otimes I_d)Z$.
Since $(W^{t_k}-\J)$ annihilates consensus components, we have
\[
(X_{k+1})^\perp = ((W^{t_k}-\J)\otimes I_d)(\tilde X_k+S_k),
\]
and by Lemma~\ref{lem:consensus} (applied in stacked form) and the triangle inequality,
\[
\| (X_{k+1})^\perp\|
\le \rho^{t_k}\|(\tilde X_k+S_k)^\perp\|
\le \rho^{t_k}\big(\|\tilde X_k^\perp\|+\|S_k^\perp\|\big)
\le \rho^{t_k}\big(\|\tilde X_k^\perp\|+\|S_k\|\big).
\]
Divide by $\sqrt N$ to obtain the claim.
\end{proof}

\begin{lemma}[Post-mixing primal dispersion decay]\label{lem:DX_post_decay}
Under Assumption~\ref{ass:standing} and schedule~\eqref{eq:log_schedule}, there exists a constant $C_X^{+}>0$ such that for all $k\ge 0$,
\[
D(X_k)\le \frac{C_X^{+}}{(k+1)^{p}}.
\]
\end{lemma}
\begin{proof}
From Lemma~\ref{lem:DX_post} and $\rho^{t_k}\le (k+2)^{-p}$ under \eqref{eq:log_schedule}, it remains to bound
$D(\tilde X_k)$ and $\|S_k\|/\sqrt N$ uniformly.
By Lemma~\ref{lem:lyap_bound}, each $\tilde x_{i,k}$ is a convex combination of $\{x_{j,k}\}$, hence $\|\tilde x_{i,k}-x_\star\|\le D$ and
\[
D(\tilde X_k)\le \max_{1\le i\le N}\|\tilde x_{i,k}-\bar x_k\|
\le \max_i \|\tilde x_{i,k}-x_\star\|+\|\bar x_k-x_\star\|\le 2D.
\]
Moreover, Lemma~\ref{lem:step_bound_app} yields $\|S_k\|/\sqrt N\le \sqrt{G_g/M}$, where $G_g$ is defined in \eqref{eq:step_unif_app}.
Therefore,
\[
D(X_{k+1})\le (k+2)^{-p}\Big(2D+\sqrt{\frac{G_g}{M}}\Big),
\]
and the claim follows with $C_X^{+}:=2D+\sqrt{G_g/M}$ after shifting indices.
\end{proof}

\subsection{Primal dispersion recursion (explicit)}

Recall the stacked post-mixing update is
\[
X_{k+1}=(W^{t_k}\otimes I_d)(\tilde X_k+S_k),
\qquad
\tilde X_{k+1}=(W^{\tau_{k+1}}\otimes I_d)X_{k+1}.
\]
\begin{lemma}[Primal dispersion recursion]\label{lem:DX_rec}
For all $k\ge0$,
\[
D(\tilde X_{k+1})
\le \rho^{\tau_{k+1}}\rho^{t_k}\Big(D(\tilde X_k)+\frac{\|S_k\|}{\sqrt N}\Big).
\]
Also,
\[
D(X_{k+1})
\le \rho^{t_k}\Big(D(\tilde X_k)+\frac{\|S_k\|}{\sqrt N}\Big).
\]
\end{lemma}
\begin{proof}
By Lemma~\ref{lem:mix_contract} applied to $Z=\tilde X_k+S_k$,
\[
D(X_{k+1})=D((W^{t_k}\otimes I_d)(\tilde X_k+S_k))\le \rho^{t_k}D(\tilde X_k+S_k).
\]
Next we bound $D(\tilde X_k+S_k)$ by the triangle inequality in the Hilbert space:
\begin{align*}
D(\tilde X_k\!+\!S_k)
&=\tfrac1{\sqrt N}\|(\Pmat\!\otimes\! I_d)(\tilde X_k\!+\!S_k)\|\\
&\le \tfrac1{\sqrt N}\|(\Pmat\!\otimes\! I_d)\tilde X_k\|\\
&\quad+\tfrac1{\sqrt N}\|(\Pmat\!\otimes\! I_d)S_k\|.
\end{align*}
The first term equals $D(\tilde X_k)$.
For the second, $\|((\Pmat)\otimes I_d)S_k\|\le \|S_k\|$, hence it is bounded by $\|S_k\|/\sqrt N$.
This proves the post-mixing inequality.
Finally, apply Lemma~\ref{lem:mix_contract} to the pre-mixing step:
$D(\tilde X_{k+1})\le \rho^{\tau_{k+1}}D(X_{k+1})$ and combine.
\end{proof}

\subsection{Gradient-tracker dispersion recursion (increment-based, fully expanded)}

Write stacked variables
\begin{gather*}
G_k:=[g_{1,k};\dots;g_{N,k}],\;
\tilde G_k:=[\tilde g_{1,k};\dots;\tilde g_{N,k}],\\
\Delta\nabla F_k:=\nabla F(X_{k+1})-\nabla F(X_k).
\end{gather*}
Algorithm~\ref{alg:disgrem} implies the stacked recursion
\begin{gather}
G_{k+1}\!=\!(W^{t_k}\!\otimes\! I_d)(\tilde G_k\!+\!\Delta\nabla F_k),\label{eq:G_update_stack}\\
\tilde G_{k+1}\!=\!(W^{\tau_{k+1}}\!\otimes\! I_d)G_{k+1}.\notag
\end{gather}

\begin{lemma}[Gradient-tracker dispersion recursion (no heterogeneity constants)]\label{lem:DG_rec}
For all $k\ge0$,
\begin{equation}\label{eq:DG_rec}
D(\tilde G_{k+1})
\le \rho^{\tau_{k+1}}\rho^{t_k}\Big(D(\tilde G_k)+\frac{L_1}{\sqrt N}\|X_{k+1}-X_k\|\Big).
\end{equation}
In particular, using Lemma~\ref{lem:step_bound_app},
\begin{equation}\label{eq:DG_rec_const}
D(\tilde G_{k+1})
\le \rho^{\tau_{k+1}}\rho^{t_k}\Big(D(\tilde G_k)+2L_1D\Big).
\end{equation}
\end{lemma}
\begin{proof}
Starting from \eqref{eq:G_update_stack}, substitute $G_{k+1}$ into $\tilde G_{k+1}$:
\begin{align*}
\tilde G_{k+1}
&=(W^{\tau_{k+1}}\!\otimes\! I_d)(W^{t_k}\!\otimes\! I_d)(\tilde G_k\!+\!\Delta\nabla F_k)\\
&=(W^{\tau_{k+1}+t_k}\!\otimes\! I_d)(\tilde G_k\!+\!\Delta\nabla F_k).
\end{align*}
Apply $(\Pmat\otimes I_d)$:
\begin{align*}
&(\Pmat\!\otimes\! I_d)\tilde G_{k+1}\\
&\;=(\Pmat W^{\tau_{k+1}+t_k}\!\otimes\! I_d)(\tilde G_k+\Delta\nabla F_k)\\
&\;=((W^{\tau_{k+1}+t_k}\!-\!\J)\!\otimes\! I_d)(\tilde G_k+\Delta\nabla F_k),
\end{align*}
since $\Pmat W^m=W^m-\J$ for any $m\ge0$.
Taking norms and using $\|W^{m}-\J\|_2\le \rho^{m}$ from Lemma~\ref{lem:consensus},
\[
\|((\Pmat)\otimes I_d)\tilde G_{k+1}\|
\le \rho^{\tau_{k+1}+t_k}\,\|(\tilde G_k+\Delta\nabla F_k)\|.
\]
Divide by $\sqrt N$ and bound the RHS by dispersion plus increment:
\begin{align*}
D(\tilde G_{k+1})
&=\tfrac1{\sqrt N}\|(\Pmat\!\otimes\! I_d)\tilde G_{k+1}\|\\
&\le \rho^{\tau_{k+1}+t_k}\bigl(
\underbrace{\tfrac1{\sqrt N}\|(\Pmat\!\otimes\! I_d)\tilde G_k\|}_{=D(\tilde G_k)}\\
&\qquad+\tfrac1{\sqrt N}\|\Delta\nabla F_k\|\bigr),
\end{align*}
where we used $\|((\Pmat)\otimes I)u\|\le \|u\|$.
It remains to bound $\|\Delta\nabla F_k\|$.
By block structure,
\begin{align*}
\|\Delta\nabla F_k\|^2
&=\textstyle\sum_{i=1}^N\|\nabla f_i(x_{i,k+1})\!-\!\nabla f_i(x_{i,k})\|^2\\
&\le \textstyle\sum_{i=1}^N (L_1\|x_{i,k+1}\!-\!x_{i,k}\|)^2
=L_1^2\|X_{k+1}\!-\!X_k\|^2.
\end{align*}
Taking square roots gives $\|\Delta\nabla F_k\|\le L_1\|X_{k+1}-X_k\|$ and yields \eqref{eq:DG_rec}.
Finally, apply Lemma~\ref{lem:step_bound_app} to get $\|X_{k+1}-X_k\|\le 2\sqrt N D$ and thus \eqref{eq:DG_rec_const}.
\end{proof}

\subsection{Hessian-tracker dispersion recursion (increment-based, fully expanded)}

Define stacked Hessian trackers $\mathcal H_k:=[H_{1,k};\dots;H_{N,k}]\in\R^{Nd^2}$ and
$\tilde{\mathcal H}_k:=[\tilde H_{1,k};\dots;\tilde H_{N,k}]$.
Define the stacked Hessian increment
\[
\Delta\nabla^2 F_k:=\nabla^2 F(X_{k+1})-\nabla^2 F(X_k),
\]
viewed as a block vector in $\R^{Nd^2}$ with Frobenius norm per block.
The Hessian-tracker recursion is
\begin{gather}
\widehat{\mathcal H}_{k+1}\!=\!(W^{t_k}\!\otimes\! I_{d^2})(\tilde{\mathcal H}_k\!+\!\Delta\nabla^2 F_k),\label{eq:H_update_stack}\\
\mathcal H_{k+1}\!=\!\mathrm{sym}(\widehat{\mathcal H}_{k+1}),\quad
\tilde{\mathcal H}_{k+1}\!=\!(W^{\tau_{k+1}}\!\otimes\! I_{d^2})\mathcal H_{k+1}.\notag
\end{gather}

\begin{lemma}[Hessian-tracker dispersion recursion (no heterogeneity constants)]\label{lem:DH_rec}
For all $k\ge0$,
\begin{equation}\label{eq:DH_rec}
D_H(\tilde{\mathcal H}_{k+1})
\le \rho^{\tau_{k+1}}\rho^{t_k}\Big(D_H(\tilde{\mathcal H}_k)+\frac{\sqrt d\,L_2}{\sqrt N}\|X_{k+1}-X_k\|\Big).
\end{equation}
In particular, using Lemma~\ref{lem:step_bound_app},
\begin{equation}\label{eq:DH_rec_const}
D_H(\tilde{\mathcal H}_{k+1})
\le \rho^{\tau_{k+1}}\rho^{t_k}\Big(D_H(\tilde{\mathcal H}_k)+2\sqrt d\,L_2D\Big).
\end{equation}
\end{lemma}
\begin{proof}
Because $\mathrm{sym}(\cdot)$ is linear and non-expansive in Frobenius norm (it is an orthogonal projection),
applying $\mathrm{sym}$ cannot increase dispersion:
\[
D_H(\mathcal H_{k+1})=D_H(\mathrm{sym}(\widehat{\mathcal H}_{k+1}))\le D_H(\widehat{\mathcal H}_{k+1}).
\]
Now, by the same algebra as in Lemma~\ref{lem:DG_rec}, we have
\begin{gather*}
\tilde{\mathcal H}_{k+1}
=(W^{\tau_{k+1}}\!\otimes\! I_{d^2})\mathcal H_{k+1}\;\Rightarrow\\
D_H(\tilde{\mathcal H}_{k+1})\le \rho^{\tau_{k+1}}D_H(\mathcal H_{k+1})
\le \rho^{\tau_{k+1}}D_H(\widehat{\mathcal H}_{k+1}),
\end{gather*}
where we used Lemma~\ref{lem:mix_contract} and the inequality above.
From \eqref{eq:H_update_stack},
\[
\widehat{\mathcal H}_{k+1}\!=\!(W^{t_k}\!\otimes\! I_{d^2})(\tilde{\mathcal H}_k\!+\!\Delta\nabla^2 F_k),
\]
hence Lemma~\ref{lem:mix_contract} yields
\begin{align*}
D_H(\widehat{\mathcal H}_{k+1})
&\le \rho^{t_k}D_H(\tilde{\mathcal H}_k\!+\!\Delta\nabla^2 F_k)\\
&\le \rho^{t_k}\!\Big(D_H(\tilde{\mathcal H}_k)\!+\!\tfrac1{\sqrt N}\|\Delta\nabla^2 F_k\|_F\Big),
\end{align*}
because $\|((\Pmat)\otimes I)\cdot\|_F\le \|\cdot\|_F$.
Combine the last two displays to get \eqref{eq:DH_rec} once we bound $\|\Delta\nabla^2F_k\|_F$.
By block structure and the inequality $\|A\|_F\le \sqrt d\,\|A\|_2$,
\begin{align*}
&\|\Delta\nabla^2 F_k\|_F^2\\
&=\textstyle\sum_{i=1}^N\|\nabla^2 f_i(x_{i,k+1})\!-\!\nabla^2 f_i(x_{i,k})\|_F^2\\
&\le \textstyle\sum_{i=1}^N d\,\|\nabla^2 f_i(x_{i,k+1})\!-\!\nabla^2 f_i(x_{i,k})\|_2^2.
\end{align*}
Using $L_2$-Lipschitzness of $\nabla^2 f_i$,
\[
\|\nabla^2 f_i(x_{i,k+1})\!-\!\nabla^2 f_i(x_{i,k})\|_2\le L_2\|x_{i,k+1}\!-\!x_{i,k}\|,
\]
hence
\begin{align*}
\|\Delta\nabla^2 F_k\|_F^2
&\le \textstyle\sum_{i=1}^N d L_2^2\|x_{i,k+1}\!-\!x_{i,k}\|^2\\
&= dL_2^2\|X_{k+1}\!-\!X_k\|^2.
\end{align*}
Taking square roots gives $\|\Delta\nabla^2 F_k\|_F\le \sqrt d\,L_2\,\|X_{k+1}-X_k\|$,
which yields \eqref{eq:DH_rec}. Finally use Lemma~\ref{lem:step_bound_app} to obtain \eqref{eq:DH_rec_const}.
\end{proof}

\subsection{Polynomial decay under a logarithmic schedule}

Fix $p>2$ and set the explicit schedule
\begin{equation}\label{eq:log_schedule}
\tau_k=t_k=\left\lceil \frac{p\log(k+2)}{-\log\rho}\right\rceil,\qquad k\ge 0,
\end{equation}
so that $\rho^{\tau_k}=\rho^{t_k}\le (k+2)^{-p}$.

\begin{lemma}[A comparison principle]\label{lem:compare}
Let $\{u_k\}_{k\ge0}$ satisfy $u_{k+1}\le a_k(u_k+b)$ with $b\ge0$, where
$a_k\le q<1$ for all $k$ and $a_k\le (k+2)^{-p}$ for all $k$ with $p>2$.
Then there exists $C>0$ such that $u_k\le C/(k+2)^{p-1}$ for all $k\ge0$.
\end{lemma}
\begin{proof}
Since $a_k\le q<1$, we have $u_{k+1}\le q(u_k+b)$ and hence by induction
\[
u_k\le q^k u_0+\sum_{j=0}^{k-1} q^{k-j}qb \le u_0+\frac{qb}{1-q}=:u_{\max}\quad \forall k\ge 0.
\]
Using also $a_k\le (k+2)^{-p}$, it follows that
\[
u_{k+1}\le a_k(u_k+b)\le (k+2)^{-p}(u_{\max}+b)\le (k+2)^{-(p-1)}(u_{\max}+b),
\]
which proves the claim with $C:=u_{\max}+b$.
\end{proof}

\begin{proposition}[Dispersion decay]\label{prop:dispersion_decay}
Under Assumption~\ref{ass:standing} and schedule~\eqref{eq:log_schedule},
there exist constants $C_X,C_G,C_H>0$ such that for all $k\ge0$,
\[
D(\tilde X_k)\le \frac{C_X}{(k+2)^{p-1}},
\qquad
D(\tilde G_k)\le \frac{C_G}{(k+2)^{p-1}},
\qquad
D_H(\tilde{\mathcal H}_k)\le \frac{C_H}{(k+2)^{p-1}}.
\]
\end{proposition}
\begin{proof}
Apply Lemma~\ref{lem:compare} to the recursion in Lemma~\ref{lem:DX_rec} and the bounds \eqref{eq:DG_rec_const} and \eqref{eq:DH_rec_const},
using the bounded forcing constants $b_X:=\sup_k\|S_k\|/\sqrt N<\infty$ (Lemma~\ref{lem:step_bound_app}),
$b_G:=2L_1D$, and $b_H:=2\sqrt d\,L_2D$.
\end{proof}

\subsection{Burn-in conditions for Proposition~\ref{prop:inexact-rn}}

We now show that for \emph{any} target $\varepsilon>0$, there exists a finite
$K_0(\varepsilon)$ such that whenever $\|g_k\|\ge\varepsilon$ and $k\ge K_0(\varepsilon)$,
the three hypotheses of Proposition~\ref{prop:inexact-rn} hold with $\eta=1/12$.
\begin{proposition}[Burn-in conditions for any stationarity level]\label{prop:burnin}
Fix $p>2$ and the schedule~\eqref{eq:log_schedule}.
For every $\varepsilon>0$ there exists an integer $K_0(\varepsilon)\ge0$ such that
for every $k\ge K_0(\varepsilon)$ with $\|g_k\|=\|\nabla f(\bar x_k)\|\ge \varepsilon$,
all three items of Proposition~\ref{prop:inexact-rn} hold with $\eta:=1/12$.

\vspace{1pt}
\noindent\textbf{Moreover (auxiliary conditions).}
With the same choice of $K_0(\varepsilon)$, we may further guarantee that for all $k\ge K_0(\varepsilon)$,
\begin{equation}\label{eq:burnin_aux_enforce}
\Delta_k^g \le \varepsilon
\qquad\text{and}\qquad
\bar\delta_k \le \sqrt{\varepsilon},
\end{equation}
where $\Delta_k^g:=\frac1N\sum_{i=1}^N\|\tilde g_{i,k}-\tilde{\bar g}_k\|$ and $\bar\delta_k:=\frac1N\sum_{i=1}^N\delta_{i,k}$.
\end{proposition}

\begin{proof}
Fix $\varepsilon>0$ and set $\eta:=1/12$.
We will produce an explicit index $K_0(\varepsilon)$ such that for every $k\ge K_0(\varepsilon)$ with $\|g_k\|\ge \varepsilon$,
all three conditions of Proposition~\ref{prop:inexact-rn} hold.

\vspace{1pt}
\noindent\textbf{Step 1: post-mixing disagreement and bridge terms.}
Recall $D_X(k)=\|X_k^\perp\|/\sqrt N = D(X_k)$.
By Lemma~\ref{lem:DX_post_decay}, there exists $C_X^{+}>0$ such that for all $k\ge 0$,
\begin{equation}\label{eq:DX_post_decay_app}
D_X(k)\le \frac{C_X^{+}}{(k+1)^p}.
\end{equation}
Therefore, for all $k\ge 0$,
\begin{equation}\label{eq:bridge_targets_app}
L_1D_X(k)\le \frac{L_1C_X^{+}}{(k+1)^p},
\qquad
L_2D_X(k)\le \frac{L_2C_X^{+}}{(k+1)^p}.
\end{equation}

\vspace{1pt}
\noindent\textbf{Step 2: tracker dispersions.}
From Proposition~\ref{prop:dispersion_decay}, there exist constants $C_G,C_H>0$ such that for all $k\ge 0$,
\begin{equation}\label{eq:GH_decay_app}
D(\tilde G_k)\le \frac{C_G}{(k+2)^{p-1}},
\qquad
D_H(\tilde{\mathcal H}_k)\le \frac{C_H}{(k+2)^{p-1}}.
\end{equation}
By Cauchy--Schwarz and $\|\cdot\|_2\le \|\cdot\|_F$,
\[
\Delta_k^g
=\frac1N\sum_{i=1}^N\|\tilde g_{i,k}-\tilde{\bar g}_k\|
\le \sqrt{\frac1N\sum_{i=1}^N\|\tilde g_{i,k}-\tilde{\bar g}_k\|^2}
= D(\tilde G_k),
\]
and
\begin{align*}
\Delta_k^H
&=\tfrac1N\textstyle\sum_{i}\big\|\mathrm{sym}(\tilde H_{i,k})\!-\!\mathrm{sym}(\tilde{\bar H}_k)\big\|_2\\
&\le \tfrac1N\textstyle\sum_{i}\big\|\mathrm{sym}(\tilde H_{i,k})\!-\!\mathrm{sym}(\tilde{\bar H}_k)\big\|_F\\
&\le D_H(\tilde{\mathcal H}_k).
\end{align*}
Hence, using Proposition~\ref{prop:dispersion_decay},
\begin{equation}\label{eq:Delta_decay_app}
\Delta_k^g\le \frac{C_G}{(k+2)^{p-1}},
\qquad
\Delta_k^H\le \frac{C_H}{(k+2)^{p-1}}.
\end{equation}

\vspace{1pt}
\noindent\textbf{Step 3: a usable lower bound on the reference step.}
Assume $\|g_k\|\ge \varepsilon$ and suppose $k$ is large enough such that $L_1D_X(k)\le \varepsilon/2$.
Then Lemma~\ref{lem:bridge-grad} implies
$\|\tilde{\bar g}_k\|\ge \varepsilon/2$.
Moreover, since $A_k^{\mathrm{ref}}\preceq (M_{H,\max}+\tilde{\bar\lambda}_k)I$,
\[
\|s_k^{\mathrm{ref}}\|
\ge \frac{\|\tilde{\bar g}_k\|}{M_{H,\max}+\tilde{\bar\lambda}_k}
\ge \frac{\varepsilon/2}{M_{H,\max}+\sqrt{MG_g}+G_H}
=:\underline s(\varepsilon),
\]
where we used $\tilde{\bar\lambda}_k=\frac{1}{N}\sum_i(\lambda_{i,k}+\delta_{i,k})\le \sqrt{MG_g}+G_H$, with $G_g$ from Lemma~\ref{lem:step_bound_app} and $G_H$ from Lemma~\ref{lem:tilde_hess_bounded}.

\vspace{1pt}
\noindent\textbf{Step 4: ensure dispersion control (Item~1).}
We first verify the relative dispersion condition \eqref{eq:rel_grad_disp} with $\alpha_d=1/2$.
Since $D(\tilde G_k)=\sqrt{\frac1N\sum_i\|\tilde g_{i,k}-\tilde{\bar g}_k\|^2}$, we have
\[
\max_{1\le i\le N}\|\tilde g_{i,k}-\tilde{\bar g}_k\|
\le \sqrt{N}\,D(\tilde G_k).
\]
From Step 3, for all $k$ large enough with $\|g_k\|\ge \varepsilon$ we have $\|\tilde{\bar g}_k\|\ge \varepsilon/2$.
Thus, choosing $k$ large enough so that $\sqrt N\,D(\tilde G_k)\le \frac12\|\tilde{\bar g}_k\|$,
condition \eqref{eq:rel_grad_disp} holds with $\alpha_d=1/2$.
Using Lemma~\ref{lem:step-disp} and Lemma~\ref{lem:lambda_dispersion} with $\alpha_d:=1/2$,
for all sufficiently large $k$ with $\|g_k\|\ge \varepsilon$ we have
\[
\|\bar s_k-s_k^{\mathrm{ref}}\|
\le \frac{6}{\sqrt{M\varepsilon}}\Delta_k^g+\frac{2}{M}\Delta_k^H+\frac{4}{M}\bar\delta_k.
\]
By Lemma~\ref{lem:delta_control}, $\bar\delta_k\le \Delta_k^H+L_2D_X(k)$.
Combining with \eqref{eq:Delta_decay_app}--\eqref{eq:bridge_targets_app}, there exists $K_1(\varepsilon)$ such that for all $k\ge K_1(\varepsilon)$,
\[
\|\bar s_k-s_k^{\mathrm{ref}}\|\le \frac{\eta}{16}\cdot \frac{\lambda_k}{M_{H,\max}+\lambda_k}\,\|s_k^{\mathrm{ref}}\|.
\]
This implies $\|s_k^{\mathrm{ref}}\|\le 2\|\bar s_k\|$ by the triangle inequality, and hence Item~1 of Proposition~\ref{prop:inexact-rn} holds.

\vspace{1pt}
\noindent\textbf{Step 5: ensure Items~2--3.}
Since $\|s_k^{\mathrm{ref}}\|\le 2\|\bar s_k\|$ and $\lambda_k\ge \sqrt{M\varepsilon}$, it is enough to ensure
\[
L_1D_X(k)\le \frac{\eta}{16}\lambda_k\|s_k^{\mathrm{ref}}\|,
\qquad
L_2D_X(k)\le \frac{\eta}{8}\lambda_k.
\]
By \eqref{eq:bridge_targets_app} and $\|s_k^{\mathrm{ref}}\|\ge \underline s(\varepsilon)$, there exists $K_2(\varepsilon)$ such that both inequalities hold for all $k\ge K_2(\varepsilon)$.

\vspace{1pt}
\noindent\textbf{Step X: ensure $\Delta_k^g\le \varepsilon$.}
By \eqref{eq:Delta_decay_app}, $\Delta_k^g\le \frac{C_G}{(k+2)^{p-1}}$.
Define
\[
K_\Delta(\varepsilon):=\min\Big\{k\ge 0:\ \frac{C_G}{(k+2)^{p-1}}\le \varepsilon\Big\}.
\]
Then for all $k\ge K_\Delta(\varepsilon)$, we have $\Delta_k^g\le \varepsilon$.

\vspace{1pt}
\noindent\textbf{Step Y: ensure $\bar\delta_k\le \sqrt{\varepsilon}$ (convex case).}
By Lemma~\ref{lem:delta_control}, we have for all $k$,
\[
\bar\delta_k\le \Delta_k^H + L_2D_X(k).
\]
Using \eqref{eq:Delta_decay_app} and \eqref{eq:DX_post_decay_app}, we obtain for all $k\ge 0$,
\[
\bar\delta_k
\le \frac{C_H}{(k+2)^{p-1}}+\frac{L_2C_X^{+}}{(k+1)^p}
\le \frac{\widetilde C_\delta}{(k+1)^{p-1}},
\]
where $\widetilde C_\delta:=C_H+L_2C_X^{+}$.
Define
\[
K_\delta(\varepsilon):=\min\Big\{k\ge 1:\ \widetilde C_\delta (k+1)^{-(p-1)}\le \sqrt{\varepsilon}\Big\}.
\]
Then for all $k\ge K_\delta(\varepsilon)$, we have $\bar\delta_k\le \sqrt{\varepsilon}$.

\vspace{1pt}
Finally set
\[
K_0(\varepsilon):=\max\{K_1(\varepsilon),K_2(\varepsilon),K_\Delta(\varepsilon),K_\delta(\varepsilon)\}.
\]
Then for every $k\ge K_0(\varepsilon)$ with $\|g_k\|\ge \varepsilon$, all three items of Proposition~\ref{prop:inexact-rn} hold with $\eta=1/12$,
and \eqref{eq:burnin_aux_enforce} also holds.
\end{proof}

\begin{theorem}[End-to-end complexity under an explicit schedule]\label{thm:end2end}
Assume Assumption~\ref{ass:standing}. Fix $p>2$ and run \textsc{DisGrem} with the explicit schedule~\eqref{eq:log_schedule}.
For $\varepsilon>0$, define
\[
K(\varepsilon):=\min\{k\ge 0:\ \|\nabla f(\bar x_k)\|\le \varepsilon\}.
\]
Then there exists a constant $C_K>0$ such that
\[
K(\varepsilon)\le K_0(\varepsilon)+C_K\,\varepsilon^{-1},
\]
where $K_0(\varepsilon)$ is as in Proposition~\ref{prop:burnin}.
\end{theorem}

\begin{theorem}[Total mixing rounds]\label{thm:comm}
Assume the hypotheses of Theorem~\ref{thm:end2end}. Then
\[
\sum_{k=0}^{K(\varepsilon)-1}(\tau_k+t_k)
=\cO\!\Big((K_0(\varepsilon)+\varepsilon^{-1})\log(K_0(\varepsilon)+\varepsilon^{-1})\Big).
\]
\end{theorem}

\begin{remark}[Explicit order of the burn-in $K_0(\varepsilon)$]\label{rem:burnin_order_app}
The burn-in index $K_0(\varepsilon)$ in Theorem~\ref{thm:end2end} is
determined by requiring the tracker dispersions to decay below
$\cO(\sqrt\varepsilon)$ (see Proposition~\ref{prop:burnin}, Steps~3--4).
Under the logarithmic schedule~\eqref{eq:log_schedule} with
$\tau_k,t_k=\lceil p\log(k+2)/(-\log\rho)\rceil$,
Proposition~\ref{prop:dispersion_decay} yields
$D(\tilde G_k),D_H(\tilde{\mathcal H}_k)=\cO(k^{-(p-1)})$.
It therefore suffices to choose $k$ such that $k^{-(p-1)}\lesssim\sqrt\varepsilon$,
giving
\[
K_0(\varepsilon)=\cO\!\big(\varepsilon^{-1/(2(p-1))}\big).
\]
Since $p>2$ in the standing logarithmic schedule, we have $p-1>1$
and hence $1/(2(p-1))<1/2<1$.
In particular, $K_0(\varepsilon)=o(\varepsilon^{-1})$ and therefore
\emph{does not dominate} the $\cO(\varepsilon^{-1})$ post-burn-in term.
The total iteration complexity in Theorem~\ref{thm:end2end}
(and communication complexity in Theorem~\ref{thm:comm}) has the
post-burn-in phase as its leading term:
\[
K(\varepsilon)=\cO(\varepsilon^{-1}),
\qquad
\sum_{k=0}^{K(\varepsilon)-1}(\tau_k+t_k)
=\cO(\varepsilon^{-1}\log\varepsilon^{-1}).
\]
As a concrete illustration, consider $p=3$ and a target stationarity
$\varepsilon=10^{-8}$.
Then $K_0(\varepsilon)=\cO(\varepsilon^{-1/4})=\cO(10^{2})$,
while the post-burn-in phase requires $\cO(\varepsilon^{-1})=\cO(10^{8})$
iterations.
Hence the burn-in overhead is negligible---roughly $0.001\%$
of the total budget.
Even for very loose networks ($\rho\approx0.97$), larger
problem-dependent constants may increase $K_0$ by a moderate
multiplicative factor, but the $o(\varepsilon^{-1})$ scaling ensures
the burn-in remains a lower-order term.
\end{remark}

% ============================================================

\section{Proofs}\label{app:proofs}
% ============================================================

% ── Lyapunov boundedness (full proof) ───────────────────────────
\begin{proof}[Proof of Lemma~\ref{lem:lyap_bound}]
We establish iterate boundedness via a self-contained \emph{invariant-set
construction} that avoids any circular dependence on the quantity $D$
being proved finite.

\vspace{1pt}
\noindent\textbf{Step 1: parameterized constants on compact sets.}
For any radius $R>0$, the sublevel set
$S(R):=\{x:\,f(x)\le f_\star+R\}$
is \emph{compact} by coercivity (Assumption~\ref{ass:standing}(3)).
Define the \emph{$R$-dependent} Lipschitz and boundedness constants:
\begin{gather*}
  L_1(R):=\sup_{\substack{x,y\in S(R)\\x\ne y}}\frac{\|\nabla f(x)-\nabla f(y)\|}{\|x-y\|},\\
  M_H(R):=\sup_{x\in S(R)}\max_i\|\nabla^2 f_i(x)\|_2,\\
  G_\nabla(R):=\sup_{x\in S(R)}\|\nabla f(x)\|.
\end{gather*}
All three are finite because each $f_i$ is twice continuously differentiable
and $S(R)$ is compact; they are non-decreasing in $R$.
Write $D_S(R):=\sup\{\|x-x_\star\|:\,x\in S(R)\}$.

\vspace{1pt}
\noindent\textbf{Step 2: define the parameterized invariant set.}
Fix $R>0$ (to be determined).
For a primal dispersion budget $B_X>0$ and a gradient-tracker
dispersion budget $B_G>0$, define
\[
  \mathcal S(R,B_X,B_G) \;:=\;
  \bigl\{\,(X_k,\tilde G_k)\;:\;
  f(\bar x_k)\le f_\star+R,\;\;
  D(X_k)\le B_X,\;\;
  D(\tilde G_k)\le B_G
  \,\bigr\}.
\]
We will choose $B_X,B_G$ as the geometric-series fixed points of the
dispersion recursions evaluated at radius $R$, and then pick $R$ large
enough to make $\mathcal S$ forward-invariant.

\vspace{1pt}
\noindent\textbf{Step 3: one-step function-value bound (no descent required).}
Suppose the current state lies in $\mathcal S(R,B_X,B_G)$.
Every local step satisfies $\|s_{i,k}\|\le\sqrt{\|\tilde g_{i,k}\|/M}$
by Lemma~\ref{lem:M-vs-s}; this bound holds \emph{unconditionally}
and does not require the inexact Newton condition.
By $\bar x_{k+1}=\bar x_k+\bar s_k$ (Lemma~\ref{lem:avg-pres}) and
$L_1$-smoothness (Assumption~\ref{ass:standing}(1)),
\[
  f(\bar x_{k+1})-f(\bar x_k)
  \;\le\;
  \langle \nabla f(\bar x_k),\,\bar s_k\rangle
  +\tfrac{L_1(R)}{2}\|\bar s_k\|^2.
\]
From the local solve (Lemma~\ref{lem:M-vs-s}),
$\|s_{i,k}\|\le\sqrt{\|\tilde g_{i,k}\|/M}$ unconditionally.
On the invariant set $\mathcal S$, the average-preservation identity
(Lemma~\ref{lem:avg-track-g}) gives
$\|\tilde{\bar g}_k\|=\|\nabla f(\bar x_k)\|\le G_\nabla(R)$,
and the gradient-tracker dispersion budget yields
$\|\tilde g_{i,k}\|\le \|\tilde{\bar g}_k\|+\sqrt{N}\,D(\tilde G_k)
\le G_\nabla(R)+\sqrt{N}\,B_G=:G_g(R,B_G)<\infty$.
Consequently $\|\bar s_k\|\le\frac{1}{N}\sum_i\|s_{i,k}\|
\le\sqrt{G_g/M}$.
All quantities entering $G_g$ depend only on $R$ (through the
sublevel set $S(R)$) and the dispersion budgets $B_X,B_G$;
in particular, the iterate-diameter constant $D$ being established
does \emph{not} appear, so no circularity arises.
Thus
\begin{equation}\label{eq:rphi_R}
  f(\bar x_{k+1})-f(\bar x_k)
  \;\le\;
  G_\nabla(R)\sqrt{\frac{G_g}{M}}+\frac{L_1(R)}{2}\frac{G_g}{M}
  \;=:\; R_\Phi(R,B_X,B_G) \;<\;\infty.
\end{equation}
We emphasize that this is a \emph{one-step upper bound}, not a descent
guarantee; we do \emph{not} require $R_\Phi\le 0$.

\vspace{1pt}
\noindent\textbf{Step 4: dispersion recursions.}
With $\tau_k,t_k\ge1$ and $q:=\rho^2\in[0,1)$:

\emph{Primal dispersion} (Lemma~\ref{lem:DX_rec} combined with the
pre-mixing contraction $D(\tilde X_k)\le \rho^{\tau_k}D(X_k)$ and the
step bound $\|S_k\|/\sqrt{N}\le \sqrt{G_g/M}$ from
Lemma~\ref{lem:step_bound_app}):
\[
  D(X_{k+1})^2
  \;\le\; q\Bigl(\rho D(X_k)+\sqrt{\tfrac{G_g}{M}}\Bigr)^2
  \;\le\; q^2 D(X_k)^2 + 2q^{3/2}D(X_k)\sqrt{\tfrac{G_g}{M}}+q\tfrac{G_g}{M}.
\]

\emph{Gradient-tracker dispersion} (Lemma~\ref{lem:DG_rec}
and $\|\nabla^2 f_i(x_{i,k})\|_2\le M_H(R)$ on the compact set):
\begin{align*}
  D(\tilde G_{k+1})^2
  &\le q^2 D(\tilde G_k)^2\\
  &\;+ 4q^2 M_H(R)(D_S(R){+}\sqrt{N}B_X)D(\tilde G_k)\\
  &\;+ 4q^2 M_H(R)^2(D_S(R){+}\sqrt{N}B_X)^2.
\end{align*}

\emph{Steady-state dispersion bounds.}
Set $B_X$ and $B_G$ as the unique positive solutions of
$B_X^2 = q^2 B_X^2 + 2q^{3/2}B_X\sqrt{G_g/M}+q\,G_g/M$
and the analogous equation for $B_G$; explicitly,
\begin{equation}\label{eq:disp_fp}
  B_X
  = \frac{\rho\sqrt{G_g/M}}{1-q},
  \qquad
  B_G
  = \frac{2qM_H(R)(D_S(R)+\sqrt{N}B_X)}{1-q}.
\end{equation}
Both are finite and depend only on $R$ and problem parameters.
Provided $D(X_0)\le B_X$ and $D(\tilde G_0)\le B_G$
(achievable by increasing $R$), the dispersion recursions
preserve $D(X_k)\le B_X$ and $D(\tilde G_k)\le B_G$ by induction.

\vspace{1pt}
\noindent\textbf{Step 5: Lyapunov function and forward invariance.}
Define
$\mathcal L_k := (f(\bar x_k)-f_\star) + c_1 D(X_k)^2 + c_2 D(\tilde G_k)^2$
with
\[
  c_1 := \frac{R_\Phi(R)}{(1-q^2)\,B_X^2},
  \qquad
  c_2 := \frac{R_\Phi(R)}{(1-q^2)\,B_G^2}.
\]
Combining the per-step bounds from Steps~3--4 and applying Young's
inequality to the cross terms (details below), one verifies
\begin{equation}\label{eq:lyap_descent_full}
  \mathcal L_{k+1}
  \;\le\;
  \mathcal L_k
  + R_\Phi
  - c_1(1-q^2)D(X_k)^2
  - c_2(1-q^2)D(\tilde G_k)^2
  + c_1\Delta_1 + c_2\Delta_2,
\end{equation}
where $\Delta_1,\Delta_2$ collect the sub-leading mixed terms.
By the choices of $c_1,c_2$ and the Young bounds
\begin{align*}
  2q^{3/2}D(X_k)\sqrt{G_g/M}
  &\le \tfrac{1-q^2}{2}D(X_k)^2 + \tfrac{2q^3 G_g/M}{1-q^2},\\
  4q^2 M_H(D_S{+}\sqrt{N}B_X)D(\tilde G_k)
  &\le \tfrac{1-q^2}{2}D(\tilde G_k)^2
  + \tfrac{8q^4 M_H^2(D_S{+}\sqrt{N}B_X)^2}{1-q^2},
\end{align*}
substituting the Young upper bounds into~\eqref{eq:lyap_descent_full}
and grouping terms yields
\begin{align*}
  \mathcal L_{k+1}
  &\le
  \mathcal L_k + R_\Phi
  - \tfrac{c_1(1{-}q^2)}{2}D(X_k)^2\\
  &\quad- \tfrac{c_2(1{-}q^2)}{2}D(\tilde G_k)^2
  + R_{\mathrm{const}},
\end{align*}
where $R_{\mathrm{const}}:=c_1\!\cdot\!\tfrac{2q^3 G_g/M}{1{-}q^2}
+c_2\!\cdot\!\tfrac{8q^4 M_H^2(D_S{+}\sqrt{N}B_X)^2}{1{-}q^2}$.
By the definitions of $c_1$ and $c_2$,
\[
  \tfrac{c_1(1{-}q^2)}{2}\,B_X^2
  = \tfrac{R_\Phi}{2},
  \qquad
  \tfrac{c_2(1{-}q^2)}{2}\,B_G^2
  = \tfrac{R_\Phi}{2}.
\]
When $D(X_k)\le B_X$ and $D(\tilde G_k)\le B_G$, the negative terms
therefore satisfy
$-\tfrac{c_1(1{-}q^2)}{2}B_X^2-\tfrac{c_2(1{-}q^2)}{2}B_G^2
=-R_\Phi$.
Furthermore, $R_{\mathrm{const}}\le c_1(1{-}q^2)B_X^2/2+c_2(1{-}q^2)B_G^2/2=R_\Phi$
holds by the fixed-point equations~\eqref{eq:disp_fp}
(which encode precisely this balance).
Hence,
\[
  \mathcal L_{k+1}
  \;\le\;
  \mathcal L_k
  + R_\Phi - R_\Phi + R_{\mathrm{const}} - R_{\mathrm{const}}
  \;=\; \mathcal L_k.
\]

\emph{Closing the invariance.}
Choose $R$ large enough so that
$\mathcal L_0\le R$,
$D(X_0)\le B_X(R)$, and
$D(\tilde G_0)\le B_G(R)$.
Such an $R$ exists because $B_X(R),B_G(R)$ grow with $R$
while $\mathcal L_0,D(X_0),D(\tilde G_0)$ are fixed.
Then the state starts in $\mathcal S(R,B_X,B_G)$
and remains there for all $k\ge0$ by the one-step invariance above.
In particular, $\sup_{k\ge0}\mathcal L_k\le \mathcal L_0=:\bar{\mathcal L}<\infty$.

\vspace{1pt}
\noindent\textbf{Step 6: deducing iterate boundedness.}
Since $f(\bar x_k)-f_\star\le\bar{\mathcal L}$, we have
$\bar x_k\in S(\bar{\mathcal L})$,
which is bounded by coercivity (Assumption~\ref{ass:standing}(3)).
Let $R_S:=D_S(\bar{\mathcal L})<\infty$.
Each local iterate satisfies
\[
  \|x_{i,k}-\bar x_k\|\le\sqrt{N}\,D(X_k)\le\sqrt{N}\,B_X,
\]
and hence
\[
  \|x_{i,k}-x_\star\|
  \;\le\;
  \|\bar x_k-x_\star\|+\|x_{i,k}-\bar x_k\|
  \;\le\;
  R_S + \sqrt{N}\,B_X
  \;=:\; D.
\]
Pre-mixed iterates $\tilde x_{i,k}$ are convex combinations of
$\{x_{j,k}\}$ and thus satisfy the same bound.
This establishes~\eqref{eq:bounded_levelset} and completes the proof.
\end{proof}

% ── Tracker spectral-norm boundedness (needed for burn-in / delta) ──
\begin{lemma}[Uniform boundedness of Hessian-tracker dispersion and stabilizer]\label{lem:tilde_hess_bounded}
Assume Assumption~\ref{ass:standing} and that $\tau_k,t_k\ge 1$ for all $k$.
Then there exists a finite constant $B_H>0$ such that
\[
\sup_{k\ge 0} D_H(\tilde{\mathcal H}_k)\le B_H.
\]
Consequently, there exists a finite constant $G_H>0$ such that
\[
\sup_{k\ge 0}\max_{1\le i\le N}\big\|\mathrm{sym}(\tilde H_{i,k})\big\|_2 \le G_H,
\qquad
\sup_{k\ge 0}\bar\delta_k \le G_H.
\]
\end{lemma}

\begin{proof}
From Lemma~\ref{lem:DH_rec} and $\tau_{k+1},t_k\ge 1$, we have
\[
D_H(\tilde{\mathcal H}_{k+1})
\le \rho^{\tau_{k+1}}\rho^{t_k}\Big(D_H(\tilde{\mathcal H}_k)+2\sqrt d\,L_2D\Big)
\le \rho^2\Big(D_H(\tilde{\mathcal H}_k)+2\sqrt d\,L_2D\Big).
\]
Let $q:=\rho^2\in[0,1)$ and $b_H:=2\sqrt d\,L_2D$. Unrolling yields
\[
D_H(\tilde{\mathcal H}_k)\le q^k D_H(\tilde{\mathcal H}_0)+\frac{qb_H}{1-q}\quad \forall k\ge 0,
\]
so we may take $B_H:=D_H(\tilde{\mathcal H}_0)+qb_H/(1-q)$.

Next, by Lemma~\ref{lem:avg-track-h},
$\mathrm{sym}(\tilde{\bar H}_k)=\bar H_k=\frac1N\sum_{i=1}^N\nabla^2 f_i(x_{i,k})$,
hence $\|\mathrm{sym}(\tilde{\bar H}_k)\|_2\le M_{H,\max}$.
For any $i$,
\[
\|\mathrm{sym}(\tilde H_{i,k})\|_2
\le \|\mathrm{sym}(\tilde{\bar H}_k)\|_2\\
\quad + \|\mathrm{sym}(\tilde H_{i,k})\!-\!\mathrm{sym}(\tilde{\bar H}_k)\|_2
\le M_{H,\max}\!+\!\sqrt{N}\,B_H.
\]
Define $G_H:=M_{H,\max} + \sqrt{N}\,B_H$.

Finally, $\delta_{i,k}=\max\{0,-\lambda_{\min}(\mathrm{sym}(\tilde H_{i,k}))\}
\le \|\mathrm{sym}(\tilde H_{i,k})\|_2\le G_H$,
and averaging yields $\bar\delta_k\le G_H$.
\end{proof}

\begin{proof}[Proof of Lemma~\ref{lem:taylor3}]
Fix $x,s$ and define $\phi(t):=h(x+ts)$ for $t\in[0,1]$.
Then $\phi'(t)=\ip{\nabla h(x+ts)}{s}$ and $\phi''(t)=s^\top\nabla^2 h(x+ts)s$.
By $L_2$-Lipschitzness of $\nabla^2 h$,
\begin{align*}
|\phi''(t)\!-\!\phi''(0)|
&=|s^\top(\nabla^2 h(x\!+\!ts)\!-\!\nabla^2 h(x))s|\\
&\le \|\nabla^2 h(x\!+\!ts)\!-\!\nabla^2 h(x)\|_2\|s\|^2\\
&\le L_2 t\|s\|^3.
\end{align*}
Hence $\phi''(t)\le \phi''(0)+L_2 t\|s\|^3$.
Integrating twice,
\begin{align*}
\phi(1)&=\phi(0)+\phi'(0)+\int_0^1\!(1\!-\!t)\phi''(t)\,dt\\
&\le \phi(0)+\phi'(0)+\!\int_0^1\!(1\!-\!t)\big(\phi''(0)+L_2 t\|s\|^3\big)dt.
\end{align*}
Compute $\int_0^1(1-t)\,dt=\tfrac12$ and $\int_0^1(1-t)t\,dt=\tfrac16$ to obtain
\[
\phi(1)\le \phi(0)+\phi'(0)+\tfrac12\phi''(0)+\tfrac{L_2}{6}\|s\|^3.
\]
Substitute $\phi(0)=h(x)$, $\phi'(0)=\ip{\nabla h(x)}{s}$, and $\phi''(0)=s^\top\nabla^2 h(x)s$.
\end{proof}

\begin{proof}[Proof of Lemma~\ref{lem:consensus}]
Because $W$ is symmetric and doubly stochastic, it is diagonalizable with eigenvalues
$1=\lambda_1>\lambda_2\ge\cdots\ge\lambda_N\ge -1$ and eigenvector $\one$ for $\lambda_1$.
The projector $\J=\frac1N\one\one^\top$ is the orthogonal projector onto $\mathrm{span}\{\one\}$ and satisfies $W\J=\J W=\J$.
We show $W^t-\J=(W-\J)^t$ by induction: for $t=1$ trivial. If true for $t$, then
\begin{align*}
(W\!-\!\J)^{t+1}&=(W\!-\!\J)^t(W\!-\!\J)\\
&=(W^t\!-\!\J)(W\!-\!\J)\\
&=W^{t+1}\!-\!W^t\J\!-\!\J W\!+\!\J^2\\
&=W^{t+1}\!-\!\J,
\end{align*}
using $W^t\J=\J$, $\J W=\J$, and $\J^2=\J$.
Thus $\|W^t-\J\|_2=\|(W-\J)^t\|_2\le \|W-\J\|_2^t=\rho^t$.
For stacked vectors, $\|(A\otimes I_d)\|_2=\|A\|_2$.
\end{proof}

\begin{proof}[Proof of Lemma~\ref{lem:gap-grad}]
By $L$-smoothness,
\[
h(y)\le h(x)+\ip{\nabla h(x)}{y-x}+\frac{L}{2}\|y-x\|^2.
\]
Choose $y=x-\frac1L\nabla h(x)$ to obtain
\[
h\Big(x-\frac1L\nabla h(x)\Big)\le h(x)-\frac{1}{2L}\|\nabla h(x)\|^2.
\]
Since $x_\star$ minimizes $h$, $h(x_\star)\le h(x-\frac1L\nabla h(x))$. Rearranging yields the claim.
\end{proof}

\begin{proof}[Proof of Lemma~\ref{lem:avg-pres}]
Because $W^{\tau_k}$ is doubly stochastic,
\[
\tilde{\bar x}_k=\frac1N\sum_{i=1}^N\sum_{j=1}^N [W^{\tau_k}]_{ij}x_{j,k}
=\frac1N\sum_{j=1}^N\Big(\sum_{i=1}^N [W^{\tau_k}]_{ij}\Big)x_{j,k}
=\frac1N\sum_{j=1}^N x_{j,k}=\bar x_k.
\]
Similarly, $W^{t_k}$ is doubly stochastic and $y_{i,k+1}=\tilde x_{i,k}+s_{i,k}$, hence
\begin{align*}
\bar x_{k+1}
&=\tfrac1N\textstyle\sum_{i=1}^N x_{i,k+1}
=\tfrac1N\textstyle\sum_{i,j} [W^{t_k}]_{ij}y_{j,k+1}\\
&=\tfrac1N\textstyle\sum_{j=1}^N y_{j,k+1}
=\tilde{\bar x}_k+\bar s_k=\bar x_k+\bar s_k.
\end{align*}
\end{proof}

\begin{proof}[Proof of Lemma~\ref{lem:avg-track-g}]
Average the tracker update in Algorithm~\ref{alg:disgrem}:
\[
\bar g_{k+1}
=\frac1N\sum_{i=1}^N g_{i,k+1}
=\frac1N\sum_{i=1}^N\sum_{j=1}^N [W^{t_k}]_{ij}\Big(\tilde g_{j,k}+\nabla f_j(x_{j,k+1})-\nabla f_j(x_{j,k})\Big).
\]
Swap sums and use column-stochasticity $\sum_i [W^{t_k}]_{ij}=1$:
\begin{align*}
\bar g_{k+1}&=\tfrac1N\textstyle\sum_{j=1}^N\!\Big(\tilde g_{j,k}+\nabla f_j(x_{j,k+1})-\nabla f_j(x_{j,k})\Big)\\
&=\tilde{\bar g}_k+\tfrac1N\textstyle\sum_{j=1}^N\!\big(\nabla f_j(x_{j,k+1})-\nabla f_j(x_{j,k})\big).
\end{align*}
Pre-mixing preserves averages, hence $\tilde{\bar g}_k=\bar g_k$.
By induction with $\bar g_0=\frac1N\sum_i \nabla f_i(x_{i,0})$, we obtain
$\bar g_k=\frac1N\sum_i \nabla f_i(x_{i,k})$ for all $k$, and thus $\tilde{\bar g}_k=\bar g_k$.
\end{proof}

\begin{proof}[Proof of Lemma~\ref{lem:avg-track-h}]
The proof is identical to Lemma~\ref{lem:avg-track-g} with gradients replaced by Hessians, using linearity of $\mathrm{sym}$.
Averaging the Hessian update gives
\[
\bar H_{k+1}=\tilde{\bar H}_k+\frac1N\sum_{i=1}^N\big(\nabla^2 f_i(x_{i,k+1})-\nabla^2 f_i(x_{i,k})\big),
\]
and $\tilde{\bar H}_k=\bar H_k$. Induction yields $\bar H_k=\frac1N\sum_i \nabla^2 f_i(x_{i,k})$.
\end{proof}

\begin{proof}[Proof of Lemma~\ref{lem:M-vs-s}]
As in Appendix~\ref{app:dispersion}, the local system satisfies
$(\mathrm{sym}(\tilde H_{i,k})+(\lambda_{i,k}+\delta_{i,k})I)\succeq \lambda_{i,k}I$.
Thus $\lambda_{i,k}\|s_{i,k}\|\le \|\tilde g_{i,k}\|$.
Since $\lambda_{i,k}^2=M\|\tilde g_{i,k}\|$, we have $\lambda_{i,k}\|s_{i,k}\|\le \lambda_{i,k}^2/M$,
hence $M\|s_{i,k}\|\le \lambda_{i,k}$. Using $M\ge L_2$ gives $L_2\|s_{i,k}\|\le\lambda_{i,k}$.
\end{proof}

\begin{proof}[Proof of Lemma~\ref{lem:avg-regime}]
From Lemma~\ref{lem:M-vs-s}, $M\|s_{i,k}\|\le \lambda_{i,k}$ for all $i$, hence
\[
\frac1N\sum_{i=1}^N\lambda_{i,k}\ge \frac{M}{N}\sum_{i=1}^N\|s_{i,k}\|\ge M\Big\|\frac1N\sum_{i=1}^N s_{i,k}\Big\|=M\|\bar s_k\|,
\]
using $\|\frac1N\sum_i s_i\|\le \frac1N\sum_i\|s_i\|$.
Since $M\ge L_2$, $L_2\|\bar s_k\|\le \frac1N\sum_i\lambda_{i,k}$.
Finally $\tilde{\bar\lambda}_k=\frac1N\sum_i(\lambda_{i,k}+\delta_{i,k})\ge \frac1N\sum_i\lambda_{i,k}$.
\end{proof}

\begin{proof}[Proof of Lemma~\ref{lem:bridge-grad}]
By Lemma~\ref{lem:avg-track-g}, $\tilde{\bar g}_k=\frac1N\sum_{i=1}^N\nabla f_i(x_{i,k})$.
Therefore
\begin{align*}
\|\nabla f(\bar x_k)\!-\!\tilde{\bar g}_k\|
&=\Big\|\tfrac1N\textstyle\sum_{i=1}^N\!\big(\nabla f_i(\bar x_k)\!-\!\nabla f_i(x_{i,k})\big)\Big\|\\
&\le \tfrac1N\textstyle\sum_{i=1}^N \|\nabla f_i(\bar x_k)\!-\!\nabla f_i(x_{i,k})\|.
\end{align*}
Using $L_1$-Lipschitzness, $\|\nabla f_i(\bar x_k)-\nabla f_i(x_{i,k})\|\le L_1\|x_{i,k}-\bar x_k\|$.
Average and apply Cauchy--Schwarz:
\[
\frac1N\sum_{i=1}^N\|x_{i,k}-\bar x_k\|
\le \sqrt{\frac1N\sum_{i=1}^N\|x_{i,k}-\bar x_k\|^2}=D_X(k).
\]
Combine.
\end{proof}

\begin{proof}[Proof of Lemma~\ref{lem:bridge-hess}]
By Lemma~\ref{lem:avg-track-h}, $\tilde{\bar H}_k=\frac1N\sum_i \nabla^2 f_i(x_{i,k})$.
Thus
\[
\nabla^2 f(\bar x_k)-\mathrm{sym}(\tilde{\bar H}_k)
=\frac1N\sum_{i=1}^N\big(\nabla^2 f_i(\bar x_k)-\nabla^2 f_i(x_{i,k})\big),
\]
and so
\begin{align*}
&\|\nabla^2 f(\bar x_k)\!-\!\mathrm{sym}(\tilde{\bar H}_k)\|_2\\
&\;\le \tfrac1N\textstyle\sum_{i=1}^N\|\nabla^2 f_i(\bar x_k)\!-\!\nabla^2 f_i(x_{i,k})\|_2\\
&\;\le \tfrac{L_2}{N}\textstyle\sum_{i=1}^N\|x_{i,k}\!-\!\bar x_k\|.
\end{align*}
Apply Cauchy--Schwarz to the last average to obtain $L_2 D_X(k)$.
\end{proof}

\begin{proof}[Proof of Lemma~\ref{lem:delta_control}]
Let $A:=\mathrm{sym}(\tilde H_{i,k})$ and $B:=\nabla^2 f(\bar x_k)\succeq0$.
By Weyl's inequality,
\[
\lambda_{\min}(A)\ge \lambda_{\min}(B)-\|A-B\|_2\ge -\|A-B\|_2.
\]
Therefore $-\lambda_{\min}(A)\le \|A-B\|_2$, and since $\delta_{i,k}=\max\{0,-\lambda_{\min}(A)\}$,
we obtain $0\le \delta_{i,k}\le \|A-B\|_2$, proving the first claim.
For the second,
\[
\|A-B\|_2\le \|A-\mathrm{sym}(\tilde{\bar H}_k)\|_2+\|\mathrm{sym}(\tilde{\bar H}_k)-B\|_2.
\]
Average over $i$ to get $\bar\delta_k\le \Delta_k^H+\|\mathrm{sym}(\tilde{\bar H}_k)-\nabla^2 f(\bar x_k)\|_2$,
and apply Lemma~\ref{lem:bridge-hess}.
\end{proof}

\begin{proof}[Proof of Lemma~\ref{lem:step-disp}]
Fix $k$ and abbreviate $A_i:=A_{i,k}$, $A:=A_k^{\mathrm{ref}}$, $\tilde g_i:=\tilde g_{i,k}$, $\tilde{\bar g}:=\tilde{\bar g}_k$.
Then $s_i=-A_i^{-1}\tilde g_i$ and $s^{\mathrm{ref}}=-A^{-1}\tilde{\bar g}$.
Decompose:
\[
s_i-s^{\mathrm{ref}}
=-A_i^{-1}(\tilde g_i-\tilde{\bar g}) + (A^{-1}-A_i^{-1})\tilde{\bar g}.
\]
Since $A_i\succeq \tilde\lambda_{i,k}I\succeq \underline\lambda_k I$, we have $\|A_i^{-1}\|_2\le 1/\underline\lambda_k$.
Thus
\[
\|s_i-s^{\mathrm{ref}}\|
\le \frac1{\underline\lambda_k}\|\tilde g_i-\tilde{\bar g}\|+\|A^{-1}-A_i^{-1}\|_2\,\|\tilde{\bar g}\|.
\]
Use the resolvent identity $A^{-1}-A_i^{-1}=A_i^{-1}(A_i-A)A^{-1}$ to get
\[
\|A^{-1}-A_i^{-1}\|_2\le \|A_i^{-1}\|_2\,\|A_i-A\|_2\,\|A^{-1}\|_2
\le \frac1{\underline\lambda_k}\cdot \frac1{\tilde{\bar\lambda}_k}\,\|A_i-A\|_2,
\]
since $A\succeq \tilde{\bar\lambda}_k I$ implies $\|A^{-1}\|_2\le 1/\tilde{\bar\lambda}_k$.
Now
\[
A_i-A=\big(\mathrm{sym}(\tilde H_{i,k})-\mathrm{sym}(\tilde{\bar H}_k)\big)+(\tilde\lambda_{i,k}-\tilde{\bar\lambda}_k)I,
\]
so $\|A_i-A\|_2\le \|\mathrm{sym}(\tilde H_{i,k})-\mathrm{sym}(\tilde{\bar H}_k)\|_2+|\tilde\lambda_{i,k}-\tilde{\bar\lambda}_k|$.
Combine, then average over $i$ using $\|\bar s-s^{\mathrm{ref}}\|\le \frac1N\sum_i\|s_i-s^{\mathrm{ref}}\|$.
\end{proof}

\begin{proof}[Proof of Lemma~\ref{lem:lambda_dispersion}]
Write $\tilde\lambda_{i,k}-\tilde{\bar\lambda}_k=(\lambda_{i,k}-\bar\lambda_k)+(\delta_{i,k}-\bar\delta_k)$,
where $\bar\lambda_k:=\frac1N\sum_i\lambda_{i,k}$.
Then
\[
\Delta_k^\lambda
=\frac1N\sum_i|\tilde\lambda_{i,k}-\tilde{\bar\lambda}_k|
\le \underbrace{\frac1N\sum_i|\lambda_{i,k}-\bar\lambda_k|}_{=:T_1}
+\underbrace{\frac1N\sum_i|\delta_{i,k}-\bar\delta_k|}_{=:T_2}.
\]
\emph{Step 1: bound $T_2$.}
Since $\delta_{i,k}\ge0$, a standard counting argument yields
\[
\sum_{i=1}^N|\delta_{i,k}-\bar\delta_k|
=2\sum_{i:\delta_{i,k}>\bar\delta_k}(\delta_{i,k}-\bar\delta_k)
\le 2\sum_{i:\delta_{i,k}>\bar\delta_k}\delta_{i,k}
\le 2\sum_{i=1}^N\delta_{i,k}=2N\bar\delta_k,
\]
so $T_2\le 2\bar\delta_k$.

\emph{Step 2: bound $T_1$.}
For any scalars $\{u_i\}$ with $\bar u=\frac1N\sum_i u_i$,
$|u_i-\bar u|=\big|\frac1N\sum_j(u_i-u_j)\big|\le \frac1N\sum_j|u_i-u_j|$.
Averaging over $i$ gives
\[
T_1\le \frac{1}{N^2}\sum_{i=1}^N\sum_{j=1}^N|\lambda_{i,k}-\lambda_{j,k}|.
\]
Let $a_i:=\|\tilde g_{i,k}\|$. Then
\[
|\lambda_{i,k}-\lambda_{j,k}|
=\sqrt{M}\,|\sqrt{a_i}-\sqrt{a_j}|
=\sqrt{M}\,\frac{|a_i-a_j|}{\sqrt{a_i}+\sqrt{a_j}}.
\]
Moreover $|a_i-a_j|\le \|\tilde g_{i,k}-\tilde g_{j,k}\|$.
By triangle inequality,
$\|\tilde g_{i,k}-\tilde g_{j,k}\|\le \|\tilde g_{i,k}-\tilde{\bar g}_k\|+\|\tilde g_{j,k}-\tilde{\bar g}_k\|$,
hence
\[
\frac1{N^2}\sum_{i,j}\|\tilde g_{i,k}-\tilde g_{j,k}\|
\le \frac1{N^2}\sum_{i,j}\big(\|\tilde g_{i,k}-\tilde{\bar g}_k\|+\|\tilde g_{j,k}-\tilde{\bar g}_k\|\big)
=2\Delta_k^g.
\]
Under \eqref{eq:rel_grad_disp}, $a_i\ge \|\tilde{\bar g}_k\|-\|\tilde g_{i,k}-\tilde{\bar g}_k\|\ge (1-\alpha_d)\|\tilde{\bar g}_k\|$ for all $i$,
so $\sqrt{a_i}+\sqrt{a_j}\ge 2\sqrt{(1-\alpha_d)\|\tilde{\bar g}_k\|}$.
Therefore
\[
T_1\le \sqrt{M}\cdot \frac{2\Delta_k^g}{2\sqrt{(1-\alpha_d)\|\tilde{\bar g}_k\|}}
=\frac{\sqrt{M}}{\sqrt{1-\alpha_d}}\cdot \frac{\Delta_k^g}{\sqrt{\|\tilde{\bar g}_k\|}}.
\]
Combine with $T_2\le 2\bar\delta_k$.
\end{proof}

\begin{proof}[Proof of Proposition~\ref{prop:inexact-rn}]
Write
\[
r_k=(\nabla^2 f(\bar x_k)+\lambda_k I)\bar s_k+g_k.
\]
Add and subtract the reference step:
\[
r_k=(\nabla^2 f(\bar x_k)+\lambda_k I)(\bar s_k-s_k^{\mathrm{ref}})
+(\nabla^2 f(\bar x_k)+\lambda_k I)s_k^{\mathrm{ref}}+g_k.
\]
Using $(\mathrm{sym}(\tilde{\bar H}_k)+\lambda_k I)s_k^{\mathrm{ref}}=-\tilde{\bar g}_k$, we have
\[
(\nabla^2 f(\bar x_k)+\lambda_k I)s_k^{\mathrm{ref}}+g_k
=(\nabla^2 f(\bar x_k)-\mathrm{sym}(\tilde{\bar H}_k))s_k^{\mathrm{ref}}+(g_k-\tilde{\bar g}_k).
\]
Therefore
\[
r_k=(\nabla^2 f(\bar x_k)+\lambda_k I)(\bar s_k-s_k^{\mathrm{ref}})
+(\nabla^2 f(\bar x_k)-\mathrm{sym}(\tilde{\bar H}_k))s_k^{\mathrm{ref}}+(g_k-\tilde{\bar g}_k).
\]
Take norms:
\[
\|r_k\|
\le \|\nabla^2 f(\bar x_k)\!+\!\lambda_k I\|_2\|\bar s_k\!-\!s_k^{\mathrm{ref}}\|\\
+\|\nabla^2 f(\bar x_k)\!-\!\mathrm{sym}(\tilde{\bar H}_k)\|_2\|s_k^{\mathrm{ref}}\|
+\|g_k\!-\!\tilde{\bar g}_k\|.
\]
By \eqref{eq:MHmax_def}, $\|\nabla^2 f(\bar x_k)\|_2\le M_{H,\max}$ on the bounded level set, hence
$\|\nabla^2 f(\bar x_k)+\lambda_k I\|_2\le M_{H,\max}+\lambda_k$.
Using item (1),
\begin{align*}
&\|\nabla^2 f(\bar x_k)\!+\!\lambda_k I\|_2\,\|\bar s_k\!-\!s_k^{\mathrm{ref}}\|\\
&\;\le (M_{H,\max}\!+\!\lambda_k)\!\cdot\!\tfrac{\eta}{8}\!\cdot\!\tfrac{\lambda_k}{M_{H,\max}+\lambda_k}\|\bar s_k\|
=\tfrac{\eta}{8}\lambda_k\|\bar s_k\|.
\end{align*}
Next, item (3) and Lemma~\ref{lem:bridge-hess} imply
$\|\nabla^2 f(\bar x_k)-\mathrm{sym}(\tilde{\bar H}_k)\|_2\le \frac{\eta}{8}\lambda_k$.
Moreover,
\[
\|s_k^{\mathrm{ref}}\|\le \|\bar s_k\|+\|\bar s_k-s_k^{\mathrm{ref}}\|
\le \Big(1+\frac{\eta}{8}\Big)\|\bar s_k\|\le \frac98\|\bar s_k\|,
\]
since $\eta\le 1/12$ implies $1+\eta/8\le 9/8$.
Therefore the second term is bounded by
\[
\|\nabla^2 f(\bar x_k)-\mathrm{sym}(\tilde{\bar H}_k)\|_2\,\|s_k^{\mathrm{ref}}\|
\le \frac{\eta}{8}\lambda_k\cdot \frac98\|\bar s_k\|=\frac{9\eta}{64}\lambda_k\|\bar s_k\|.
\]
Finally, item (2) and Lemma~\ref{lem:bridge-grad} imply $\|g_k-\tilde{\bar g}_k\|\le \frac{\eta}{8}\lambda_k\|\bar s_k\|$.
Summing yields
\[
\|r_k\|\le \left(\frac{\eta}{8}+\frac{9\eta}{64}+\frac{\eta}{8}\right)\lambda_k\|\bar s_k\|
=\frac{25\eta}{64}\lambda_k\|\bar s_k\|\le \eta\lambda_k\|\bar s_k\|.
\]
The inequality $L_2\|\bar s_k\|\le \lambda_k$ is Lemma~\ref{lem:avg-regime}.
\end{proof}

\begin{proof}[Proof of Lemma~\ref{lem:descent}]
Apply Lemma~\ref{lem:taylor3} to $f$ at $(\bar x_k,\bar s_k)$:
\[
f(\bar x_k+\bar s_k)\le f(\bar x_k)+\ip{g_k}{\bar s_k}+\frac12\bar s_k^\top\nabla^2 f(\bar x_k)\bar s_k+\frac{L_2}{6}\|\bar s_k\|^3.
\]
Let $d_k:=(\nabla^2 f(\bar x_k)+\lambda_k I)\bar s_k+g_k$.
Then $\ip{g_k}{\bar s_k}=-\bar s_k^\top\nabla^2 f(\bar x_k)\bar s_k-\lambda_k\|\bar s_k\|^2+\ip{d_k}{\bar s_k}$.
Substitute:
\[
f(\bar x_{k+1})\le f(\bar x_k)-\frac12\bar s_k^\top\nabla^2 f(\bar x_k)\bar s_k-\lambda_k\|\bar s_k\|^2+\ip{d_k}{\bar s_k}+\frac{L_2}{6}\|\bar s_k\|^3.
\]
Since $\nabla^2 f(\bar x_k)\succeq 0$, drop the nonpositive term. Also
$\ip{d_k}{\bar s_k}\le \|d_k\|\,\|\bar s_k\|\le \eta\lambda_k\|\bar s_k\|^2$
and $L_2\|\bar s_k\|\le \lambda_k$ implies $\frac{L_2}{6}\|\bar s_k\|^3\le \frac16\lambda_k\|\bar s_k\|^2$.
Combine.
\end{proof}

\begin{proof}[Proof of Lemma~\ref{lem:grad_bound}]
By the integral form of Taylor's theorem,
\[
\nabla f(\bar x_k+\bar s_k)=g_k+\nabla^2 f(\bar x_k)\bar s_k+r_k,
\qquad \|r_k\|\le \frac{L_2}{2}\|\bar s_k\|^2.
\]
Using $d_k=(\nabla^2 f(\bar x_k)+\lambda_k I)\bar s_k+g_k$, we have
$g_k+\nabla^2 f(\bar x_k)\bar s_k=-\lambda_k\bar s_k+d_k$.
Thus
\[
\|g_{k+1}\|\le \lambda_k\|\bar s_k\|+\|d_k\|+\|r_k\|
\le \lambda_k\|\bar s_k\|+\eta\lambda_k\|\bar s_k\|+\frac{L_2}{2}\|\bar s_k\|^2.
\]
Finally $L_2\|\bar s_k\|\le \lambda_k$ implies $\frac{L_2}{2}\|\bar s_k\|^2\le \frac12\lambda_k\|\bar s_k\|$.
\end{proof}

\begin{lemma}[Regularizer scaling after burn-in]\label{lem:lambda_scaling}
Assume Assumption~\ref{ass:standing} and the schedule~\eqref{eq:log_schedule}.
Fix any $\varepsilon>0$ and let $K_0(\varepsilon)$ be as in Proposition~\ref{prop:burnin}.
Then there exists a constant $C_\lambda>0$ such that for all $k\ge K_0(\varepsilon)$ with $\|g_k\|\ge \varepsilon$,
\[
\lambda_k\le C_\lambda\,\sqrt{\|g_k\|}.
\]
\end{lemma}
\begin{proof}
Recall $\lambda_k=\tilde{\bar\lambda}_k=\frac1N\sum_{i=1}^N(\lambda_{i,k}+\delta_{i,k})$ with
$\lambda_{i,k}=\sqrt{M\|\tilde g_{i,k}\|}$ and $\delta_{i,k}\ge 0$.
By Jensen's inequality (concavity of $x\mapsto\sqrt{x}$),
\[
\frac1N\sum_{i=1}^N \sqrt{\|\tilde g_{i,k}\|}
\le \sqrt{\frac1N\sum_{i=1}^N \|\tilde g_{i,k}\|}.
\]
Moreover, by the triangle inequality and the definition of $\Delta_k^g$,
\[
\frac1N\sum_{i=1}^N \|\tilde g_{i,k}\|
\le \|\tilde{\bar g}_k\|+\Delta_k^g.
\]
By Lemma~\ref{lem:bridge-grad}, $\|\tilde{\bar g}_k-g_k\|\le L_1D_X(k)$.
Using Item~2 of Proposition~\ref{prop:inexact-rn} (which holds for all $k\ge K_0(\varepsilon)$ with $\|g_k\|\ge \varepsilon$),
together with $\|(H_k+\lambda_k I)\bar s_k\|=\|g_k-r_k\|\ge \lambda_k\|\bar s_k\|$ and $\|r_k\|\le \eta\lambda_k\|\bar s_k\|$,
we obtain $(1-\eta)\lambda_k\|\bar s_k\|\le \|g_k\|$ and hence
\[
L_1D_X(k)\le \frac{\eta}{8}\lambda_k\|\bar s_k\|\le \frac{\eta}{8(1-\eta)}\|g_k\|.
\]
Therefore,
\[
\|\tilde{\bar g}_k\|\le \|g_k\|+L_1D_X(k)\le \Big(1+\frac{\eta}{8(1-\eta)}\Big)\|g_k\|.
\]
In addition, Proposition~\ref{prop:burnin} ensures $\Delta_k^g\le \varepsilon\le \|g_k\|$ for all $k\ge K_0(\varepsilon)$ with $\|g_k\|\ge \varepsilon$.
Thus,
\[
\frac1N\sum_{i=1}^N \|\tilde g_{i,k}\|
\le \Big(2+\frac{\eta}{8(1-\eta)}\Big)\|g_k\|.
\]
Hence,
\[
\lambda_k
= \frac1N\sum_{i=1}^N\big(\sqrt{M\|\tilde g_{i,k}\|}+\delta_{i,k}\big)
\le \sqrt{M}\sqrt{\frac1N\sum_{i=1}^N \|\tilde g_{i,k}\|}+\bar\delta_k
\le C_0\sqrt{\|g_k\|}+\bar\delta_k,
\]
with $C_0:=\sqrt{M\big(2+\frac{\eta}{8(1-\eta)}\big)}$.
Finally, Proposition~\ref{prop:burnin} also ensures $\bar\delta_k\le \sqrt{\varepsilon}\le \sqrt{\|g_k\|}$ on the same index set.
Absorbing this into the constant yields $\lambda_k\le (C_0+1)\sqrt{\|g_k\|}$, completing the proof.
\end{proof}

\begin{proof}[Proof of Lemma~\ref{lem:steady}]
Assume $\eta\le 1/12$, so $c_d:=1-\eta-\frac16\ge \frac34$ in Lemma~\ref{lem:descent}.
Also from Lemma~\ref{lem:grad_bound}, $C_g:=1+\eta+\frac12\le 2$.
For $k\in\mathcal I$, $\|g_{k+1}\|\ge \frac14\|g_k\|$ and $\|g_{k+1}\|\le C_g\lambda_k\|\bar s_k\|$ imply
$\|\bar s_k\|\ge \|g_k\|/(4C_g\lambda_k)$.
Then Lemma~\ref{lem:descent} yields
\[
\Phi_k-\Phi_{k+1}\ge c_d\lambda_k\|\bar s_k\|^2
\ge c_d\lambda_k\left(\frac{\|g_k\|}{4C_g\lambda_k}\right)^2
=\frac{c_d}{16C_g^2}\frac{\|g_k\|^2}{\lambda_k}.
\]
By the additional assumption in Lemma~\ref{lem:steady}, we have $\lambda_k\le C_\lambda\sqrt{\|g_k\|}$.
Therefore,
\[
\frac{\|g_k\|^2}{\lambda_k}\ge \frac{1}{C_\lambda}\|g_k\|^{3/2}.
\]
Finally, convexity and the bounded level set imply $\Phi_k\le \ip{g_k}{\bar x_k-x_\star}\le D\|g_k\|$, hence
$\|g_k\|^{3/2}\ge D^{-3/2}\Phi_k^{3/2}$.
Collect constants into $\tau$.
\end{proof}

\begin{proof}[Proof of Lemma~\ref{lem:solve32}]
If $\Phi_k=0$ for some $k$, then $\Phi_{k'}=0$ for all $k'\ge k$ and the claim is trivial.
Assume $\Phi_k>0$ for all $k$.

From $\Phi_{k+1}\le \Phi_k-\tau \Phi_k^{3/2}=\Phi_k(1-\tau\sqrt{\Phi_k})$, we consider two cases.

\emph{Case 1: $\tau\sqrt{\Phi_k}\ge 1$.}
Then $\Phi_{k+1}\le 0$. Since $\Phi_{k+1}\ge 0$, we must have $\Phi_{k+1}=0$,
contradicting $\Phi_{k+1}>0$.
Hence this case cannot occur under the standing assumption $\Phi_k>0$.

Therefore, necessarily $\tau\sqrt{\Phi_k}\in(0,1)$ for all $k$.

Let $\psi_k:=\Phi_k^{-1/2}$. Then
\[
\psi_{k+1}=\Phi_{k+1}^{-1/2}\ge \Phi_k^{-1/2}(1\!-\!\tau\sqrt{\Phi_k})^{-1/2}\\
=\psi_k(1\!-\!u_k)^{-1/2},\; u_k:=\tau\sqrt{\Phi_k}\in(0,1).
\]
Using convexity of $(1-u)^{-1/2}$ on $[0,1)$ and $(1-u)^{-1/2}\ge 1+\frac12u$, we obtain
\[
\psi_{k+1}\ge \psi_k\Big(1+\frac12 u_k\Big)=\psi_k+\frac{\tau}{2}.
\]
Thus $\psi_k\ge \psi_0+\frac{\tau}{2}k\ge \frac{\tau}{2}(k+2)$ after adjusting constants, and hence
\[
\Phi_k=\psi_k^{-2}\le \frac{4}{\tau^2(k+2)^2}.
\]
\end{proof}

\begin{proof}[Proof of Theorem~\ref{thm:main}]
Fix $\eta\le 1/12$ and assume Proposition~\ref{prop:inexact-rn} holds for all $k\ge 0$.
Define $\Phi_k:=f(\bar x_k)-f_\star$ and $g_k:=\nabla f(\bar x_k)$.

\vspace{1pt}
\noindent\textbf{Step 0: $\Phi_k$ is nonincreasing.}
By Proposition~\ref{prop:inexact-rn}, the hypotheses of Lemma~\ref{lem:descent} hold at every iteration.
Hence, with $c_d:=1-\eta-\frac16>0$,
\begin{equation}\label{eq:Phi_monotone}
\Phi_{k+1}\le \Phi_k - c_d\,\lambda_k\|\bar s_k\|^2\le \Phi_k,
\qquad \forall k\ge 0.
\end{equation}
In particular, $\{\Phi_k\}$ is nonincreasing.

\vspace{1pt}
\noindent\textbf{Step 1: a uniform gradient growth bound.}
Let $r_k:=(\nabla^2 f(\bar x_k)+\lambda_k I)\bar s_k+g_k$ as defined in the paper.
From Proposition~\ref{prop:inexact-rn}, $\|r_k\|\le \eta\lambda_k\|\bar s_k\|$.
Moreover,
\[
\lambda_k\|\bar s_k\|
\le \|(\nabla^2 f(\bar x_k)+\lambda_k I)\bar s_k\|
=\|g_k-r_k\|
\le \|g_k\|+\|r_k\|
\le \|g_k\|+\eta\lambda_k\|\bar s_k\|.
\]
Rearranging yields
\begin{equation}\label{eq:lambda_s_upper}
\lambda_k\|\bar s_k\|\le \frac{1}{1-\eta}\|g_k\|,\qquad \forall k\ge 0.
\end{equation}
By Lemma~\ref{lem:grad_bound} (which applies at every iteration under Proposition~\ref{prop:inexact-rn}),
\[
\|g_{k+1}\|\le C_g\,\lambda_k\|\bar s_k\|,
\qquad
C_g:=1+\eta+\frac12.
\]
Combining with \eqref{eq:lambda_s_upper} gives the uniform growth bound
\begin{equation}\label{eq:grad_growth}
\|g_{k+1}\|\le \kappa_{\mathrm{inc}}\|g_k\|,
\qquad
\kappa_{\mathrm{inc}}:=\frac{C_g}{1-\eta}.
\end{equation}
Since $\eta\le 1/12$, we have $C_g\le 2$ and thus $\kappa_{\mathrm{inc}}\le 24/11<4$.

\vspace{1pt}
\noindent\textbf{Step 2: decay along the steady-step subsequence.}
Recall the index sets
\[
\mathcal I:=\{k\ge 0:\ \|g_{k+1}\|\ge \tfrac14\|g_k\|\},
\qquad
\mathcal S:=\{k\ge 0:\ \|g_{k+1}\|< \tfrac14\|g_k\|\}.
\]
Let $n_k:=|\mathcal I\cap\{0,1,\dots,k-1\}|$ be the number of steady indices among the first $k$ iterations.
List these indices increasingly as $0\le i_0<i_1<\cdots<i_{n_k-1}\le k-1$ (if $n_k=0$ skip this step).
Define the subsequence $u_j:=\Phi_{i_j}$.
By Lemma~\ref{lem:steady}, for each steady index $i_j\in\mathcal I$,
\[
\Phi_{i_j}-\Phi_{i_j+1}\ge \tau\,\Phi_{i_j}^{3/2}=\tau\,u_j^{3/2}.
\]
Using monotonicity \eqref{eq:Phi_monotone}, we have $\Phi_{i_{j+1}}\le \Phi_{i_j+1}$.
Therefore,
\[
u_{j+1}=\Phi_{i_{j+1}}
\le \Phi_{i_j+1}
\le \Phi_{i_j}-\tau\,\Phi_{i_j}^{3/2}
= u_j-\tau\,u_j^{3/2}.
\]
Applying Lemma~\ref{lem:solve32} to $\{u_j\}$ yields
\begin{equation}\label{eq:u_bound}
u_j\le \frac{4}{\tau^2(j+2)^2},\qquad \forall j\ge 0.
\end{equation}
Since $\Phi_k\le \Phi_{i_{n_k-1}}=u_{n_k-1}$ by monotonicity, we obtain
\begin{equation}\label{eq:Phi_by_nk}
\Phi_k\le \frac{4}{\tau^2(n_k+1)^2}\le \frac{4}{\tau^2(n_k+2)^2},
\qquad \forall k\ge 1.
\end{equation}

\vspace{1pt}
\noindent\textbf{Step 3: two-case bound yielding $\cO(\varepsilon^{-1})$ iteration complexity.}
Let $m_k:=|\mathcal S\cap\{0,1,\dots,k-1\}|=k-n_k$.
We consider two cases.

\emph{Case 1: $n_k\ge k/2$.}
Then \eqref{eq:Phi_by_nk} gives
\[
\Phi_k\le \frac{4}{\tau^2(n_k+2)^2}\le \frac{16}{\tau^2(k+2)^2}.
\]

\emph{Case 2: $n_k< k/2$ (hence $m_k>k/2$).}
For every sharp index $\ell\in\mathcal S$, we have $\|g_{\ell+1}\|\le \frac14\|g_\ell\|$ by definition.
For every other index, we have the growth bound \eqref{eq:grad_growth}.
Therefore, after $k$ iterations,
\[
\|g_k\|
\le \kappa_{\mathrm{inc}}^{\,n_k}\left(\frac14\right)^{m_k}\|g_0\|
\le \kappa_{\mathrm{inc}}^{\,k/2}\left(\frac14\right)^{k/2}\|g_0\|
=\left(\frac{\kappa_{\mathrm{inc}}}{4}\right)^{k/2}\|g_0\|.
\]
Since $\kappa_{\mathrm{inc}}/4<1$, the bounded level set and convexity imply
\[
\Phi_k=f(\bar x_k)-f_\star
\le \ip{g_k}{\bar x_k-x_\star}
\le \|\bar x_k-x_\star\|\,\|g_k\|
\le D\,\|g_k\|.
\]
Therefore,
\[
\Phi_k\le D\,\|g_0\|\left(\frac{\kappa_{\mathrm{inc}}}{4}\right)^{k/2}.
\]
Because $\sup_{k\ge 0}(k+2)^2\left(\frac{\kappa_{\mathrm{inc}}}{4}\right)^{k/2}<\infty$, there exists a finite constant $C_{\exp}>0$ such that
\[
\Phi_k\le \frac{C_{\exp}}{(k+2)^2}\qquad\forall k\ge 0.
\]

Combining both cases, there exists $C_\Phi>0$ such that for all $k\ge 0$,
\[
\Phi_k\le \frac{C_\Phi}{(k+2)^2},
\]
which proves $\Phi_k=\cO(1/k^2)$.

\vspace{1pt}
\noindent\textbf{Step 4: gradient rate and $\varepsilon$-complexity.}
By Lemma~\ref{lem:gap-grad} with $h=f$ and $L=L_1$,
\[
\|g_k\|^2\le 2L_1\Phi_k\le \frac{2L_1 C_\Phi}{(k+2)^2},
\]
hence $\|g_k\|\le \sqrt{2L_1 C_\Phi}/(k+2)=\cO(1/k)$.
Therefore, to ensure $\|g_k\|\le \varepsilon$, it suffices to take
\[
k\ \ge\ \frac{\sqrt{2L_1 C_\Phi}}{\varepsilon}-2
=\cO(\varepsilon^{-1}).
\]
This proves items 2--3.
\end{proof}

\begin{proof}[Proof of Lemma~\ref{lem:delta_o_g}]
From Lemma~\ref{lem:delta_control}, $\bar\delta_k\le \Delta_k^H+L_2D_X(k)$.

\vspace{1pt}
\noindent\textbf{Bound $\Delta_k^H$.}
By Jensen's inequality and $\|\cdot\|_2\le \|\cdot\|_F$,
\[
\Delta_k^H\le D_H(\tilde{\mathcal H}_k).
\]
We apply the Hessian-tracker dispersion recursion (Lemma~\ref{lem:DH_rec}),
noting that $\mathcal H_k$ is obtained from $\tilde{\mathcal H}_{k-1}$
via the post-mixing update at iteration $k{-}1$:
\begin{align*}
D_H(\tilde{\mathcal H}_k)
&\le \rho^{\tau_k}D_H(\mathcal H_k)\\
&\le \rho^{\tau_k}\rho^{t_{k-1}}\bigl(D_H(\tilde{\mathcal H}_{k-1})
\!+\!2\sqrt{d}L_2D\bigr)\\
&\le \rho^{\tau_k}(B_H\!+\!2\sqrt{d}L_2D),
\end{align*}
where $B_H:=\sup_{j\ge 0}D_H(\tilde{\mathcal H}_j)<\infty$ is the uniform bound
established in Lemma~\ref{lem:tilde_hess_bounded}.
In particular, since $B_H$ is finite, we obtain
\[
\Delta_k^H\le \rho^{\tau_k}\big(B_H+2\sqrt{d}\,L_2D\big)
=:\hat C_H\,\rho^{\tau_k}.
\]

\vspace{1pt}
\noindent\textbf{Bound $D_X(k)$.}
By Lemma~\ref{lem:DX_post},
\[
D_X(k)=D(X_k)\le \rho^{t_{k-1}}\Big(D(\tilde X_{k-1})+\frac{\|S_{k-1}\|}{\sqrt N}\Big).
\]
Moreover, $D(\tilde X_{k-1})\le 2D$ and Lemma~\ref{lem:step_bound_app} gives $\|S_{k-1}\|/\sqrt N\le \sqrt{G_g/M}$.
Thus
\[
D_X(k)\le \rho^{t_{k-1}}\Big(2D+\sqrt{\frac{G_g}{M}}\Big).
\]

\vspace{1pt}
Therefore,
\[
\bar\delta_k\le \hat C_H\,\rho^{\tau_k}
+L_2\rho^{t_{k-1}}\Big(2D+\sqrt{\frac{G_g}{M}}\Big)
\le C_\delta\max\{\rho^{\tau_k},\rho^{t_{k-1}}\}.
\]
Using \eqref{eq:delta_schedule}, we obtain $\bar\delta_k\le \|g_k\|^{1+\gamma}$ for all sufficiently large $k$.
\end{proof}

% (Full proof inserted; see v4 script source for details.)
\begin{proof}[Proof of Theorem~\ref{thm:superlinear}]
Let $x_k:=\bar x_k$, $g_k:=\nabla f(x_k)$, and $H_k:=\nabla^2 f(x_k)$.
Under Assumption~\ref{ass:strong}, for all $k$ we have $H_k\succeq \mu I$ on the level set.

By definition, $r_k:=(H_k+\lambda_k I)\bar s_k+g_k$ and Proposition~\ref{prop:inexact-rn} gives for all large $k$
\[
(H_k+\lambda_k I)\bar s_k=-g_k+r_k,\qquad \|r_k\|\le \eta\lambda_k\|\bar s_k\|,\qquad L_2\|\bar s_k\|\le \lambda_k,
\]
with $\eta\in(0,1/12]$. Moreover, Item~3 of Proposition~\ref{prop:inexact-rn} implies $D_X(k)=\mathcal O(\lambda_k)$.
Under the additional mixing schedule \eqref{eq:delta_schedule} and the bound
$D_X(k)\le \rho^{t_{k-1}}\big(2D+\sqrt{G_g/M}\big)$ (see Lemma~\ref{lem:DX_post} and Lemma~\ref{lem:step_bound_app}),
we further have $D_X(k)=o(\lambda_k)$ as $k\to\infty$.
Lemma~\ref{lem:bridge-grad} yields
$\|\tilde{\bar g}_k-g_k\|\le L_1D_X(k)$, hence $\|\tilde{\bar g}_k\|=\|g_k\|+o(\lambda_k)$.
Combining with Lemma~\ref{lem:delta_o_g} and $\lambda_k=\tilde{\bar\lambda}_k$ gives $\lambda_k=\cO(\sqrt{\|g_k\|})$.

Since $H_k+\lambda_k I\succeq \mu I$,
\[
\mu\|\bar s_k\|\le \|(H_k+\lambda_k I)\bar s_k\|=\|g_k-r_k\|\le \|g_k\|+\eta\lambda_k\|\bar s_k\|.
\]
For large $k$, $\eta\lambda_k\le \mu/2$, so $\|\bar s_k\|\le \frac{2}{\mu}\|g_k\|$.

Using $g_{k+1}=g_k+H_k\bar s_k+e_k$ with $\|e_k\|\le \frac{L_2}{2}\|\bar s_k\|^2$ and $g_k+H_k\bar s_k=-\lambda_k\bar s_k+r_k$,
\[
\|g_{k+1}\|\le (1+\eta)\lambda_k\|\bar s_k\|+\frac{L_2}{2}\|\bar s_k\|^2
\le \cO(\|g_k\|^{3/2})+\cO(\|g_k\|^2)=\cO(\|g_k\|^{3/2}),
\]
which proves the claim.
\end{proof}



\begin{proof}[Proof of Theorem~\ref{thm:end2end}]
Fix $\varepsilon>0$ and let $K(\varepsilon)$ be the first index with $\|g_{K(\varepsilon)}\|\le \varepsilon$.
Then $\|g_k\|\ge \varepsilon$ for all $0\le k\le K(\varepsilon)-1$.
Let $j:=k-K_0(\varepsilon)\ge 0$ and consider the shifted sequence $\{\bar x_{K_0(\varepsilon)+j}\}_{j\ge 0}$.
On the index set $\{K_0(\varepsilon),\dots,K(\varepsilon)-1\}$, Proposition~\ref{prop:inexact-rn} holds with $\eta=1/12$,
so the proof of Theorem~\ref{thm:main} applies to the shifted sequence and yields that there exists a constant
$C_{\nabla}>0$ such that
\[
\|g_{K_0(\varepsilon)+j}\|\le \frac{C_{\nabla}}{j+2},\qquad \forall j\ge 0.
\]
Thus $\|g_k\|\le \varepsilon$ once $j\ge C_{\nabla}/\varepsilon-2$, i.e., $k\ge K_0(\varepsilon)+C_{\nabla}/\varepsilon-2$,
which gives $K(\varepsilon)\le K_0(\varepsilon)+C_K\varepsilon^{-1}$ for some constant $C_K>0$.
\end{proof}

\begin{proof}[Proof of Theorem~\ref{thm:comm}]
Under \eqref{eq:log_schedule}, $\tau_k=t_k=\lceil p\log(k+2)/(-\log\rho)\rceil\le c_0+c_1\log(k+2)$.
Hence $\sum_{k=0}^{K(\varepsilon)-1}(\tau_k+t_k)=\cO(K(\varepsilon)\log(K(\varepsilon)+2))$.
Combine with Theorem~\ref{thm:end2end}.
\end{proof}

\end{document}